% Sensibilisation sur les effets de certaines substances sur la transmission synaptique — SVT Tle D — Cours signé OnBuch+
% Programme officiel MINESEC. Compiler avec : tectonic N09-transmission-synaptique-substances.tex
\newcommand{\DOCMATIERE}{SVT}
\newcommand{\DOCNIVEAU}{Tle D}
\newcommand{\DOCMODULE}{Module I : Le Monde Vivant — Les synapses}
\newcommand{\DOCLECON}{N09}
\newcommand{\DOCTITRE}{Sensibilisation sur les effets de certaines substances sur la transmission synaptique}
% OnBuch+ — préambule « cours signés » ENRICHI (compilé avec tectonic / XeLaTeX)
% Système visuel « Pop » d'OnBuch+ : fond crème (jamais blanc), bordures encre
% épaisses, ombres « dures » décalées, titres Archivo Black en capitales.
% Version MATHÉMATIQUES : graphes (pgfplots/tikz), boîtes pédagogiques
% (définition, propriété, méthode, attention, le savais-tu, exercice résolu,
% exercice, TP, bilan), page de garde et signature OnBuch+ ancrée en bas.
%
% Macros attendues dans main.tex AVANT \input{preamble} :
%   \DOCMATIERE  \DOCNIVEAU  \DOCMODULE  \DOCLECON (numéro)  \DOCTITRE
\documentclass[11pt]{article}

\usepackage{fontspec}
\usepackage[french]{babel}
\usepackage[a4paper,margin=2cm,top=3.4cm,bottom=2.8cm,headheight=1.4cm,footskip=1.3cm]{geometry}
\usepackage{amsmath,amssymb,amsfonts}
\usepackage{mathtools} % \xmapsto, \coloneqq... souvent utilisés par les modèles
\usepackage{cancel} % simplifications barrées (bilans de forces)
% \abs{z} (module, valeur absolue) et \norm{u} (norme), utilisés par les modèles
\providecommand{\abs}{}\let\abs\relax\DeclarePairedDelimiter{\abs}{\lvert}{\rvert}
\providecommand{\norm}{}\let\norm\relax\DeclarePairedDelimiter{\norm}{\lVert}{\rVert}
\usepackage[table]{xcolor} % [table] : fournit aussi \rowcolors (alternance de lignes)
\usepackage{pagecolor}
\usepackage{tikz}
\usetikzlibrary{arrows.meta,positioning,calc,shapes.geometric,shapes.misc,shapes.arrows,decorations.pathmorphing,decorations.pathreplacing,decorations.markings,patterns,shadows,mindmap,trees,fit,backgrounds,matrix,intersections,angles,quotes}
\usepackage{pgfplots}
\usepackage[version=4]{mhchem} % équations nucléaires \ce{^{235}_{92}U}
\usepackage[european,siunitx]{circuitikz} % schémas électriques
\pgfplotsset{compat=1.18}
\usepgfplotslibrary{fillbetween,groupplots} % aires entre courbes (calcul intégral)
\usepackage{enumitem}
\usepackage{fancyhdr}
% [most] charge toutes les bibliothèques tcolorbox (listings, minted, raster,
% documentation...) et gonfle inutilement la mémoire TeX sur un document long
% (~300 boîtes) : on ne charge que ce qui est réellement utilisé.
\usepackage[skins,breakable]{tcolorbox}
\usepackage{graphicx}
\usepackage{colortbl}
\usepackage{booktabs}
\usepackage{tabularx}
\usepackage{diagbox} % case d'en-tête diagonale des tableaux à double entrée
% Colonnes extensibles centrée / gauche / droite (C, L, R), souvent utilisées par les modèles
\newcolumntype{C}{>{\centering\arraybackslash}X}
\newcolumntype{L}{>{\raggedright\arraybackslash}X}
\newcolumntype{R}{>{\raggedleft\arraybackslash}X}
\usepackage{array}
\usepackage{multicol}
\usepackage{siunitx}
% parse-numbers=false : accepte aussi \SI{2,5 \times 10^{-3}}{...}
\sisetup{locale=FR,output-decimal-marker={,},group-minimum-digits=4,parse-numbers=false}
\DeclareSIUnit\ohm{\ensuremath{\symup{\Omega}}} % U+2126 absent de la police
\DeclareSIUnit\division{div} % réglages d'oscilloscope (ms/div, V/div)
\DeclareSIUnit\tr{tr} % tours (vitesse de rotation, ex. \SI{1500}{\tr\per\minute})
\DeclareSIUnit\molar{M} % concentration molaire (ex. \SI{3}{\molar}, \SI{1}{\micro\molar})
\DeclareSIUnit\an{an} % année (le modèle confond parfois avec \year, un compteur TeX) — ex. \SI{5}{\kilo\meter\per\an}
\DeclareSIUnit\FCFA{FCFA} % franc CFA (ex. \SI{500}{\FCFA\per\kilo\gram})
\DeclareSIUnit\degreeBrix{\textdegree Brix} % teneur en sucre d'un sirop/jus (réfractométrie)
\DeclareSIUnit\month{mois} % le modèle utilise parfois \month, un compteur TeX, comme unité de durée
\usepackage{qrcode}
% \degree : commande fréquemment utilisée par les modèles pour un angle ou une
% température en dehors de \SI{}{} (non fournie par les paquets chargés).
\providecommand{\degree}{^{\circ}}
% \qed (fin de démonstration), utilisé par les modèles sans amsthm
\providecommand{\qed}{\hfill$\square$}
% \barre : double barre des valeurs interdites dans un tableau de variations (tkz-tab)
\providecommand{\barre}{\ensuremath{\|}}
% environnement proof (amsthm non chargé) utilisé par les modèles
\ifdefined\proof\else\newenvironment{proof}[1][Démonstration]{\par\noindent\textit{#1.}\ }{\qed\par}\fi
\usepackage{lastpage}
\usepackage{hyperref}
\hypersetup{hidelinks}

% ============================================================
% Palette « Pop » OnBuch+ — mêmes hex que la classe P (plus_ui.dart)
% ============================================================
\definecolor{poporange}{HTML}{F59321}
\definecolor{popdark}{HTML}{E07A0C}
\definecolor{popcream}{HTML}{FFF6E8}
\definecolor{popink}{HTML}{1C1714}
\definecolor{poppurple}{HTML}{7A5AE0}
\definecolor{popgreen}{HTML}{2E9E6A}
\definecolor{poppink}{HTML}{E0457B}
\definecolor{popblue}{HTML}{2F7DE1}
\definecolor{popgold}{HTML}{C98A12}
\definecolor{popmuted}{HTML}{8A7F76}
\definecolor{popline}{HTML}{EADFD2}
\definecolor{popgray}{HTML}{8A7F76}
\colorlet{popgrey}{popgray}
% Alias tolérants (le modèle nomme parfois les couleurs différemment)
\colorlet{popwhite}{white}
\colorlet{popblack}{popink}
\colorlet{popred}{poppink}
\colorlet{popyellow}{popgold}
% Teintes claires (fonds de tableaux/encadrés) — le modèle nomme parfois
% les couleurs avec un suffixe L (ex. popblueL) sans les avoir définies.
\colorlet{poporangeL}{poporange!20}
\colorlet{popdarkL}{popdark!20}
\colorlet{poppurpleL}{poppurple!20}
\colorlet{popgreenL}{popgreen!20}
\colorlet{poppinkL}{poppink!20}
\colorlet{popblueL}{popblue!20}
\colorlet{popgoldL}{popgold!20}
\colorlet{popredL}{popred!20}
\colorlet{popgrayL}{popgray!20}
% Teintes claires pour fonds de tableaux / zones
\colorlet{popblueL}{popblue!10!white}
\colorlet{poporangeL}{poporange!14!white}
\colorlet{popgreenL}{popgreen!12!white}
\colorlet{poppurpleL}{poppurple!12!white}
\colorlet{poppinkL}{poppink!12!white}
\colorlet{popgoldL}{popgold!14!white}

\pagecolor{popcream}

% ============================================================
% Typographie
% ============================================================
\defaultfontfeatures{Ligatures=TeX}
\setmainfont{PlusJakartaSans-Medium.ttf}[Path=../../../fonts/,BoldFont=PlusJakartaSans-Bold.ttf]
% Maths en sans-serif (Fira Math) avec chiffres et lettres droites de Plus
% Jakarta, pour que les formules se fondent dans le texte.
\usepackage[math-style=ISO,bold-style=ISO]{unicode-math}
\setmathfont{FiraMath-Regular.otf}[Path=../../../fonts/]
\setmathfont{PlusJakartaSans-Medium.ttf}[Path=../../../fonts/,range={up/num,up/{Latin,latin},\mathup}]
\usepackage{icomma} % virgule décimale sans espace en mode math
% Caractères mathématiques Unicode absents des polices de texte (Plus Jakarta,
% Archivo Black) : rendus par leur équivalent mathématique, en texte comme en
% formule — sinon ils s'affichent en carré vide (ex. « Arithmétique dans ℤ »).
\usepackage{newunicodechar}
\newunicodechar{ℝ}{\ensuremath{\mathbb{R}}}
\newunicodechar{ℤ}{\ensuremath{\mathbb{Z}}}
\newunicodechar{ℂ}{\ensuremath{\mathbb{C}}}
\newunicodechar{ℕ}{\ensuremath{\mathbb{N}}}
\newunicodechar{ℚ}{\ensuremath{\mathbb{Q}}}
\newunicodechar{𝔻}{\ensuremath{\mathbb{D}}}
\newunicodechar{α}{\ensuremath{\alpha}}
\newunicodechar{β}{\ensuremath{\beta}}
\newunicodechar{γ}{\ensuremath{\gamma}}
\newunicodechar{δ}{\ensuremath{\delta}}
\newunicodechar{ε}{\ensuremath{\varepsilon}}
\newunicodechar{θ}{\ensuremath{\theta}}
\newunicodechar{λ}{\ensuremath{\lambda}}
\newunicodechar{μ}{\ensuremath{\mu}}
\newunicodechar{σ}{\ensuremath{\sigma}}
\newunicodechar{τ}{\ensuremath{\tau}}
\newunicodechar{φ}{\ensuremath{\varphi}}
\newunicodechar{ψ}{\ensuremath{\psi}}
\newunicodechar{ω}{\ensuremath{\omega}}
\newunicodechar{Σ}{\ensuremath{\Sigma}}
% majuscules produites par \MakeUppercase dans les titres (α-aminés -> Α)
\newunicodechar{Α}{\ensuremath{\alpha}}
\newunicodechar{Β}{\ensuremath{\beta}}
\newunicodechar{Γ}{\ensuremath{\Gamma}}
\newunicodechar{Δ}{\ensuremath{\Delta}}
\newunicodechar{Ε}{\ensuremath{\varepsilon}}
\newunicodechar{Θ}{\ensuremath{\Theta}}
\newunicodechar{Λ}{\ensuremath{\Lambda}}
\newunicodechar{Μ}{\ensuremath{\mu}}
\newunicodechar{Τ}{\ensuremath{\tau}}
\newunicodechar{Φ}{\ensuremath{\Phi}}
\newunicodechar{Ψ}{\ensuremath{\Psi}}
\newunicodechar{Ω}{\ensuremath{\Omega}}
\newunicodechar{↦}{\ensuremath{\mapsto}}
\newunicodechar{∈}{\ensuremath{\in}}
\newunicodechar{∉}{\ensuremath{\notin}}
\newunicodechar{∧}{\ensuremath{\wedge}}
\newunicodechar{⇔}{\ensuremath{\Leftrightarrow}}
\newunicodechar{⟺}{\ensuremath{\Longleftrightarrow}}
\newunicodechar{⇒}{\ensuremath{\Rightarrow}}
\newunicodechar{⟹}{\ensuremath{\Longrightarrow}}
\newunicodechar{∘}{\ensuremath{\circ}}
\newunicodechar{∪}{\ensuremath{\cup}}
\newunicodechar{∩}{\ensuremath{\cap}}
\newunicodechar{≡}{\ensuremath{\equiv}}
\newunicodechar{⊂}{\ensuremath{\subset}}
\newunicodechar{⊥}{\ensuremath{\perp}}
\newunicodechar{∃}{\ensuremath{\exists}}
\newunicodechar{∀}{\ensuremath{\forall}}
\newunicodechar{∛}{\ensuremath{\sqrt[3]{\,}}}
\newunicodechar{∜}{\ensuremath{\sqrt[4]{\,}}}
\newunicodechar{⃗}{\ensuremath{{}^{\to}}}
\newunicodechar{ⁿ}{\ifmmode^{n}\else\textsuperscript{\ensuremath{n}}\fi}
\newunicodechar{ˣ}{\ifmmode^{x}\else\textsuperscript{\ensuremath{x}}\fi}
\newunicodechar{ᵇ}{\ifmmode^{b}\else\textsuperscript{\ensuremath{b}}\fi}
\newunicodechar{ʳ}{\ifmmode^{r}\else\textsuperscript{\ensuremath{r}}\fi}
\newunicodechar{ᵘ}{\ifmmode^{u}\else\textsuperscript{\ensuremath{u}}\fi}
\newunicodechar{ʸ}{\ifmmode^{y}\else\textsuperscript{\ensuremath{y}}\fi}
\newunicodechar{ᵃ}{\ifmmode^{a}\else\textsuperscript{\ensuremath{a}}\fi}
\newunicodechar{ᵉ}{\ifmmode^{e}\else\textsuperscript{\ensuremath{e}}\fi}
\newunicodechar{ᵏ}{\ifmmode^{k}\else\textsuperscript{\ensuremath{k}}\fi}
\newunicodechar{ᵖ}{\ifmmode^{p}\else\textsuperscript{\ensuremath{p}}\fi}
\newunicodechar{ᴺ}{\ifmmode^{N}\else\textsuperscript{\ensuremath{N}}\fi}
\newunicodechar{ᶜ}{\ifmmode^{c}\else\textsuperscript{\ensuremath{c}}\fi}
\newunicodechar{ʰ}{\ifmmode^{h}\else\textsuperscript{\ensuremath{h}}\fi}
\newunicodechar{ᵠ}{\ifmmode^{q}\else\textsuperscript{\ensuremath{q}}\fi}
\newunicodechar{ₙ}{\ifmmode_{n}\else\textsubscript{\ensuremath{n}}\fi}
\newunicodechar{ₐ}{\ifmmode_{a}\else\textsubscript{\ensuremath{a}}\fi}
\newunicodechar{ᵢ}{\ifmmode_{i}\else\textsubscript{\ensuremath{i}}\fi}
\newunicodechar{ᵣ}{\ifmmode_{r}\else\textsubscript{\ensuremath{r}}\fi}
\newunicodechar{ₖ}{\ifmmode_{k}\else\textsubscript{\ensuremath{k}}\fi}
\newunicodechar{ᵦ}{\ifmmode_{\beta}\else\textsubscript{\ensuremath{\beta}}\fi}
\newunicodechar{ₓ}{\ifmmode_{x}\else\textsubscript{\ensuremath{x}}\fi}
\newfontfamily\popdisplay{ArchivoBlack-Regular.ttf}[Path=../../../fonts/]
\newfontfamily\poplogo{SpaceGrotesk-Bold.ttf}[Path=../../../fonts/]
\newfontfamily\popsemi{PlusJakartaSans-SemiBold.ttf}[Path=../../../fonts/]
\newfontfamily\popxbold{PlusJakartaSans-ExtraBold.ttf}[Path=../../../fonts/]
\newfontfamily\popxboldls{PlusJakartaSans-ExtraBold.ttf}[Path=../../../fonts/,LetterSpace=3.0]

\setlist[enumerate]{leftmargin=1.7em,itemsep=5pt,topsep=4pt}
\setlist[itemize]{leftmargin=1.7em,itemsep=4pt,topsep=4pt,label={\color{poporange}\textbullet}}
\setlength{\parindent}{0pt}
\setlength{\parskip}{5pt}
\renewcommand{\arraystretch}{1.25}

% pgfplots : style Pop par défaut
\pgfplotsset{
  every axis/.append style={axis line style={popink,line width=1pt},tick style={popink},
    grid style={popline,line width=0.5pt},label style={font=\small},tick label style={font=\footnotesize}},
  popcurve/.style={poporange,line width=1.8pt,no markers,smooth},
}
% Même style disponible dans un \draw TikZ simple (hors pgfplots)
\tikzset{popcurve/.style={poporange,line width=1.8pt}}
\tikzset{popgrid/.style={popline,very thin}}
\colorlet{popcurve}{poporange} % le nom du style est parfois utilisé comme couleur

% ============================================================
% Pill / squiggle
% ============================================================
\newcommand{\poppill}[3]{%
  \tikz[baseline=-0.55ex]\node[draw=popink,line width=1.2pt,rounded corners=2.6pt,%
    fill=#2,inner xsep=6.5pt,inner ysep=3pt]{%
    \popxboldls\fontsize{7}{7}\selectfont\color{#3}\MakeUppercase{#1}};%
}
\newcommand{\popsquiggle}[1]{%
  \tikz[baseline=-0.4ex]{
    \draw[#1,line width=2.6pt,line cap=round]
      plot [smooth] coordinates {(0,0.09) (0.34,0.01) (0.68,0.21) (1.02,0.01) (1.36,0.10)};
  }%
}
% Mot-clé surligné (marqueur orange sous le texte)
\newcommand{\cle}[1]{{\popxbold\color{popdark}#1}}

% ============================================================
% En-tête / pied
% ============================================================
\pagestyle{fancy}
\fancyhf{}
\renewcommand{\headrulewidth}{0pt}
\renewcommand{\footrulewidth}{0pt}
\lhead{\raisebox{-0.15ex}{\poplogo\fontsize{13}{13}\selectfont\color{popink}OnBuch\color{poporange}+}}
\rhead{\poppill{\DOCMATIERE\ \textbullet\ \DOCNIVEAU\ \textbullet\ Leçon \DOCLECON}{white}{popink}}
\lfoot{\popxbold\fontsize{7.6}{7.6}\selectfont\color{popmuted}\MakeUppercase{Cours signé OnBuch+}\ \textbullet\ {\popsemi\DOCTITRE}}
\rfoot{\tikz[baseline=-0.6ex]\node[rounded corners=8.5pt,fill=popink,minimum height=17pt,minimum width=17pt,inner xsep=5pt,inner sep=0]{\popxbold\fontsize{8}{8}\selectfont\color{popcream}\thepage\,/\,\pageref*{LastPage}};}

% ============================================================
% Titres
% ============================================================
\newcommand{\chaptitre}[2]{%
  \begin{tcolorbox}[enhanced,colback=poporange,colframe=popink,boxrule=2.6pt,arc=11pt,
      drop shadow=popink,left=15pt,right=15pt,top=12pt,bottom=11pt,before skip=3pt,after skip=18pt]
    \poppill{#2}{popink}{popcream}\par\vspace{9pt}
    {\raggedright\hyphenpenalty=10000\exhyphenpenalty=10000\popdisplay\fontsize{22}{24}\selectfont\color{popcream}\MakeUppercase{#1}\par}
  \end{tcolorbox}
}
% Section numérotée : pastille orange avec le numéro + titre Archivo
\newcounter{popsec}
\newcommand{\coursec}[1]{%
  \stepcounter{popsec}\setcounter{popsub}{0}%
  \par\vspace{16pt}\noindent
  \begin{minipage}{\linewidth}
  \tikz[baseline=-0.7ex]\node[draw=popink,line width=1.6pt,fill=poporange,rounded corners=4pt,minimum size=22pt,inner sep=0]{\popdisplay\fontsize{12}{12}\selectfont\color{popcream}\thepopsec};%
  \hspace{8pt}{\popdisplay\fontsize{15}{17}\selectfont\color{popink}\MakeUppercase{#1}}\par
  \vspace{4pt}\noindent{\color{poporange}\rule{36pt}{3.6pt}}
  \end{minipage}\par\nopagebreak\vspace{8pt}
}
\newcounter{popsub}
\newcommand{\courssub}[1]{%
  \stepcounter{popsub}%
  \par\vspace{9pt}\noindent{\popxbold\fontsize{11.5}{13}\selectfont\color{popink}{\color{poporange}\thepopsec.\thepopsub}\hspace{6pt}#1}\par\nopagebreak\vspace{4pt}
}

% ============================================================
% Boîtes pédagogiques — même recette : fond blanc, bordure épaisse colorée,
% ombre dure encre, badge plein. Titre optionnel [#1].
% ============================================================
\tcbset{popbase/.style={breakable,enhanced,colback=white,boxrule=2.3pt,arc=9pt,
  shadow={3pt}{-3pt}{0pt}{popink},left=14pt,right=14pt,top=11pt,bottom=11pt,before skip=11pt,after skip=15pt,
  fontupper=\popsemi\fontsize{10.6}{14.5}\selectfont\color{popink},
  fontlower=\popsemi\fontsize{10.6}{14.5}\selectfont\color{popink}}}
% Coche dessinée (le glyphe ✓ n'existe pas dans les polices du document)
\newcommand{\popcheck}[1]{\tikz[baseline=-0.5ex]\draw[#1,line width=1.6pt,line cap=round,line join=round](-0.13,0.0)--(-0.04,-0.09)--(0.13,0.1);}
\providecommand{\coche}{\popcheck{popgreen}}
\providecommand{\resultat}[1]{\par\smallskip\noindent\fcolorbox{popgreen}{popgreenL}{\parbox{\dimexpr\linewidth-2\fboxsep-2\fboxrule\relax}{\textbf{Résultat.} #1}}\par\smallskip}
\newcommand{\popboxhead}[3]{\poppill{#1}{#2}{white}\ifx\relax#3\relax\else\hspace{6pt}{\popxbold\fontsize{10.6}{12}\selectfont\color{popink}#3}\fi\par\vspace{7pt}}

\newtcolorbox{aretenir}[1][]{popbase,colframe=popgreen,before upper={\popboxhead{À retenir}{popgreen}{#1}}}
\newtcolorbox{exemplebox}[1][]{popbase,colframe=poporange,before upper={\popboxhead{Exemple}{poporange}{#1}}}
\newtcolorbox{definition}[1][]{popbase,colframe=popblue,before upper={\popboxhead{Définition}{popblue}{#1}}}
\newtcolorbox{propriete}[1][]{popbase,colframe=poppurple,before upper={\popboxhead{Propriété}{poppurple}{#1}}}
\newtcolorbox{methode}[1][]{popbase,colframe=popgold,before upper={\popboxhead{Méthode}{popgold}{#1}}}
\newtcolorbox{attention}[1][]{popbase,colframe=poppink,before upper={\popboxhead{Attention}{poppink}{#1}}}
\newtcolorbox{experience}[1][]{popbase,colframe=popblue,colback=popblueL,before upper={\popboxhead{Expérience}{popblue}{#1}}}
% « Le savais-tu ? » : carte violette pleine (contexte, culture, Cameroun)
\newtcolorbox{savaistu}[1][]{popbase,colback=poppurple,colframe=popink,
  fontupper=\popsemi\fontsize{10.4}{14.2}\selectfont\color{white},
  before upper={\poppill{Le savais-tu\,?}{white}{poppurple}\ifx\relax#1\relax\else\hspace{6pt}{\popxbold\color{white}#1}\fi\par\vspace{7pt}}}
% Exercice résolu : énoncé en haut, \tcblower puis la solution sur fond crème
\newcounter{popexo}\newcounter{popexr}
\newtcolorbox{exoresolu}[1][]{popbase,colframe=poporange,colbacklower=popcream,
  segmentation style={popink,line width=1.2pt,dashed},
  before upper={\stepcounter{popexr}\popboxhead{Exercice résolu \thepopexr}{poporange}{#1}},
  before lower={\poppill{Solution}{popink}{popcream}\par\vspace{6pt}}}
% Exercice (non résolu en place) — la correction est dans la partie Corrigés
\newtcolorbox{exercice}[1][]{popbase,colframe=popink,
  before upper={\stepcounter{popexo}\popboxhead{Exercice \thepopexo}{popink}{#1}}}
% Corrigé
\newtcolorbox{corrige}[1][]{popbase,colframe=popgreen,colback=popgreenL,
  before upper={\popboxhead{Corrigé}{popgreen}{#1}}}
% Objectifs / prérequis / bilan
\newtcolorbox{objectifs}{popbase,colframe=popink,colback=poporangeL,
  before upper={\popboxhead{Objectifs de la leçon}{popink}{}}}
\newtcolorbox{prerequis}{popbase,colframe=popink,colback=popblueL,
  before upper={\popboxhead{Prérequis}{popblue}{}}}
\newtcolorbox{bilan}{popbase,colframe=popink,colback=white,boxrule=2.8pt,
  before upper={\popboxhead{Fiche bilan}{poporange}{L'essentiel à connaître pour l'examen}}}
% Figure encadrée avec légende
\newtcolorbox{popfigure}[1][]{enhanced,breakable=false,colback=white,colframe=popline,boxrule=1.4pt,arc=8pt,
  left=10pt,right=10pt,top=10pt,bottom=8pt,before skip=10pt,after skip=12pt,halign=center,
  lower separated=false,halign lower=center,
  fontlower=\popsemi\fontsize{9}{11.5}\selectfont\color{popmuted},
  #1}
\newcommand{\legende}[1]{\par\vspace{6pt}{\popsemi\fontsize{9}{11.5}\selectfont\color{popmuted}#1}}

% ============================================================
% Signature OnBuch+
% ============================================================
\newcommand{\poplogoinline}[1]{{\poplogo\fontsize{#1}{#1}\selectfont\color{popink}OnBuch\color{poporange}+}}
\newcommand{\signatureonbuch}{%
  \begin{tcolorbox}[enhanced,colback=popink,colframe=popink,boxrule=2.2pt,arc=10pt,
      drop shadow=poporange,left=14pt,right=14pt,top=10pt,bottom=10pt,before skip=0pt,after skip=0pt]
    \tikz[baseline=-0.7ex]\node[circle,fill=popgreen,draw=popcream,line width=1.3pt,minimum size=22pt,inner sep=0]{\popcheck{white}};%
    \hspace{9pt}\begin{minipage}[c]{0.62\linewidth}
      {\popxbold\fontsize{12}{13}\selectfont\color{popcream}Signé\ \textbullet\ {\poplogo OnBuch\color{poporange}+}}\par\vspace{2pt}
      {\raggedright\popsemi\fontsize{8}{10.5}\selectfont\color{popline}\MakeUppercase{Équipe pédagogique OnBuch+}\\\MakeUppercase{Conforme au programme officiel MINESEC}\par}
    \end{minipage}\hfill
    {\poplogo\fontsize{22}{22}\selectfont\color{popcream}OnBuch\color{poporange}+}
  \end{tcolorbox}%
}

% ============================================================
% Page de garde
% #1 titre · #2 matière · #3 niveau · #4 module · #5 n° leçon · #6 durée ·
% #7 description / objectifs (texte ou itemize)
% ============================================================
\newcommand{\pagegarde}[7]{%
  \thispagestyle{empty}
  \newgeometry{a4paper,margin=1.7cm,top=1.6cm,bottom=1.4cm}%
  \begin{tikzpicture}[remember picture,overlay]
    \fill[poporange!18] ([xshift=-3.2cm,yshift=-3.6cm]current page.north east) circle (4.6cm);
    \fill[poppurple!14] ([xshift=2.4cm,yshift=7.6cm]current page.south west) circle (3.8cm);
  \end{tikzpicture}%
  {\poplogo\fontsize{30}{30}\selectfont\color{popink}OnBuch\color{poporange}+}\hfill
  \raisebox{0.6ex}{\poppill{Cours signé \textbullet\ Premium}{popink}{popcream}}\par
  \vspace{1pt}\popsquiggle{poppurple}\par
  \vspace{8pt}
  \begin{tcolorbox}[enhanced,colback=poporange,colframe=popink,boxrule=2.8pt,arc=13pt,
      drop shadow=popink,left=18pt,right=18pt,top=12pt,bottom=13pt,after skip=10pt]
    \poppill{#2 \textbullet\ #3}{popink}{popcream}\hspace{5pt}\poppill{#4}{white}{popink}\par\vspace{8pt}
    {\popdisplay\fontsize{14}{14}\selectfont\color{popink}LEÇON #5}\par\vspace{4pt}
    {\raggedright\hyphenpenalty=10000\exhyphenpenalty=10000\popdisplay\fontsize{25}{27}\selectfont\color{popcream}\MakeUppercase{#1}\par}
    \vspace{8pt}
    {\popxbold\fontsize{9.5}{9.5}\selectfont\color{popink}\MakeUppercase{Durée conseillée : #6}}
  \end{tcolorbox}
  \begin{tcolorbox}[enhanced,colback=white,colframe=popink,boxrule=2.2pt,arc=10pt,
      drop shadow=popink,left=16pt,right=16pt,top=10pt,bottom=10pt,after skip=8pt]
    \poppill{Dans cette leçon}{poppurple}{white}\par\vspace{5pt}
    \popsemi\fontsize{9.3}{12.6}\selectfont\color{popink}#7
  \end{tcolorbox}
  \noindent\begin{minipage}[t]{0.487\linewidth}
    \begin{tcolorbox}[enhanced,colback=popgreen,colframe=popink,boxrule=2.2pt,arc=10pt,
        drop shadow=popink,left=11pt,right=11pt,top=10pt,bottom=10pt,height=3.1cm,valign=center]
      \tcbox[colback=white,colframe=popink,boxrule=1.3pt,arc=5pt,boxsep=0pt,nobeforeafter,
        left=3pt,right=3pt,top=3pt,bottom=3pt,valign=center]{\qrcode[height=1.9cm]{https://play.google.com/store/apps/details?id=com.luvvix.onbuch&hl=fr}}%
      \hspace{8pt}\begin{minipage}[c]{0.5\linewidth}
        {\popdisplay\fontsize{11}{12.5}\selectfont\color{white}\MakeUppercase{Télécharge OnBuch}}\par\vspace{4pt}
        {\raggedright\popsemi\fontsize{8.4}{11}\selectfont\color{white}Gratuit \textbullet\ Google Play\\Tuteur IA, annales, concours, résultats\par}
      \end{minipage}
    \end{tcolorbox}
  \end{minipage}\hfill
  \begin{minipage}[t]{0.487\linewidth}
    \begin{tcolorbox}[enhanced,colback=poppurple,colframe=popink,boxrule=2.2pt,arc=10pt,
        drop shadow=popink,left=11pt,right=11pt,top=10pt,bottom=10pt,height=3.1cm,valign=center]
      \tikz[baseline=-0.7ex]\node[circle,fill=white,draw=popink,line width=1.3pt,minimum size=1.6cm,inner sep=0]{\popdisplay\fontsize{24}{24}\selectfont\color{poppurple}+};%
      \hspace{8pt}\begin{minipage}[c]{0.6\linewidth}
        {\popdisplay\fontsize{11}{12.5}\selectfont\color{white}\MakeUppercase{Passe à OnBuch+}}\par\vspace{4pt}
        {\raggedright\popsemi\fontsize{8.4}{11}\selectfont\color{white}Tous les cours signés, Tuteur IA illimité, fiches PDF — onglet OnBuch+ de l'app\par}
      \end{minipage}
    \end{tcolorbox}
  \end{minipage}
  \vspace{8pt}
  \popcontenu
  \vfill
  \signatureonbuch
  \restoregeometry
  \newpage
}

% Rangée « Ce cours contient » (page de garde)
\newcommand{\poptile}[3]{% #1 couleur #2 gros texte #3 légende
  \begin{tcolorbox}[enhanced,colback=#1,colframe=popink,boxrule=2pt,arc=9pt,shadow={2.5pt}{-2.5pt}{0pt}{popink},
      left=6pt,right=6pt,top=9pt,bottom=9pt,halign=center,valign=center,height=1.75cm,width=\linewidth,nobeforeafter]
    {\popdisplay\fontsize{12.5}{13}\selectfont\color{white}\MakeUppercase{#2}}\par\vspace{3pt}
    {\centering\popsemi\fontsize{7.8}{9.5}\selectfont\color{white}#3\par}
  \end{tcolorbox}}
\newcommand{\popcontenu}{%
  \par\noindent{\popxboldls\fontsize{7.5}{7.5}\selectfont\color{popmuted}CE COURS SIGNÉ CONTIENT}\par\vspace{7pt}
  \noindent\begin{minipage}[t]{0.235\linewidth}\poptile{popblue}{Cours}{complet, illustré, pas à pas}\end{minipage}\hfill
  \begin{minipage}[t]{0.235\linewidth}\poptile{poporange}{Méthodes}{et exercices résolus}\end{minipage}\hfill
  \begin{minipage}[t]{0.235\linewidth}\poptile{poppink}{Exercices}{type examen + corrigés}\end{minipage}\hfill
  \begin{minipage}[t]{0.235\linewidth}\poptile{popgreen}{Fiche bilan}{l'essentiel à retenir}\end{minipage}\par}

% Sommaire
\newtcolorbox{sommaire}{enhanced,colback=white,colframe=popink,boxrule=2.2pt,arc=10pt,
  drop shadow=popink,left=18pt,right=18pt,top=13pt,bottom=13pt,before skip=6pt,after skip=20pt,
  before upper={\poppill{Sommaire}{poppurple}{white}\par\vspace{9pt}\popsemi\fontsize{10.6}{14.5}\selectfont\color{popink}}}

% Signature de fin de cours — poussée tout en bas de la dernière page
\newcommand{\findecours}{%
  \par\vfill\nopagebreak
  \begin{center}\popsquiggle{poporange}\end{center}\vspace{4pt}
  \signatureonbuch
}
\AtBeginDocument{\def\llbracket{\mathopen{[\mkern-3.5mu[}}\def\rrbracket{\mathclose{]\mkern-3.5mu]}}} % intervalles d'entiers (glyphe ⟦ absent de la police)
% Glyphes absents de Fira Math : redéfinis à partir de glyphes disponibles
\AtBeginDocument{%
  \def\setminus{\mathbin{\backslash}}\let\smallsetminus\setminus
  \def\vdots{\mathord{\vbox{\baselineskip3.2pt\lineskiplimit0pt\kern4pt\hbox{.}\hbox{.}\hbox{.}}}}%
  \def\ddots{\mathinner{\mkern1mu\raise6.5pt\hbox{.}\mkern2mu\raise3.6pt\hbox{.}\mkern2mu\raise0.7pt\hbox{.}\mkern1mu}}%
  \def\triangle{\mathord{\scalebox{0.9}{$\symup{\Delta}$}}}%
  \def\nabla{\mathord{\raisebox{\depth}{\scalebox{1}[-1]{$\symup{\Delta}$}}}}%
  \def\checkmark{\popcheck{popdark}}%
  \def\star{\mathbin{\ast}}\def\bigstar{\mathord{\ast}}%
  \def\diamond{\mathbin{\scalebox{0.6}{$\Diamond$}}}%
  \def\bigcup{\mathop{\mathchoice{\raisebox{-0.3em}{\scalebox{1.7}{$\cup$}}}{\scalebox{1.25}{$\cup$}}{\cup}{\cup}}\displaylimits}%
  \def\bigcap{\mathop{\mathchoice{\raisebox{-0.3em}{\scalebox{1.7}{$\cap$}}}{\scalebox{1.25}{$\cap$}}{\cap}{\cap}}\displaylimits}%
  \def\lesseqqgtr{\mathrel{\lessgtr}}\def\bigoplus{\mathop{\oplus}\displaylimits}%
}
\providecommand{\ssi}{si et seulement si} % abréviation fréquente
\AtBeginDocument{\providecommand{\wideparen}{\overparen}} % arc de cercle

%% FIGFIX-BEGIN v5 — correctifs globaux des figures (docs/onbuch-generation/tools/figfix)
\usepackage{adjustbox}\usepackage{etoolbox}\tcbuselibrary{raster}
% toute figure TikZ/pgfplots plus large que la ligne est réduite à la largeur disponible
\BeforeBeginEnvironment{tikzpicture}{\begin{adjustbox}{max width=\linewidth}}
\AfterEndEnvironment{tikzpicture}{\end{adjustbox}}
% légendes pgfplots lisibles par-dessus les courbes
\pgfplotsset{every axis/.append style={legend style={fill=white,fill opacity=0.92,text opacity=1,draw=black!15,inner sep=3pt}}}
% étiquettes d'arêtes (syntaxe edge["..."]) avec fond blanc
\tikzset{every edge quotes/.append style={fill=white,inner sep=1.2pt}}
%% FIGFIX-END
\begin{document}

\pagegarde{Sensibilisation sur les effets de certaines substances sur la transmission synaptique}{SVT}{Tle D}{Module I : Le Monde Vivant — Les synapses}{N09}{4 h}{Cette leçon étudie la structure, l'ultrastructure et le fonctionnement des synapses, les mécanismes de sommation qui font du neurone un véritable intégrateur, puis analyse comment certaines substances (drogues, médicaments, toxines) perturbent la transmission synaptique. Elle s'appuie sur des expériences historiques (Loewi) et l'exploitation de documents pour construire les savoirs exigibles au Baccalauréat.
\vspace{4pt}
\begin{multicols}{2}\small
\begin{itemize}
\item À la fin de cette leçon, je sais schématiser et légender la structure et l'ultrastructure d'une synapse chimique.
\item À la fin de cette leçon, je sais distinguer sur des documents les synapses neuro-neuroniques, neuro-musculaires (plaque motrice) et neuro-glandulaires.
\item À la fin de cette leçon, je sais identifier sur des électronographies les synapses à transmission électrique et à transmission chimique.
\item À la fin de cette leçon, je sais interpréter les expériences de Loewi et en déduire la nature chimique de la transmission synaptique.
\item À la fin de cette leçon, je sais expliquer les étapes du fonctionnement d'une synapse chimique et d'une synapse électrique.
\item À la fin de cette leçon, je sais identifier sur des enregistrements les sommations spatiale, temporelle et spatio-temporelle et expliquer le rôle intégrateur du neurone.
\end{itemize}\end{multicols}}

\begin{sommaire}
\begin{multicols}{2}
\begin{enumerate}
\item I. Structure et types anatomiques de synapses
\item II. Nature de la transmission synaptique : l'expérience de Loewi
\item III. Fonctionnement des synapses chimiques et électriques
\item IV. Sommations et rôle intégrateur du neurone
\item V. Effets des substances sur la transmission synaptique
\item VI. Synthèse et sensibilisation
\item Activité d'intégration
\item Méthodes et exercices résolus
\item Exercices
\item Corrigés des exercices
\item Fiche bilan
\end{enumerate}
\end{multicols}
\end{sommaire}

\begin{objectifs}
\begin{itemize}
\item À la fin de cette leçon, je sais schématiser et légender la structure et l'ultrastructure d'une synapse chimique.
\item À la fin de cette leçon, je sais distinguer sur des documents les synapses neuro-neuroniques, neuro-musculaires (plaque motrice) et neuro-glandulaires.
\item À la fin de cette leçon, je sais identifier sur des électronographies les synapses à transmission électrique et à transmission chimique.
\item À la fin de cette leçon, je sais interpréter les expériences de Loewi et en déduire la nature chimique de la transmission synaptique.
\item À la fin de cette leçon, je sais expliquer les étapes du fonctionnement d'une synapse chimique et d'une synapse électrique.
\item À la fin de cette leçon, je sais identifier sur des enregistrements les sommations spatiale, temporelle et spatio-temporelle et expliquer le rôle intégrateur du neurone.
\item À la fin de cette leçon, je sais expliquer le mode d'action de certaines substances (curare, atropine, cocaïne, nicotine, alcool, insecticides...) sur la transmission synaptique.
\item À la fin de cette leçon, je sais exploiter des résultats expérimentaux et des courbes pour argumenter sur les dangers des drogues.
\end{itemize}
\end{objectifs}

\begin{prerequis}
\begin{itemize}
\item Le neurone : structure (corps cellulaire, dendrites, axone) et rôle dans la conduction de l'influx nerveux (classe de Première).
\item Le message nerveux : potentiel de repos et potentiel d'action (notions vues en début d'année de Terminale).
\item Notion de réflexe et d'arc réflexe (Première).
\item Notions de membrane plasmique, de perméabilité ionique et de transport actif (Première).
\item Lecture et interprétation de courbes et de tableaux de résultats expérimentaux.
\end{itemize}
\end{prerequis}

\newpage

\coursec{I. Structure et types anatomiques de synapses}

\textbf{Situation d'accroche.} À Yaoundé, un jeune mototaximan consomme régulièrement du tramadol pour « tenir » les longues journées de travail ; un soir, il est admis en urgence à l'hôpital central, inconscient, la respiration ralentie. Au marché Mokolo, on vend des cigarettes et de l'alcool que beaucoup de jeunes élèves testent par curiosité. Pourquoi ces substances modifient-elles le comportement, les réflexes, voire la conscience ? Comment une simple molécule peut-elle bloquer ou amplifier le dialogue entre nos neurones ?

\textbf{Problématique.} Cette leçon nous conduit au cœur des \cle{synapses}, ces zones de contact où s'échangent les messages nerveux, pour comprendre comment drogues, médicaments et toxines perturbent leur fonctionnement. Mais avant d'étudier ces perturbations, il faut d'abord répondre à une question plus fondamentale : comment les neurones sont-ils organisés les uns par rapport aux autres, et à quoi ressemblent leurs zones de contact ?

\courssub{Mise en évidence des synapses : le neurone, une cellule indépendante}

\textbf{Observation.} Au microscope optique, une coupe de moelle épinière colorée (par exemple à l'hématéline-éosine ou par imprégnation argentique) montre des corps cellulaires de neurones étoilés, munis de longs prolongements. En suivant ces prolongements à fort grossissement, on constate un fait capital : les extrémités des axones d'un neurone viennent s'appliquer contre le corps cellulaire ou les dendrites d'un autre neurone, \emph{sans jamais fusionner avec lui}. Il existe toujours une discontinuité, une frontière entre les deux cellules.

\textbf{Un débat historique éclairant.} À la fin du XIX\textsuperscript{e} siècle, deux écoles s'affrontent. Pour Camillo Golgi, le système nerveux forme un \emph{réseau continu} : les prolongements des neurones seraient soudés les uns aux autres, comme les fils d'un filet. Pour Santiago Ramón y Cajal, qui utilise la méthode de coloration inventée par Golgi lui-même (imprégnation au nitrate d'argent), chaque neurone est une \emph{cellule indépendante}, une unité anatomique et fonctionnelle : c'est la \cle{théorie neuronale}. L'observation attentive de milliers de préparations donne raison à Cajal : les neurones se touchent mais ne se confondent pas. Le microscope électronique, au XX\textsuperscript{e} siècle, tranchera définitivement en montrant la membrane de chaque cellule et l'espace qui les sépare.

\begin{savaistu}[Golgi et Cajal, Nobel 1906]
Camillo Golgi et Santiago Ramón y Cajal ont reçu \emph{ensemble} le prix Nobel de physiologie ou médecine en 1906, alors qu'ils défendaient des théories opposées ! Golgi avait inventé la coloration, Cajal l'avait utilisée pour démontrer que les neurones sont des cellules indépendantes (théorie neuronale). C'est un bel exemple de démarche scientifique : une technique nouvelle permet de trancher un débat théorique par l'observation.
\end{savaistu}

\textbf{Interprétation.} Si les neurones sont des cellules indépendantes, alors le message nerveux doit être \emph{transmis} d'une cellule à l'autre au niveau de zones de contact spécialisées. Ces zones de contact portent un nom : les synapses (du grec \emph{sunapsis}, « jonction »).

\begin{definition}[Synapse]
Une \cle{synapse} est la zone de contact fonctionnel entre un neurone et une cellule cible (un autre neurone, une cellule musculaire ou une cellule glandulaire), où le message nerveux est transmis d'une cellule à l'autre. On distingue l'\cle{élément présynaptique} (le neurone qui émet le message) et l'\cle{élément postsynaptique} (la cellule qui le reçoit).
\end{definition}

\begin{attention}[Le neurone n'est pas un fil continu]
Erreur fréquente : croire que l'influx nerveux « passe » d'un neurone à l'autre comme le courant dans un fil électrique. Non ! Chaque neurone est une cellule close, limitée par sa membrane. La transmission synaptique est une \emph{étape de communication} entre deux cellules distinctes, avec des mécanismes spécifiques que nous étudierons.
\end{attention}

\courssub{Structure d'une synapse chimique vue au microscope optique}

Au microscope optique, l'extrémité de l'axone présynaptique se ramifie en fines arborisations dont chaque ramification se termine par un renflement appelé \cle{bouton synaptique} (ou terminaison axonique). Ce bouton, d'un diamètre de l'ordre de \SI{1}{\micro\metre}, vient s'appliquer contre la membrane de l'élément postsynaptique : corps cellulaire ou dendrite d'un neurone, fibre musculaire, cellule glandulaire. Un même neurone postsynaptique peut recevoir des centaines, voire des milliers de boutons synaptiques issus de neurones différents : on comprend déjà qu'il pourra « intégrer » de nombreuses informations (nous y reviendrons à la section IV).

Mais le microscope optique ne permet pas de voir ce qui se passe \emph{à l'intérieur} du bouton ni \emph{entre} les deux cellules. Pour cela, il faut le microscope électronique.

\courssub{Ultrastructure d'une synapse chimique au microscope électronique}

\textbf{Observation (document type électronographie).} Une coupe ultrafine d'une synapse, observée au microscope électronique à transmission, révèle une organisation très ordonnée :

\begin{itemize}
\item le \cle{bouton synaptique} contient de nombreuses \cle{vésicules synaptiques} (petits sacs membranaires de \SIrange{30}{50}{\nano\metre} de diamètre) et des \cle{mitochondries}, signe d'une intense activité énergétique ;
\item la membrane du bouton, en regard de la cellule cible, est épaissie : c'est la \cle{membrane présynaptique} ;
\item entre les deux cellules existe un espace intercellulaire étroit et constant, la \cle{fente synaptique}, large de \SIrange{20}{40}{\nano\metre} ;
\item la membrane de la cellule cible, en regard du bouton, est la \cle{membrane postsynaptique} ; elle porte des \cle{récepteurs spécifiques} (protéines membranaires) capables de fixer les substances libérées par l'élément présynaptique.
\end{itemize}

\begin{popfigure}
\begin{center}
\begin{tikzpicture}[scale=1.0, every node/.style={font=\small}]
% Bouton synaptique
\fill[poporangeL] (-3.2,0.6) .. controls (-3.4,2.6) and (-1.5,3.6) .. (0,3.6)
  .. controls (1.5,3.6) and (3.4,2.6) .. (3.2,0.6) .. controls (1.5,0.2) and (-1.5,0.2) .. (-3.2,0.6);
\draw[popdark, thick] (-3.2,0.6) .. controls (-3.4,2.6) and (-1.5,3.6) .. (0,3.6)
  .. controls (1.5,3.6) and (3.4,2.6) .. (3.2,0.6);
% Axone
\draw[popdark, thick] (-0.8,3.6) -- (-0.8,4.6);
\draw[popdark, thick] (0.8,3.6) -- (0.8,4.6);
\node[popdark] at (0,4.9) {Axone présynaptique};
% Vésicules
\foreach \x/\y in {-1.8/2.6, -0.9/3.0, 0.1/2.7, 1.1/3.0, 1.9/2.5, -0.4/2.1, 0.7/2.2, -1.3/1.9, 1.5/1.8}{
  \draw[popblue, fill=popblueL] (\x,\y) circle (0.16);}
% Mitochondries
\draw[popgreen, fill=popgreenL, thick] (-2.2,1.2) ellipse (0.55 and 0.28);
\draw[popgreen] (-2.55,1.2) .. controls (-2.3,1.35) and (-2.1,1.05) .. (-1.85,1.2);
\draw[popgreen, fill=popgreenL, thick] (2.3,1.1) ellipse (0.55 and 0.28);
\draw[popgreen] (1.95,1.1) .. controls (2.2,1.25) and (2.4,0.95) .. (2.65,1.1);
% Membrane présynaptique
\draw[poporange, line width=2.2pt] (-3.2,0.6) .. controls (-1.5,0.2) and (1.5,0.2) .. (3.2,0.6);
% Fente synaptique
\fill[popcream] (-3.4,-0.15) rectangle (3.4,0.35);
\draw[<->, popmuted, thick] (3.55,0.55) -- (3.55,-0.6);
\node[popmuted, right, align=left] at (3.65,-0.05) {\SIrange{20}{40}{\nano\metre}};
% Membrane postsynaptique
\draw[poppurple, line width=2.2pt] (-3.4,-0.6) -- (3.4,-0.6);
% Récepteurs
\foreach \x in {-2.4,-1.2,0,1.2,2.4}{
  \draw[poppurple, fill=poppinkL, thick] (\x,-0.6) ++(-0.13,-0.28) rectangle ++(0.26,0.28);
  \draw[poppurple, thick] (\x,-0.6) -- (\x,-0.32);}
% Cellule postsynaptique
\fill[poppurpleL, opacity=0.5] (-3.4,-2.6) rectangle (3.4,-0.9);
\draw[popdark, thick] (-3.4,-0.9) -- (3.4,-0.9);
\node[popdark] at (0,-1.8) {Élément postsynaptique (dendrite, corps cellulaire, muscle, glande)};
% Légendes fléchées
\node[popblue, right] at (4.1,2.7) {1. Vésicules synaptiques};
\draw[->, popblue] (4.05,2.7) -- (2.1,2.6);
\node[popgreen, right] at (4.1,1.5) {2. Mitochondrie};
\draw[->, popgreen] (4.05,1.5) -- (2.9,1.15);
\node[poporange, right] at (4.1,0.75) {3. Membrane présynaptique};
\draw[->, poporange] (4.05,0.7) -- (2.6,0.4);
\node[popmuted, left, align=right] at (-4.1,-0.05) {4. Fente\\ synaptique};
\draw[->, popmuted] (-4.0,-0.05) -- (-2.9,0.1);
\node[poppurple, left, align=right] at (-4.1,-1.1) {5. Récepteurs\\ postsynaptiques};
\draw[->, poppurple] (-4.0,-1.1) -- (-2.5,-0.75);
\node[popdark, font=\small\itshape] at (-1.6,1.15) {Bouton synaptique};
\end{tikzpicture}
\end{center}
\legende{Ultrastructure d'une synapse chimique vue au microscope électronique. 1 : vésicules synaptiques contenant le neuromédiateur ; 2 : mitochondries fournissant l'énergie ; 3 : membrane présynaptique ; 4 : fente synaptique (\SIrange{20}{40}{\nano\metre}) ; 5 : récepteurs spécifiques de la membrane postsynaptique.}
\end{popfigure}

\begin{attention}[La fente synaptique n'est pas un vide artéfactuel]
Erreur très fréquente au Bac : présenter la fente synaptique comme un « trou » ou un artefact de préparation. C'est un \emph{espace réel} de \SIrange{20}{40}{\nano\metre}, rempli de liquide interstitiel. Conséquence capitale : dans une synapse chimique, le message \emph{ne franchit pas} cet espace sous forme électrique ; il le franchit sous forme de molécules chimiques (neuromédiateurs) qui diffusent dans la fente.
\end{attention}

\textbf{Le contenu des vésicules (savoir admis).} Les vésicules synaptiques contiennent des molécules synthétisées par le neurone présynaptique : les \cle{neuromédiateurs}. Les principaux sont l'\cle{acétylcholine}, la \cle{noradrénaline}, la \cle{dopamine}, la \cle{sérotonine} et le \cle{GABA} (acide gamma-aminobutyrique). Chaque type de neurone libère en général un neuromédiateur déterminé.

\begin{definition}[Neuromédiateur]
Un \cle{neuromédiateur} est une substance chimique synthétisée par le neurone présynaptique, stockée dans les vésicules synaptiques, libérée dans la fente synaptique et capable d'agir sur des récepteurs spécifiques de la membrane postsynaptique pour y modifier l'activité électrique.
\end{definition}

\textbf{Interprétation fonctionnelle.} La présence de nombreuses mitochondries dans le bouton synaptique s'explique : la synthèse des neuromédiateurs, leur libération et leur recyclage sont des processus qui consomment beaucoup d'ATP. La polarisation de la synapse (vésicules d'un côté, récepteurs de l'autre) explique pourquoi la transmission synaptique est \emph{unidirectionnelle} : le message ne peut passer que de l'élément présynaptique vers l'élément postsynaptique.

\begin{methode}[Schématiser une synapse au Bac]
Pour obtenir tous les points lors d'une schématisation de synapse :
\begin{enumerate}
\item Représenter le \textbf{bouton synaptique} (renflement de l'axone) contenant des \textbf{vésicules} (petits cercles) et au moins une \textbf{mitochondrie} ;
\item Tracer la \textbf{membrane présynaptique} et, en regard, la \textbf{membrane postsynaptique}, séparées par la \textbf{fente synaptique} (espace blanc étroit) ;
\item Figurer des \textbf{récepteurs} sur la membrane postsynaptique (petites protéines enchâssées) ;
\item Légender proprement chaque structure avec des traits de légende qui ne se croisent pas, et donner un titre au schéma.
\end{enumerate}
Ne jamais représenter les vésicules dans l'élément postsynaptique : c'est une erreur qui coûte des points !
\end{methode}

\courssub{Les types anatomiques de synapses}

Selon la nature des deux éléments en contact, on distingue trois grands types anatomiques de synapses.

\textbf{a) Les synapses neuro-neuroniques.} Elles relient deux neurones. Selon la partie du neurone postsynaptique contactée, on parle de synapses :
\begin{itemize}
\item \cle{axo-dendritiques} : le bouton synaptique s'applique sur une dendrite (cas le plus fréquent) ;
\item \cle{axo-somatiques} : le bouton s'applique directement sur le corps cellulaire ;
\item \cle{axo-axoniques} : le bouton s'applique sur l'axone (souvent sur le segment initial) d'un autre neurone, ce qui permet de moduler sa libération de neuromédiateur.
\end{itemize}

\textbf{b) Les synapses neuro-musculaires : la plaque motrice.} Au niveau d'un muscle squelettique (par exemple les muscles du bras d'un porteur de l'école primaire de Mvog-Ada), l'axone du motoneurone se divise en une \cle{arborisation terminale} dont chaque rameau innerve une fibre musculaire. La zone de contact, appelée \cle{plaque motrice} (ou jonction neuro-musculaire), présente une organisation remarquable : l'extrémité axonique repose dans des \cle{gouttières synaptiques} creusées à la surface de la fibre musculaire, et la membrane musculaire forme de profonds \cle{replis sous-neuraux} (jonctions sous-neurales) dont les crêtes sont très riches en \cle{récepteurs à l'acétylcholine}. Cette organisation en gouttières et replis augmente considérablement la surface de réception, garantissant l'efficacité de la transmission : chaque influx nerveux déclenche normalement une contraction.

\textbf{c) Les synapses neuro-glandulaires.} Elles relient un neurone (souvent du système nerveux autonome) à une \cle{cellule glandulaire} : par exemple les cellules de la médullo-surrénale qui libèrent l'adrénaline, ou les cellules des glandes salivaires. Le bouton synaptique libère son neuromédiateur au voisinage de la cellule glandulaire, dont la membrane porte les récepteurs correspondants ; la cellule répond en sécrétant son produit (hormone, salive, sueur...).

\begin{popfigure}
\begin{center}
\begin{tikzpicture}[scale=0.92, every node/.style={font=\small}]
% ===== 1. Neuro-neuronique =====
\begin{scope}[shift={(-5.2,0)}]
\node[popdark, font=\small\bfseries] at (0,3.4) {a) Neuro-neuronique};
% axone présynaptique
\draw[poporange, very thick] (-1.6,2.6) -- (0.2,1.6);
\fill[poporangeL, draw=popdark] (0.2,1.6) circle (0.28);
% corps cellulaire postsynaptique
\fill[poppurpleL, draw=popdark, thick] (0.6,0.2) circle (0.75);
\fill[poppurple, draw=popdark] (0.6,0.2) circle (0.25);
% dendrite
\draw[popdark, thick] (0.1,0.85) -- (-0.3,1.5);
\draw[popdark, thick] (1.1,0.85) -- (1.5,1.5);
\node[popmuted, align=center] at (0,-1.2) {bouton sur dendrite\\ ou corps cellulaire\\ d'un neurone};
\end{scope}
% ===== 2. Neuro-musculaire =====
\begin{scope}[shift={(0,0)}]
\node[popdark, font=\small\bfseries] at (0,3.4) {b) Neuro-musculaire};
% axone moteur
\draw[poporange, very thick] (-1.8,2.8) -- (-0.2,1.8);
\draw[poporange, thick] (-0.2,1.8) -- (0.5,1.3);
\draw[poporange, thick] (-0.2,1.8) -- (1.1,1.5);
\fill[poporangeL, draw=popdark] (0.5,1.3) circle (0.16);
\fill[poporangeL, draw=popdark] (1.1,1.5) circle (0.16);
% fibre musculaire avec replis
\draw[popgreen, fill=popgreenL, thick] (-1.4,0.9) -- (1.9,0.9)
  .. controls (1.95,0.7) .. (1.9,0.5) -- (-1.4,0.5) -- cycle;
\foreach \x in {0.3,0.6,0.9,1.2}{
  \draw[popgreen, thick] (\x,0.9) -- (\x,1.12);}
\node[popmuted, align=center] at (0.3,-1.2) {plaque motrice : gouttières\\ et replis riches en\\ récepteurs à l'acétylcholine};
\node[popgreen, font=\footnotesize\itshape] at (0.3,0.28) {fibre musculaire};
\end{scope}
% ===== 3. Neuro-glandulaire =====
\begin{scope}[shift={(5.2,0)}]
\node[popdark, font=\small\bfseries] at (0,3.4) {c) Neuro-glandulaire};
% axone
\draw[poporange, very thick] (-1.6,2.6) -- (0,1.5);
\fill[poporangeL, draw=popdark] (0,1.5) circle (0.24);
% cellule glandulaire
\fill[popblueL, draw=popdark, thick] (0.5,0.2) circle (0.7);
\foreach \x/\y in {0.3/0.4, 0.7/0.1, 0.4/-0.1}{
  \fill[popblue] (\x,\y) circle (0.08);}
\node[popmuted, align=center] at (0,-1.2) {bouton au contact\\ d'une cellule\\ glandulaire};
\end{scope}
\end{tikzpicture}
\end{center}
\legende{Les trois types anatomiques de synapses. a) Synapse neuro-neuronique (axo-dendritique ou axo-somatique) ; b) synapse neuro-musculaire (plaque motrice) : l'arborisation terminale de l'axone moteur repose dans des gouttières synaptiques, la membrane musculaire forme des replis sous-neuraux riches en récepteurs à l'acétylcholine ; c) synapse neuro-glandulaire.}
\end{popfigure}

\begin{exemplebox}[La plaque motrice, une synapse à haut rendement]
Quand tu lèves la main pour répondre à une question, l'ordre part du cortex cérébral, descend le long de la moelle épinière, emprunte un motoneurone dont l'axone atteint les muscles du bras. Au niveau de chaque plaque motrice, l'acétylcholine libérée se fixe sur les récepteurs des replis sous-neuraux et déclenche la contraction des fibres musculaires. Une seule fibre musculaire ne possède qu'\emph{une} plaque motrice, mais un motoneurone peut innerver de quelques fibres (muscles des doigts, précision fine) à plus de mille fibres (muscles de la cuisse, puissance).
\end{exemplebox}

\begin{exoresolu}[Identifier les éléments d'une synapse sur un document]
\textbf{Énoncé.} Sur une électronographie de synapse, on observe : une expansion terminale d'axone contenant des vésicules de \SI{40}{\nano\metre} et des mitochondries ; un espace intercellulaire de \SI{30}{\nano\metre} ; une membrane sous-jacente épaissie appartenant à une dendrite. 1) Nommer les structures observées. 2) Préciser le type anatomique de cette synapse. 3) Expliquer pourquoi la transmission y est unidirectionnelle.
\tcblower
\textbf{Solution détaillée.}
1) L'expansion terminale est le \textbf{bouton synaptique} ; les vésicules sont les \textbf{vésicules synaptiques} (contenant le neuromédiateur) ; l'espace de \SI{30}{\nano\metre} est la \textbf{fente synaptique} (valeur conforme aux \SIrange{20}{40}{\nano\metre} attendus) ; la membrane sous-jacente est la \textbf{membrane postsynaptique} porteuse de récepteurs.

2) L'élément postsynaptique étant une \emph{dendrite}, il s'agit d'une synapse \textbf{neuro-neuronique axo-dendritique}.

3) La transmission est unidirectionnelle car les \textbf{vésicules de neuromédiateur} ne se trouvent que du côté présynaptique et les \textbf{récepteurs spécifiques} que du côté postsynaptique : cette polarisation structurale impose un sens unique de fonctionnement, de l'élément présynaptique vers l'élément postsynaptique.
\end{exoresolu}

\begin{aretenir}[L'essentiel de la section I]
\begin{itemize}
\item Les neurones sont des cellules \textbf{indépendantes} (théorie neuronale de Cajal) qui communiquent au niveau de zones de contact : les \textbf{synapses}.
\item Une synapse chimique comprend : le \textbf{bouton synaptique} (vésicules + mitochondries), la \textbf{membrane présynaptique}, la \textbf{fente synaptique} (\SIrange{20}{40}{\nano\metre}) et la \textbf{membrane postsynaptique} porteuse de \textbf{récepteurs}.
\item Les vésicules contiennent un \textbf{neuromédiateur} (acétylcholine, dopamine, sérotonine, GABA...).
\item Trois types anatomiques : \textbf{neuro-neuroniques} (axo-dendritiques, axo-somatiques, axo-axoniques), \textbf{neuro-musculaires} (plaque motrice) et \textbf{neuro-glandulaires}.
\end{itemize}
\end{aretenir}

Maintenant que la structure des synapses est connue, une question décisive se pose : le message nerveux franchit-il la fente synaptique sous forme électrique ou sous forme chimique ? C'est l'expérience historique d'Otto Loewi, étudiée dans la section suivante, qui apportera la réponse.

\coursec{II. Nature de la transmission synaptique : l'expérience de Loewi}

Dans la section précédente, nous avons décrit la structure des synapses : entre le bouton synaptique du neurone présynaptique et la membrane postsynaptique, il existe une \cle{fente synaptique} étroite (environ 20 nm), que les membranes ne franchissent pas. Une question fondamentale se pose alors immédiatement : comment le message nerveux, qui voyage sous forme électrique le long des fibres nerveuses, peut-il \emph{traverser} cet espace ? La réponse à cette question n'est pas tombée du ciel : elle a été établie par une expérience élégante, devenue un grand classique du Baccalauréat, l'\cle{expérience de Loewi} (1921). C'est elle que nous allons analyser pas à pas, en suivant la démarche d'investigation scientifique.

\courssub{Le problème posé : transmission électrique ou chimique ?}

Au début du XX\textsuperscript{e} siècle, les physiologistes savent que l'influx nerveux est un phénomène de nature électrique (variations du potentiel de membrane). Mais au niveau de la synapse, deux hypothèses s'affrontent :

\begin{itemize}
\item \textbf{Hypothèse 1 — transmission électrique :} le courant électrique passerait directement de la cellule présynaptique à la cellule postsynaptique, comme le courant passe d'un fil à un autre.
\item \textbf{Hypothèse 2 — transmission chimique :} l'arrivée de l'influx provoquerait la libération, par l'élément présynaptique, d'une \cle{substance chimique} (un \cle{neuromédiateur}) qui diffuserait dans la fente synaptique et agirait sur l'élément postsynaptique.
\end{itemize}

Pour trancher entre ces deux hypothèses, il faut un dispositif expérimental capable de \emph{séparer} l'effet nerveux de tout contact direct. C'est exactement ce que réalise Otto Loewi, physiologiste autrichien, en 1921, en travaillant sur un matériel bien choisi : le \cle{cœur de grenouille}.

\begin{methode}[Interpréter une expérience au Bac : la démarche à suivre]
Face à tout document expérimental (au Baccalauréat comme en TP), rédige ton analyse en suivant toujours les six étapes de la démarche d'investigation scientifique :
\begin{enumerate}
\item \textbf{Problème} : quelle question l'expérience cherche-t-elle à résoudre ?
\item \textbf{Hypothèses} : quelles réponses sont envisageables ?
\item \textbf{Protocole} : que fait l'expérimentateur, et quelles sont les variables contrôlées ?
\item \textbf{Résultats} : que constate-t-on, de façon objective (sans interpréter) ?
\item \textbf{Interprétation} : que signifient ces résultats, quelle hypothèse est validée ou rejetée ?
\item \textbf{Conclusion} : quelle notion générale peut-on dégager ?
\end{enumerate}
Une erreur fréquente consiste à confondre \emph{résultat} (ce qu'on observe) et \emph{interprétation} (ce qu'on en déduit). Le correcteur attend les deux, clairement séparés.
\end{methode}

\courssub{Le dispositif expérimental de Loewi (1921)}

\textbf{Matériel biologique.} Le cœur de grenouille présente un avantage décisif : prélevé et placé dans du \cle{liquide de Ringer} (solution physiologique dont la composition en sels minéraux — \ce{NaCl}, \ce{KCl}, \ce{CaCl2} — maintient les cellules en survie), il continue de battre \emph{spontanément} pendant plusieurs heures, car son rythme est automatique (il naît dans le tissu cardiaque lui-même). De plus, l'innervation du cœur de grenouille est simple : le \cle{nerf vague} (nerf parasympathique, aussi appelé nerf pneumogastrique ou nerf X) \emph{ralentit} les battements cardiaques lorsqu'on le stimule.

\textbf{Le dispositif.} Loewi utilise \emph{deux} cœurs de grenouille :

\begin{itemize}
\item le \textbf{cœur A} est prélevé \emph{avec son nerf vague} et perfusé avec du liquide de Ringer ;
\item le \textbf{cœur B} est prélevé \emph{sans son nerf vague} : on dit qu'il est \cle{dénervé}. Il est placé dans un second récipient contenant du liquide de Ringer.
\end{itemize}

\begin{experience}[L'expérience de Loewi : protocole et observations]
\textbf{Problème :} le message nerveux franchit-il la jonction nerf-cœur sous forme électrique ou sous forme chimique ?

\textbf{Protocole :}
\begin{enumerate}
\item On enregistre la fréquence de battement des deux cœurs : elle est normale et comparable (environ 40 à 60 battements par minute).
\item On stimule électriquement le \textbf{nerf vague du cœur A}.
\item On recueille le \textbf{liquide de perfusion du cœur A} et on le transfère dans le récipient contenant le \textbf{cœur B dénervé}.
\end{enumerate}

\textbf{Observations (résultats) :}
\begin{enumerate}
\item Dès la stimulation de son nerf vague, le \textbf{cœur A ralentit} : c'est le résultat attendu, le nerf vague est cardiomodérateur.
\item Quelques instants après le transfert du liquide, le \textbf{cœur B ralentit à son tour}, alors qu'il n'a reçu \emph{aucune stimulation nerveuse} (il est dénervé) et \emph{aucun contact} avec le cœur A : seul du liquide a été transféré.
\end{enumerate}

\textbf{Interprétation :} la stimulation du nerf vague a provoqué, au niveau de la jonction entre le nerf et le muscle cardiaque, la \textbf{libération d'une substance chimique} dans le liquide de perfusion. Cette substance, transportée par le liquide, a atteint le cœur B et y a reproduit l'effet du nerf vague : le ralentissement. Loewi nomma cette substance \emph{Vagusstoff} (« substance du vague ») ; elle fut identifiée quelques années plus tard comme étant l'\cle{acétylcholine}.

\textbf{Conclusion :} au niveau de la jonction nerf-cœur, la transmission du message nerveux est de \textbf{nature chimique} : elle s'effectue par libération d'un neuromédiateur.
\end{experience}

\begin{popfigure}
\begin{tikzpicture}[scale=0.92, every node/.style={font=\small}]
% Coeur A
\draw[fill=poppinkL, draw=poppink, thick] (0,0.4) ellipse (1.1 and 0.9);
\node[font=\bfseries] at (0,0.4) {C\oe ur A};
\node[font=\scriptsize] at (0,-0.15) {innervé};
% nerf vague
\draw[very thick, poppurple] (-2.8,1.6) -- (-1.0,1.1);
\node[poppurple, font=\scriptsize, align=center] at (-2.6,1.95) {nerf vague\\ (stimulé)};
\draw[->, very thick, poporange] (-2.2,2.6) -- (-2.2,1.75);
\node[poporange, font=\scriptsize] at (-2.2,2.85) {stimulation};
% liquide A
\draw[popblue!50, fill=popblueL] (-1.5,-0.9) rectangle (1.5,-1.5);
\node[font=\scriptsize, popblue] at (0,-1.2) {liquide de Ringer A};
% Coeur B
\draw[fill=poppinkL, draw=poppink, thick] (6,0.4) ellipse (1.1 and 0.9);
\node[font=\bfseries] at (6,0.4) {C\oe ur B};
\node[font=\scriptsize] at (6,-0.15) {dénervé};
% liquide B
\draw[popblue!50, fill=popblueL] (4.5,-0.9) rectangle (7.5,-1.5);
\node[font=\scriptsize, popblue] at (6,-1.2) {liquide de Ringer B};
% transfert
\draw[->, very thick, popgreen] (1.6,-1.2) .. controls (3,-2.1) .. (4.4,-1.2);
\node[popgreen, font=\scriptsize, align=center] at (3,-2.25) {transfert du liquide\\ (contenant l'acétylcholine)};
% effets
\node[font=\scriptsize, align=center, popdark] at (0,1.7) {ralentissement\\ immédiat};
\node[font=\scriptsize, align=center, popdark] at (6,1.7) {ralentissement\\ retardé};
\end{tikzpicture}
\legende{Dispositif de l'expérience de Loewi (1921). Le cœur A, prélevé avec son nerf vague, et le cœur B, dénervé, sont maintenus en survie dans du liquide de Ringer. La stimulation du nerf vague du cœur A le ralentit ; le transfert du liquide de perfusion de A vers B provoque, avec un léger retard, le ralentissement du cœur B : preuve de la libération d'une substance chimique (l'acétylcholine) par les terminaisons nerveuses.}
\end{popfigure}

\courssub{Analyse quantitative des résultats}

Pour bien comprendre la logique de l'expérience, il est utile de représenter l'évolution des fréquences cardiaques des deux cœurs au cours du temps. Avant la stimulation, les deux cœurs battent à une fréquence de base voisine (environ \num{50} battements par minute). Après stimulation du nerf vague du cœur A, la fréquence du cœur A chute rapidement ; celle du cœur B ne chute qu'\emph{après le transfert du liquide}, avec un retard correspondant au temps de diffusion de la substance.

\begin{popfigure}
\begin{tikzpicture}
\begin{axis}[
width=0.85\linewidth, height=6.2cm,
xlabel={Temps (min)}, ylabel={Fréquence cardiaque (battements/min)},
xmin=0, xmax=12, ymin=0, ymax=60,
xtick={0,2,4,6,8,10,12},
ytick={0,10,20,30,40,50,60},
legend style={at={(0.98,0.97)}, anchor=north east, font=\scriptsize},
grid=both, grid style={popline!40},
]
% Coeur A
\addplot[very thick, poppink, smooth] coordinates {(0,50) (2,50) (3,50) (3.5,38) (4,28) (5,24) (6,24) (8,26) (10,32) (12,38)};
\addlegendentry{C\oe ur A (nerf vague stimulé à $t = 3$ min)}
% Coeur B
\addplot[very thick, popblue, smooth, dashed] coordinates {(0,50) (2,50) (4,50) (5,50) (6,50) (6.5,42) (7,32) (8,27) (9,27) (10,29) (12,34)};
\addlegendentry{C\oe ur B (transfert du liquide à $t = 6$ min)}
% lignes verticales
\draw[poporange, thick, dashed] (axis cs:3,0) -- (axis cs:3,58);
\draw[popgreen, thick, dashed] (axis cs:6,0) -- (axis cs:6,58);
\node[poporange, font=\scriptsize, anchor=west] at (axis cs:3.1,56) {stimulation};
\node[popgreen, font=\scriptsize, anchor=west] at (axis cs:6.1,56) {transfert};
\end{axis}
\end{tikzpicture}
\legende{Évolution comparée des fréquences cardiaques des cœurs A et B au cours de l'expérience de Loewi (valeurs indicatives). Le cœur A ralentit dès la stimulation de son nerf vague ($t = 3$ min). Le cœur B, dénervé, ne ralentit qu'après le transfert du liquide de perfusion de A ($t = 6$ min), avec un retard dû au temps de diffusion de l'acétylcholine.}
\end{popfigure}

\begin{attention}[Le point de contrôle capital : le cœur B est dénervé]
Le ralentissement du cœur B ne peut \emph{pas} s'expliquer par une stimulation nerveuse directe, puisque son nerf vague a été sectionné. Il ne peut pas non plus s'expliquer par un courant électrique venant du cœur A, puisque les deux cœurs ne sont en \emph{aucun contact} : seul du liquide a été transféré. La seule explication possible est donc qu'une \textbf{substance chimique}, libérée dans le liquide du cœur A lors de la stimulation nerveuse, a été véhiculée jusqu'au cœur B. C'est ce contrôle expérimental qui rend la démonstration de Loewi irréfutable. Au Bac, si tu oublies de mentionner que le cœur B est dénervé, ton interprétation perd toute sa valeur logique.
\end{attention}

\courssub{Conclusion et généralisation : la transmission chimique}

\begin{propriete}[Nature chimique de la transmission synaptique]
Au niveau des \cle{synapses chimiques}, l'arrivée de l'influx nerveux au bouton synaptique présynaptique provoque la \textbf{libération d'un neuromédiateur} (substance chimique) dans la fente synaptique. Ce neuromédiateur diffuse, se fixe sur des récepteurs spécifiques de la membrane postsynaptique et y déclenche une réponse (excitation ou inhibition). Cette propriété, établie expérimentalement par Loewi en 1921 pour la jonction nerf-cœur, a été généralisée ensuite à la quasi-totalité des synapses de l'organisme (synapses neuro-neuroniques, neuro-musculaires, neuro-glandulaires).
\end{propriete}

\textbf{Expérience complémentaire (savoir admis).} Le cœur reçoit aussi une innervation \emph{sympathique}, antagoniste de celle du nerf vague. En stimulant le \cle{nerf sympathique} d'un cœur de grenouille, on observe une \emph{accélération} des battements ; le transfert du liquide de perfusion reproduit cette accélération sur un cœur dénervé. La substance libérée est ici la \cle{noradrénaline}. Il existe donc :

\begin{itemize}
\item des \textbf{neuromédiateurs inhibiteurs} (comme l'acétylcholine sur le cœur, qui le ralentit) ;
\item des \textbf{neuromédiateurs excitateurs} (comme la noradrénaline sur le cœur, qui l'accélère).
\end{itemize}

Remarque importante : un même neuromédiateur peut être excitateur sur un effecteur et inhibiteur sur un autre. Ainsi, l'acétylcholine \emph{excite} le muscle squelettique au niveau de la plaque motrice (elle déclenche la contraction), alors qu'elle \emph{inhibe} le cœur. Tout dépend du type de récepteur présent sur la membrane postsynaptique.

\begin{exemplebox}[Quelques neuromédiateurs à connaître]
\begin{itemize}
\item L'\textbf{acétylcholine} est le neuromédiateur de la \cle{plaque motrice} (jonction neuro-musculaire des muscles squelettiques) et des synapses du système nerveux \textbf{parasympathique} (dont le nerf vague sur le cœur).
\item La \textbf{noradrénaline} est le neuromédiateur de nombreuses synapses du système nerveux \textbf{sympathique} (accélération cardiaque, dilatation des bronches, etc.).
\item D'autres neuromédiateurs seront évoqués dans la suite du cours : la dopamine, la sérotonine, le GABA, le glutamate — c'est précisément sur leurs synapses qu'agissent de nombreuses drogues et toxines (section V).
\end{itemize}
\end{exemplebox}

\begin{savaistu}[Une expérience née d'un rêve]
Otto Loewi raconta que l'idée de cette expérience lui était venue \emph{en rêve}, dans la nuit de Pâques 1921. Réveillé en sursaut, il griffonna quelques notes, mais ne parvint pas à les relire le lendemain ! Le rêve se reproduisit la nuit suivante : cette fois, il se leva immédiatement, se rendit à son laboratoire et réalisa l'expérience avant l'aube. Quinze ans plus tard, en 1936, il reçut le \textbf{prix Nobel de physiologie ou médecine} (conjointement avec Henry Dale) pour cette découverte de la transmission chimique de l'influx nerveux. Une belle leçon : en science, l'intuition compte, mais seule l'expérience rigoureuse transforme une idée en savoir.
\end{savaistu}

\courssub{Limites de la généralisation : il existe aussi des synapses électriques}

L'expérience de Loewi démontre la nature chimique de la transmission au niveau de la jonction nerf-cœur, et cette nature chimique a été confirmée pour l'immense majorité des synapses. Cependant, la généralisation n'est pas absolue : on a découvert ultérieurement des \cle{synapses électriques}, où le courant passe directement d'une cellule à l'autre à travers des \cle{jonctions communicantes} (ou \emph{gap junctions}), canaux protéiques reliant les cytoplasmes des deux cellules. Dans ce cas, la transmission est quasi instantanée, sans libération de neuromédiateur. Le tableau suivant oppose les deux modalités, que nous détaillerons dans la section III :

\begin{center}
\begin{tabularx}{\linewidth}{|l|X|X|}
\hline
\rowcolor{popblueL} \textbf{Critère} & \textbf{Synapse chimique} & \textbf{Synapse électrique} \\
\hline
Intermédiaire & Neuromédiateur libéré dans la fente synaptique & Courant ionique direct via des jonctions communicantes \\
\hline
Fente synaptique & Large (environ 20 nm) & Très étroite (environ 2 à 3 nm) \\
\hline
Vitesse de transmission & Retard synaptique (environ 0,5 à 1 ms) & Quasi instantanée \\
\hline
Sens de transmission & Unidirectionnelle & Souvent bidirectionnelle \\
\hline
Exemple & Jonction nerf-cœur, plaque motrice & Certaines synapses du système nerveux central, muscle cardiaque \\
\hline
\end{tabularx}
\end{center}

\courssub{Exercice résolu}

\begin{exoresolu}[Analyse de l'expérience de Loewi]
\textbf{Énoncé.} Deux cœurs de grenouille, A et B, sont isolés et maintenus en survie dans du liquide de Ringer. Le cœur A conserve son nerf vague ; le cœur B est dénervé. La stimulation électrique du nerf vague du cœur A provoque le ralentissement de ses battements. Le liquide de perfusion du cœur A est alors transféré sur le cœur B : on observe, après un court délai, le ralentissement du cœur B.
\begin{enumerate}
\item Quel problème cette expérience permet-elle de résoudre ?
\item Pourquoi le cœur B a-t-il été dénervé ?
\item Interpréter le ralentissement du cœur B et conclure.
\end{enumerate}
\tcblower
\textbf{Solution détaillée.}
\begin{enumerate}
\item \textbf{Problème posé :} l'expérience vise à déterminer la \emph{nature du message} qui franchit la jonction entre le nerf vague et le muscle cardiaque : ce message est-il de nature électrique ou de nature chimique ?
\item \textbf{Rôle de la dénervation du cœur B :} le cœur B est dénervé afin d'éliminer toute possibilité de stimulation nerveuse directe. De plus, les deux cœurs n'étant reliés par aucun contact, aucun courant électrique ne peut passer de A à B. Le cœur B constitue ainsi un \emph{témoin biologique} : s'il réagit, ce ne peut être qu'à une substance transportée par le liquide transféré.
\item \textbf{Interprétation et conclusion :} le ralentissement du cœur B après transfert du seul liquide de perfusion montre que la stimulation du nerf vague du cœur A a provoqué la \emph{libération d'une substance chimique} (l'acétylcholine, ou \emph{Vagusstoff}) dans le liquide. Cette substance, véhiculée par le liquide, a agi sur le cœur B en reproduisant l'effet du nerf vague. \textbf{Conclusion :} la transmission du message nerveux au niveau de la jonction nerf-cœur est de \emph{nature chimique} : elle s'effectue par libération d'un neuromédiateur dans la fente synaptique.
\end{enumerate}
\end{exoresolu}

\begin{aretenir}[L'essentiel de la section]
\begin{itemize}
\item L'\textbf{expérience de Loewi} (1921) : deux cœurs de grenouille en survie dans du liquide de Ringer ; le cœur A est innervé, le cœur B est dénervé.
\item La stimulation du nerf vague du cœur A le ralentit ; le \textbf{transfert du liquide de perfusion} de A vers B ralentit le cœur B, sans aucun contact nerveux ni électrique.
\item Interprétation : libération d'une \textbf{substance chimique}, la \emph{Vagusstoff}, identifiée comme l'\textbf{acétylcholine}.
\item Conclusion générale : la transmission synaptique est de \textbf{nature chimique} au niveau des synapses chimiques, par libération d'un \textbf{neuromédiateur} dans la fente synaptique.
\item Il existe des neuromédiateurs \textbf{excitateurs} (noradrénaline du nerf sympathique) et \textbf{inhibiteurs} (acétylcholine sur le cœur).
\item Limite : certaines synapses fonctionnent par \textbf{transmission électrique} directe (jonctions communicantes) — c'est l'objet de la section suivante.
\end{itemize}
\end{aretenir}

Maintenant que la nature chimique de la transmission est établie, reste à comprendre \emph{comment} se déroule précisément cette transmission au niveau moléculaire : libération du neuromédiateur, fixation sur les récepteurs, destinée du médiateur — et comment fonctionnent, en parallèle, les synapses électriques. C'est l'objet de la section III.

\coursec{III. Fonctionnement des synapses chimiques et électriques}

L'expérience de Loewi, étudiée dans la section précédente, a démontré de manière éclatante que la transmission du message nerveux au niveau de la synapse cœur--nerf vague est de nature \cle{chimique} : une substance diffusible, l'acétylcholine, relaie le signal d'un élément à l'autre. Mais comment cette substance est-elle libérée ? Comment agit-elle sur l'élément postsynaptique ? Et toutes les synapses fonctionnent-elles selon ce mode chimique ? C'est à ces questions que répond la présente section, en s'appuyant sur les observations en microscopie électronique et sur les enregistrements électrophysiologiques qui ont permis de reconstituer, étape par étape, le mécanisme de la transmission synaptique.

\courssub{III.1. Les étapes de la transmission chimique}

\textbf{Situation problème.}\par Au laboratoire de physiologie de l'Université de Yaoundé I, on enregistre à l'aide de microélectrodes le potentiel de membrane d'une fibre musculaire de grenouille pendant que l'on stimule le nerf moteur qui l'innerve. On observe que chaque stimulation du nerf est suivie, après un très bref retard, d'une variation du potentiel de la fibre musculaire. Que s'est-il passé dans l'intervalle, au niveau de la \cle{plaque motrice} ?

L'observation des synapses chimiques au microscope électronique révèle trois éléments structuraux essentiels : le \cle{bouton synaptique} (renflement de l'extrémité axonique présynaptique) rempli de \cle{vésicules synaptiques} et de mitochondries, la \cle{fente synaptique} (espace de 20 à \SI{40}{nm} séparant les deux éléments) et la \cle{membrane postsynaptique} portant des \cle{récepteurs} spécifiques. Ces données ultrastructurales, croisées avec les enregistrements électriques et les dosages biochimiques, ont permis d'établir la séquence suivante, en six étapes.

\begin{experience}[Reconstitution des étapes de la transmission chimique]
\textbf{Étape 1 — Arrivée du potentiel d'action.} Le potentiel d'action (PA) se propage le long de l'axone et atteint le bouton synaptique : la membrane présynaptique se \cle{dépolarisation}.

\textbf{Étape 2 — Entrée des ions calcium.} La dépolarisation provoque l'ouverture des \cle{canaux calciques voltage-dépendants} de la membrane présynaptique. Les ions \ce{Ca^{2+}}, beaucoup plus concentrés à l'extérieur, entrent massivement dans le bouton synaptique.

\textbf{Étape 3 — Exocytose du neuromédiateur.} L'élévation de la concentration intracellulaire en \ce{Ca^{2+}} déclenche la fusion des vésicules synaptiques avec la membrane présynaptique : c'est l'\cle{exocytose}. Le \cle{neuromédiateur} (acétylcholine, noradrénaline, dopamine, GABA...) est libéré dans la fente synaptique.

\textbf{Étape 4 — Fixation sur les récepteurs.} Le neuromédiateur diffuse dans la fente et se fixe sur les \cle{récepteurs postsynaptiques spécifiques}, selon un modèle clé--serrure comparable à l'interaction enzyme--substrat.

\textbf{Étape 5 — Potentiel postsynaptique.} La fixation du médiateur provoque l'ouverture de \cle{canaux ioniques} de la membrane postsynaptique : des mouvements d'ions (\ce{Na+}, \ce{Cl-}, \ce{K+}) modifient localement le potentiel de membrane : c'est le \cle{potentiel postsynaptique}.

\textbf{Étape 6 — Élimination du neuromédiateur.} Pour que le signal reste bref et que la synapse soit prête pour un nouveau message, le neuromédiateur est rapidement éliminé de la fente, soit par \cle{dégradation enzymatique} (l'\cle{acétylcholinestérase} hydrolyse l'acétylcholine en acétate et choline), soit par \cle{recapture} (recyclage) dans l'élément présynaptique.
\end{experience}

\begin{popfigure}
\begin{center}
\begin{tikzpicture}[scale=0.92, every node/.style={font=\small}]
% Membrane présynaptique
\draw[very thick, popblue] (-5,2.2) .. controls (-2,1.7) and (2,1.7) .. (5,2.2);
\draw[very thick, popblue] (-5,1.8) .. controls (-2,1.3) and (2,1.3) .. (5,1.8);
\node[popblue, anchor=south west] at (-5,2.3) {\textbf{Membrane présynaptique (bouton synaptique)}};
% Vésicules
\foreach \x/\y in {-3.2/2.6,-1.8/2.7,-0.4/2.6,1.0/2.7,2.4/2.6}{
  \draw[fill=poporangeL, draw=poporange] (\x,\y) circle (0.22);
  \fill[poporange] (\x,\y) circle (0.05);}
% Canal calcique
\draw[fill=popgreenL, draw=popgreen, thick] (-0.35,1.35) rectangle (0.35,2.15);
\node[popgreen, anchor=west] at (0.5,2.9) {\textbf{2. Canal calcique ouvert}};
\draw[->, thick, popgreen] (0,3.6) -- (0,2.3);
\node[popgreen] at (0,3.85) {\ce{Ca^{2+}}};
% Exocytose
\draw[fill=poporangeL, draw=poporange] (3.6,1.55) circle (0.22);
\draw[->, thick, poporange] (3.6,2.35) -- (3.6,1.85);
\node[poporange, anchor=west] at (3.9,2.3) {\textbf{3. Exocytose}};
% Fente synaptique
\draw[<->, popmuted] (5.3,1.5) -- (5.3,0.1);
\node[popmuted, anchor=west] at (5.4,0.8) {fente};
% Médiateurs dans la fente
\foreach \x/\y in {-2.5/0.9,-1.2/0.6,0.2/0.9,1.5/0.6,2.8/0.9}{
  \fill[poporange] (\x,\y) circle (0.07);}
\node[poporange] at (-2.5,0.2) {\textbf{4. Diffusion du neuromédiateur}};
% Membrane postsynaptique
\draw[very thick, poppurple] (-5,-0.1) -- (5,-0.1);
\draw[very thick, poppurple] (-5,-0.5) -- (5,-0.5);
\node[poppurple, anchor=north west] at (-5,-0.6) {\textbf{Membrane postsynaptique}};
% Récepteurs
\foreach \x in {-2.5,-0.5,1.5}{
  \draw[fill=poppinkL, draw=poppink, thick] (\x,-0.55) rectangle (\x+0.5,0.0);}
\node[poppink, anchor=west] at (2.2,-0.9) {\textbf{5. Récepteurs spécifiques}};
% Enzyme
\node[popdark, star, star points=6, fill=popgoldL, draw=popgold, inner sep=1.5pt] at (4.3,0.5) {};
\node[popgold!80!black, anchor=west] at (4.6,0.5) {\textbf{6. Enzyme (AChE)}};
% PA arrivant
\draw[->, very thick, popink] (-6.2,2.0) -- (-5.1,2.0);
\node[popink, anchor=east] at (-6.3,2.0) {\textbf{1. PA}};
\end{tikzpicture}
\end{center}
\legende{Schéma séquentiel du fonctionnement d'une synapse chimique. 1 : arrivée du potentiel d'action ; 2 : ouverture des canaux calciques et entrée de \ce{Ca^{2+}} ; 3 : exocytose des vésicules ; 4 : diffusion du neuromédiateur dans la fente synaptique ; 5 : fixation sur les récepteurs postsynaptiques et ouverture de canaux ioniques ; 6 : dégradation enzymatique (acétylcholinestérase, AChE) ou recapture du médiateur.}
\end{popfigure}

\textbf{Preuve expérimentale du rôle des ions calcium.} On place une préparation nerf--muscle dans un liquide physiologique normal : la stimulation du nerf déclenche une réponse du muscle (contraction et potentiel postsynaptique enregistré). On transfère ensuite la préparation dans un milieu \textbf{dépourvu de \ce{Ca^{2+}}} : la stimulation du nerf présynaptique ne provoque \textbf{plus aucune réponse postsynaptique}, alors que le potentiel d'action se propage normalement jusqu'au bouton synaptique. Si l'on injecte du \ce{Ca^{2+}} dans le bouton, la libération de médiateur reprend.

\begin{attention}[Interpréter correctement l'expérience du calcium]
L'absence de réponse postsynaptique en milieu sans \ce{Ca^{2+}} ne signifie \textbf{pas} que le calcium est le neuromédiateur. Elle signifie que l'entrée de \ce{Ca^{2+}} dans le bouton synaptique est \cle{indispensable à l'exocytose} des vésicules, donc à la libération du médiateur. Le calcium est le \emph{coupleur} entre l'influx nerveux et la sécrétion ; le messager chimique reste le neuromédiateur.
\end{attention}

\courssub{III.2. Le potentiel postsynaptique : excitation ou inhibition}

\begin{definition}[Potentiel postsynaptique]
Le \cle{potentiel postsynaptique} (PPS) est une variation \textbf{locale} et \textbf{graduée} du potentiel de membrane de l'élément postsynaptique, provoquée par la fixation du neuromédiateur sur ses récepteurs. Contrairement au potentiel d'action, il n'obéit pas à la loi du tout ou rien : son amplitude dépend de la quantité de médiateur libérée et il ne se propage pas.
\end{definition}

Selon la nature des canaux ouverts, on distingue deux types de potentiels postsynaptiques.

\begin{itemize}
\item Le \cle{potentiel postsynaptique excitateur} (\cle{PPSE}) : le récepteur activé ouvre des canaux perméables au \ce{Na+}. L'entrée d'ions sodium provoque une \textbf{dépolarisation locale} (le potentiel passe par exemple de $-70$ à $-60$ mV) qui \textbf{rapproche} la membrane du seuil d'excitation (environ $-55$ mV) : le neurone postsynaptique est \emph{facilité}.
\item Le \cle{potentiel postsynaptique inhibiteur} (\cle{PPSI}) : le récepteur activé ouvre des canaux perméables au \ce{Cl-} (entrée de chlorures) ou au \ce{K+} (sortie de potassium). Il en résulte une \textbf{hyperpolarisation} (par exemple de $-70$ à $-75$ mV) qui \textbf{éloigne} la membrane du seuil : le neurone postsynaptique est \emph{inhibé}.
\end{itemize}

\begin{popfigure}
\begin{center}
\begin{tikzpicture}
\begin{axis}[
  width=0.95\linewidth, height=6.2cm,
  xlabel={Temps (ms)}, ylabel={Potentiel de membrane (mV)},
  xmin=0, xmax=10, ymin=-90, ymax=-30,
  xtick={0,2,4,6,8,10}, ytick={-90,-70,-55,-40},
  grid=both, grid style={popline!40},
  legend style={at={(0.98,0.35)}, anchor=east, font=\small}]
\addplot[popcurve, popgreen, very thick, domain=0:10, samples=200]
  {-70 + 15*exp(-((x-4)^2)/1.6)};
\addlegendentry{PPSE (dépolarisation)}
\addplot[popcurve, poppurple, very thick, domain=0:10, samples=200]
  {-70 - 12*exp(-((x-4)^2)/1.6)};
\addlegendentry{PPSI (hyperpolarisation)}
\addplot[popink, dashed, thick, domain=0:10] {-55};
\addlegendentry{Seuil ($-55$ mV)}
\addplot[popmuted, dotted, domain=0:10] {-70};
\addlegendentry{Repos ($-70$ mV)}
\end{axis}
\end{tikzpicture}
\end{center}
\legende{Comparaison d'un PPSE et d'un PPSI enregistrés sur le même neurone postsynaptique. Le PPSE rapproche le potentiel du seuil (effet excitateur) ; le PPSI l'en éloigne (effet inhibiteur). Un PPSE isolé n'atteint généralement pas le seuil : c'est la sommation de plusieurs PPSE qui peut déclencher un potentiel d'action (voir section IV).}
\end{popfigure}

\begin{attention}[Le neuromédiateur n'est ni excitateur ni inhibiteur « en soi »]
Erreur fréquente au Baccalauréat : écrire « l'acétylcholine est un médiateur excitateur ». C'est \textbf{faux} dans l'absolu. L'effet dépend du \cle{type de récepteur postsynaptique} et du canal ionique qu'il commande. Ainsi, l'acétylcholine \textbf{excite} le muscle squelettique (PPSE à la plaque motrice, d'où la contraction) mais \textbf{inhibe} le cœur (ralentissement du rythme cardiaque observé par Loewi). Même médiateur, récepteurs différents, effets opposés.
\end{attention}

\courssub{III.3. Caractéristiques de la transmission chimique}

\begin{propriete}[Unidirectionnalité de la transmission chimique]
La transmission chimique est \cle{unidirectionnelle} : elle s'effectue toujours de l'élément présynaptique vers l'élément postsynaptique, car \textbf{seul} l'élément présynaptique possède les vésicules de neuromédiateur et \textbf{seule} la membrane postsynaptique porte les récepteurs spécifiques. Cette asymétrie structurale impose le sens unique de circulation de l'information nerveuse.
\end{propriete}

Deux autres propriétés distinguent la synapse chimique :

\begin{itemize}
\item Le \cle{délai synaptique} : entre l'arrivée du PA dans le bouton et l'apparition du PPS, il s'écoule environ \SI{0,5}{ms} à \SI{1}{ms}, temps nécessaire à l'entrée du calcium, à l'exocytose, à la diffusion du médiateur et à l'ouverture des canaux. Ce retard, mesurable à l'oscilloscope, est une signature de la transmission chimique.
\item La \cle{fatigabilité} : une stimulation répétée et prolongée épuise le stock de vésicules et de médiateur ; la réponse postsynaptique s'affaiblit puis cesse. La synapse chimique est donc le maillon fragile — mais aussi modulable — de la chaîne nerveuse.
\end{itemize}

\begin{exemplebox}[Un exemple chiffré à retenir]
Le délai synaptique d'une synapse chimique est d'environ \num{0,5} à \SI{1}{ms}, contre pratiquement \SI{0}{ms} pour une synapse électrique. Sur un circuit réflexe comportant plusieurs synapses chimiques, ces délais s'additionnent et expliquent le temps de réaction mesuré (quelques dizaines de millisecondes pour un réflexe simple).
\end{exemplebox}

\begin{savaistu}[Une enzyme éclair : l'acétylcholinestérase]
L'acétylcholinestérase, présente dans la fente synaptique de la plaque motrice, détruit l'acétylcholine en quelques millisecondes. Grâce à cette élimination quasi immédiate, chaque message nerveux reste bref et le muscle peut se relâcher entre deux influx : c'est ce qui permet des mouvements précis et rapides, comme ceux des doigts d'un joueur de \emph{mvet} ou d'un tailleur de Yaoundé. Certains insecticides (organophosphorés) bloquent cette enzyme : l'acétylcholine s'accumule alors et provoque des contractions continues, spasmes et paralysies — un exemple frappant des effets des substances sur la synapse, que nous détaillerons en section V.
\end{savaistu}

\courssub{III.4. La synapse électrique : les jonctions communicantes}

Toutes les synapses ne passent pas par un intermédiaire chimique. Sur certaines électronographies, on observe des synapses où les membranes des deux éléments sont \textbf{très rapprochées} (espace de 2 à \SI{4}{nm} seulement, contre 20 à \SI{40}{nm} pour la fente d'une synapse chimique), \textbf{sans vésicules synaptiques}. Ce sont des \cle{synapses électriques}.

Leur ultrastructure montre des \cle{jonctions communicantes} (ou \emph{gap junctions}) : des canaux protéiques transmembranaires, appelés \cle{connexons}, qui s'affrontent et relient directement les cytoplasmes des deux cellules. Les ions circulent librement d'une cellule à l'autre à travers ces tunnels.

\begin{propriete}[Propriétés de la transmission électrique]
Dans une synapse électrique, le courant ionique passe \textbf{directement} d'une cellule à l'autre par les connexons. Il en résulte :
\begin{itemize}
\item une transmission \textbf{quasi instantanée} (pas de délai synaptique mesurable) ;
\item une transmission le plus souvent \textbf{bidirectionnelle} (le courant peut passer dans les deux sens) ;
\item une transmission \textbf{infatigable} (pas de stock de médiateur à épuiser).
\end{itemize}
On trouve ces synapses là où la rapidité et la synchronisation sont vitales : cellules du \textbf{muscle cardiaque} (contraction simultanée des cardiomyocytes), \textbf{muscle lisse} des organes creux, et certains circuits nerveux d'\textbf{alerte rapide} (réflexes de fuite des invertébrés).
\end{propriete}

\begin{popfigure}
\begin{center}
\begin{tikzpicture}[scale=0.9, every node/.style={font=\small}]
% ===== Synapse électrique (gauche) =====
\begin{scope}[shift={(-4.2,0)}]
\draw[very thick, popgreen] (-2,1.2) -- (2,1.2);
\draw[very thick, popgreen] (-2,0.8) -- (2,0.8);
\draw[very thick, poppurple] (-2,-0.2) -- (2,-0.2);
\draw[very thick, poppurple] (-2,-0.6) -- (2,-0.6);
% Connexons
\foreach \x in {-1.2,0,1.2}{
  \draw[fill=popgoldL, draw=popgold, thick] (\x-0.18,0.75) rectangle (\x+0.18,-0.15);}
\draw[->, thick, popink] (-1.2,1.6) -- (-1.2,-1.0);
\draw[->, thick, popink] (0,-1.0) -- (0,1.6);
\draw[->, thick, popink] (1.2,1.6) -- (1.2,-1.0);
\node[popdark, font=\small\bfseries] at (0,2.1) {Synapse électrique};
\node[popmuted] at (0,-1.35) {jonctions communicantes (connexons)};
\draw[<->, popmuted] (2.2,0.8) -- (2.2,-0.2);
\node[popmuted, anchor=west] at (2.3,0.3) {2 à \SI{4}{nm}};
\end{scope}
% ===== Synapse chimique (droite) =====
\begin{scope}[shift={(4.2,0)}]
\draw[very thick, popblue] (-2,1.2) -- (2,1.2);
\draw[very thick, popblue] (-2,0.8) -- (2,0.8);
\draw[very thick, poppurple] (-2,-0.6) -- (2,-0.6);
\draw[very thick, poppurple] (-2,-1.0) -- (2,-1.0);
\foreach \x in {-1.3,-0.4,0.6,1.4}{
  \draw[fill=poporangeL, draw=poporange] (\x,1.55) circle (0.18);}
\foreach \x/\y in {-1.0/0.3,0.1/0.0,1.1/0.35}{
  \fill[poporange] (\x,\y) circle (0.06);}
\foreach \x in {-1.0,0.4}{
  \draw[fill=poppinkL, draw=poppink, thick] (\x,-1.05) rectangle (\x+0.4,-0.55);}
\draw[->, thick, popink] (-1.8,1.6) -- (-1.8,-1.4);
\node[popdark, font=\small\bfseries] at (0,2.1) {Synapse chimique};
\node[popmuted] at (0,-1.75) {vésicules, fente large, récepteurs};
\draw[<->, popmuted] (2.2,0.8) -- (2.2,-0.6);
\node[popmuted, anchor=west] at (2.3,0.1) {20 à \SI{40}{nm}};
\end{scope}
\end{tikzpicture}
\end{center}
\legende{Comparaison ultrastructurale. À gauche, synapse électrique : membranes accolées (2 à \SI{4}{nm}) reliées par des connexons, passage direct et bidirectionnel du courant ionique, aucune vésicule. À droite, synapse chimique : vésicules présynaptiques, fente large (20 à \SI{40}{nm}), récepteurs postsynaptiques, transmission unidirectionnelle.}
\end{popfigure}

\begin{methode}[Identifier le type de synapse sur une électronographie]
\begin{enumerate}
\item Chercher des \textbf{vésicules} dans l'élément présynaptique.
\item Mesurer (ou estimer à l'échelle) la \textbf{largeur de l'espace} intercellulaire.
\item Chercher des \textbf{récepteurs} (épaississements) sur la membrane postsynaptique.
\end{enumerate}
\textbf{Conclusion :} présence de vésicules + fente large (20 à \SI{40}{nm}) + récepteurs $\Rightarrow$ synapse \textbf{chimique}. Membranes accolées (2 à \SI{4}{nm}) + absence de vésicules + jonctions communicantes $\Rightarrow$ synapse \textbf{électrique}.
\end{methode}

\courssub{III.5. Comparaison synthétique des deux types de transmission}

\begin{center}
\rowcolors{2}{popblueL}{white}
\begin{tabularx}{\linewidth}{|l|X|X|}
\hline
\rowcolor{popblue!40} \textbf{Critère} & \textbf{Synapse chimique} & \textbf{Synapse électrique} \\
\hline
Support de la transmission & Neuromédiateur (molécule chimique) & Courant ionique direct (ions) \\
\hline
Structure clé & Vésicules, fente de 20 à \SI{40}{nm}, récepteurs & Connexons (gap junctions), espace de 2 à \SI{4}{nm} \\
\hline
Vitesse / délai & Délai synaptique de \num{0,5} à \SI{1}{ms} & Quasi instantanée ($\approx$ \SI{0}{ms}) \\
\hline
Sens de transmission & Unidirectionnelle & Souvent bidirectionnelle \\
\hline
Fatigabilité & Fatigable (épuisement des vésicules) & Infatigable \\
\hline
Exemples & Plaque motrice, synapses du SNC, synapse nerf vague--cœur & Muscle cardiaque, muscle lisse, circuits d'alerte \\
\hline
\end{tabularx}
\end{center}

\begin{exoresolu}[QCM de synthèse]
\textbf{Énoncé.} Pour chaque affirmation, indiquer si elle est vraie ou fausse, en justifiant brièvement.

\textbf{a)} Dans une synapse chimique, le neuromédiateur traverse la membrane postsynaptique pour agir à l'intérieur de la cellule.

\textbf{b)} En l'absence de \ce{Ca^{2+}} extracellulaire, la stimulation du neurone présynaptique ne provoque plus de potentiel postsynaptique.

\textbf{c)} L'acétylcholine est toujours un neuromédiateur excitateur.

\textbf{d)} La transmission dans une synapse électrique est bidirectionnelle et sans délai notable.
\tcblower
\textbf{Solution détaillée.}

\textbf{a) Faux.} Le neuromédiateur ne pénètre \textbf{pas} dans la cellule postsynaptique : il se fixe sur des récepteurs membranaires externes, ce qui provoque l'ouverture de canaux ioniques. Ce sont des \emph{ions} (\ce{Na+}, \ce{Cl-}...) qui traversent la membrane, pas le médiateur.

\textbf{b) Vrai.} L'entrée de \ce{Ca^{2+}} dans le bouton synaptique est indispensable à l'exocytose des vésicules. Sans calcium, pas de libération de médiateur, donc pas de potentiel postsynaptique — alors même que le PA atteint normalement le bouton.

\textbf{c) Faux.} L'effet d'un médiateur dépend du récepteur postsynaptique : l'acétylcholine excite le muscle squelettique (plaque motrice) mais inhibe le cœur (expérience de Loewi). Un médiateur n'est donc jamais « excitateur » ou « inhibiteur » en soi.

\textbf{d) Vrai.} Les connexons relient directement les cytoplasmes : le courant ionique passe dans les deux sens, sans étape chimique, donc sans délai synaptique mesurable.
\end{exoresolu}

\begin{aretenir}[L'essentiel de la section III]
\begin{itemize}
\item La transmission chimique se déroule en six étapes : arrivée du PA $\rightarrow$ entrée de \ce{Ca^{2+}} $\rightarrow$ exocytose $\rightarrow$ fixation sur les récepteurs $\rightarrow$ potentiel postsynaptique $\rightarrow$ élimination du médiateur (enzyme ou recapture).
\item Le \ce{Ca^{2+}} est le coupleur indispensable entre l'influx et la libération du médiateur.
\item Le PPS est local et gradué : \textbf{PPSE} = dépolarisation (entrée de \ce{Na+}) rapprochant du seuil ; \textbf{PPSI} = hyperpolarisation (entrée de \ce{Cl-} ou sortie de \ce{K+}) éloignant du seuil.
\item Transmission chimique : unidirectionnelle, délai de \num{0,5} à \SI{1}{ms}, fatigable. Transmission électrique : directe par connexons, quasi instantanée, souvent bidirectionnelle, infatigable.
\end{itemize}
\end{aretenir}

Maintenant que les mécanismes des deux types de transmission sont établis, une question demeure : un neurone postsynaptique reçoit souvent des centaines de messages, excitateurs et inhibiteurs, en provenance de nombreuses synapses. Comment « décide »-t-il d'émettre ou non un potentiel d'action ? C'est le problème des \emph{sommations} et du \emph{rôle intégrateur} du neurone, objet de la section IV.

\coursec{IV. Sommations et rôle intégrateur du neurone}

Dans la section précédente, nous avons vu comment un influx nerveux, arrivé au bouton synaptique, déclenche la libération d'un neurotransmetteur qui se lie à ses récepteurs sur la membrane postsynaptique. Cette liaison ouvre des canaux ioniques et génère un \textbf{potentiel postsynaptique} (PPS) : un \textbf{PPSE} (excitateur, dépolarisant) si des cations (\ce{Na+}, \ce{Ca^{2+}}) entrent, ou un \textbf{PPSI} (inhibiteur, hyperpolarisant ou stabilisant) si des anions (\ce{Cl-}) entrent ou des cations (\ce{K+}) sortent.

Mais une question fondamentale se pose : \emph{un unique potentiel postsynaptique suffit-il à déclencher un potentiel d'action dans le neurone postsynaptique ?} C'est à cette question que nous allons répondre en analysant des enregistrements intracellulaires classiques, ce qui nous mènera aux concepts de \textbf{sommation} et de \textbf{rôle intégrateur du neurone}.

\courssub{Un PPSE unique est généralement sous-seuil : le constat expérimental}

\begin{experience}[Enregistrement intracellulaire d'un PPSE unique]
\textbf{Matériel :} Préparation isolée de moelle épinière (grenouille ou tortue), microélectrode intracellulaire remplie de \ce{KCl} \SI{3}{M} enfoncée dans un motoneurone $\alpha$, électrode de stimulation extracellulaire sur une fibre afferente sensitive (racine dorsale), oscilloscope.

\textbf{Protocole :}
\begin{enumerate}
    \item On repère le potentiel de repos du motoneurone (environ \SI{-70}{mV}).
    \item On délivre une \textbf{stimulation unique} (impulsion brève, \SI{0,1}{ms}, intensité seuil) sur la fibre afferente unique.
    \item On enregistre la réponse électrique du motoneurone.
\end{enumerate}

\textbf{Observations (voir Figure 1) :}
\begin{itemize}
    \item Après un court délai synaptique (\SI{0,5}{à 1}{ms}), la tension de membrane s'écarte du potentiel de repos vers des valeurs moins négatives : c'est une \textbf{dépolarisation}.
    \item Cette dépolarisation est \textbf{graduée} : son amplitude (ici \SI{+10}{mV}) dépend de l'intensité de la stimulation (nombre de fibres stimulées, quantité de neurotransmetteur libéré).
    \item Elle est \textbf{transitoire} : elle monte rapidement, atteint un pic, puis redescend exponentiellement vers le potentiel de repos (constante de temps membranaire $\tau_m \approx \SI{10}{à 20}{ms}$).
    \item \textbf{Crucial :} Le sommet de la dépolarisation (\SI{-60}{mV} ici) \textbf{n'atteint pas le seuil de déclenchement} du potentiel d'action (environ \SI{-55}{mV}). Aucun potentiel d'action n'est émis.
\end{itemize}

\textbf{Interprétation :}
Un unique potentiel postsynaptique excitateur (PPSE) est un \textbf{signal local, gradué, décrémentaire} (il s'atténue en se propageant passivement le long des dendrites et du soma). Il ne se propage pas activement comme un potentiel d'action. Pour qu'un potentiel d'action naisse, la dépolarisation doit atteindre le \textbf{seuil} au niveau du \textbf{cône d'émergence} (segment initial de l'axone), zone à forte densité de canaux \ce{Na+} voltage-dépendants.

\begin{popfigure}
\begin{tikzpicture}[scale=0.9, every node/.style={font=\small}]
    % Axes
    \draw[->] (0,0) -- (10,0) node[right] {$t$ (ms)};
    \draw[->] (0,-1) -- (0,4) node[above] {$V_m$ (mV)};
    % Graduations temps
    \foreach \x/\label in {0/0, 2/2, 4/4, 6/6, 8/8, 10/10} \draw (\x, -0.1) -- (\x, 0.1) node[below] {\label};
    % Graduations tension
    \foreach \y/\label in {0/-70, 1/-60, 2/-50, 3/-40} \draw (-0.1, \y) -- (0.1, \y) node[left] {\label};
    % Ligne de repos
    \draw[dashed, popmuted] (0,0) -- (10,0) node[right, pos=0.9] {$V_{repos}$};
    % Ligne de seuil
    \draw[dashed, poporange] (0,1.5) -- (10,1.5) node[right, poporange] {Seuil (\SI{-55}{mV})};
    % PPSE unique
    \draw[popblue, very thick, domain=0.5:10, samples=200] plot (\x, {1.5*exp(-(\x-0.5)/2.5) * (1-exp(-(\x-0.5)/0.5))});
    % Annotation
    \node[popblue, align=center] at (3, 2.8) {PPSE unique\\Amplitude \SI{10}{mV}};
    \draw[popblue, ->] (3.2, 2.5) -- (2.5, 1.6);
    \node[poporange, align=center] at (7, 3.2) {Seuil non atteint :\\Pas de potentiel d'action};
\end{tikzpicture}
\legende{Figure 1 — Enregistrement intracellulaire d'un PPSE unique (sous-seuil). Le potentiel de repos est à \SI{-70}{mV}, le seuil à \SI{-55}{mV}. La dépolarisation maximale atteint \SI{-60}{mV}, insuffisante pour déclencher un potentiel d'action.}
\end{popfigure}

\end{experience}

\begin{attention}[Erreur fréquente : Confondre PPSE et Potentiel d'Action]
Ne confonds jamais un \textbf{PPSE} (potentiel postsynaptique excitateur) avec un \textbf{potentiel d'action} (influx nerveux).
\begin{itemize}
    \item \textbf{PPSE :} Signal \textbf{gradué} (amplitude variable), \textbf{local} (ne se propage pas activement), \textbf{décrémentaire} (s'atténue avec la distance), \textbf{sommatif} (s'additionne aux autres PPS). Pas de période de réfractaire.
    \item \textbf{Potentiel d'action :} Signal \textbf{tout-ou-rien} (amplitude constante, loi du tout ou rien), \textbf{propagé} sans décrément le long de l'axone, \textbf{non sommatif} (période de réfractaire absolue), \textbf{seuil} obligatoire.
\end{itemize}
\emph{Astuce Bac :} Si le sujet parle d'amplitude variable, de sommation, ou de propagation décrémentaire $\rightarrow$ c'est un PPS. Si amplitude constante, propagation active, tout-ou-rien $\rightarrow$ c'est un potentiel d'action.
\end{attention}

\courssub{La sommation temporelle : l'addition dans le temps}

Puisqu'un PPSE unique est insuffisant, comment le système nerveux fait-il pour transmettre l'information ? La réponse réside dans la \textbf{fréquence} des impulsions arrivant sur une \textbf{même} synapse.

\begin{experience}[Mise en évidence de la sommation temporelle (Eccles et al., années 1950)]
\textbf{Protocole :} Même préparation. On stimule \textbf{une seule fibre afferente} avec un train d'impulsions à \textbf{haute fréquence} (ex. \SI{100}{Hz}, soit un intervalle de \SI{10}{ms}).

\textbf{Observations (voir Figure 2) :}
\begin{itemize}
    \item Le premier PPSE est identique à celui vu précédemment (amplitude \SI{10}{mV}).
    \item Le second PPSE arrive \textbf{avant que le premier ne soit totalement revenu au repos}. Il démarre donc d'un potentiel de membrane déjà dépolarisé (ex. \SI{-65}{mV}).
    \item Les deux dépolarisations \textbf{s'additionnent} : le sommet du second PPSE atteint \SI{-55}{mV} (seuil).
    \item Un \textbf{potentiel d'action} est déclenché au cône d'émergence.
    \item Si la stimulation se poursuit, d'autres potentiels d'action suivent (fréquence codée par la fréquence des PPSE).
\end{itemize}

\begin{popfigure}
\begin{tikzpicture}[scale=0.9, every node/.style={font=\small}]
    \draw[->] (0,0) -- (12,0) node[right] {$t$ (ms)};
    \draw[->] (0,-1) -- (0,4.5) node[above] {$V_m$ (mV)};
    \foreach \x/\label in {0/0, 2/2, 4/4, 6/6, 8/8, 10/10, 12/12} \draw (\x, -0.1) -- (\x, 0.1) node[below] {\label};
    \foreach \y/\label in {0/-70, 1/-60, 1.5/-55, 2/-50, 3/-40, 4/-30} \draw (-0.1, \y) -- (0.1, \y) node[left] {\label};
    \draw[dashed, popmuted] (0,0) -- (12,0) node[right, pos=0.95] {$V_{repos}$};
    \draw[dashed, poporange] (0,1.5) -- (12,1.5) node[right, poporange] {Seuil};
    % PPSE 1
    \draw[popblue, very thick, domain=1:12, samples=200] plot (\x, {1.5*exp(-(\x-1)/2.5) * (1-exp(min(3,-(\x-1)/0.5)))});
    % PPSE 2 (starts at t=3)
    \draw[popblue, very thick, domain=3:12, samples=200] plot (\x, {1.5*exp(-(\x-3)/2.5) * (1-exp(min(3,-(\x-3)/0.5)))});
    % PPSE 3 (starts at t=5)
    \draw[popblue, very thick, domain=5:12, samples=200] plot (\x, {1.5*exp(-(\x-5)/2.5) * (1-exp(min(3,-(\x-5)/0.5)))});
    % Summation trace (conceptual thick red line showing envelope)
    \draw[poppink, ultra thick, domain=1:6, samples=100] plot (\x, {
        ( (\x>=1) ? 1.5*exp(-(\x-1)/2.5) * (1-exp(min(3,-(\x-1)/0.5))) : 0 ) +
        ( (\x>=3) ? 1.5*exp(-(\x-3)/2.5) * (1-exp(min(3,-(\x-3)/0.5))) : 0 ) +
        ( (\x>=5) ? 1.5*exp(-(\x-5)/2.5) * (1-exp(min(3,-(\x-5)/0.5))) : 0 )
    });
    % Action Potential schematic at t~5.5
    \draw[popdark, very thick] (5.5, 1.5) -- (5.6, 4) -- (5.8, -0.5) -- (6.0, 0);
    \node[popdark, align=center] at (6.5, 3.5) {Potentiel\\d'action};
    \node[popblue, align=center] at (2, 3.8) {Stimulations\\à \SI{100}{Hz}};
    \draw[popblue, ->] (2.2, 3.4) -- (2, 1.8);
    \draw[popblue, ->] (2.2, 3.4) -- (4, 1.8);
    \draw[popblue, ->] (2.2, 3.4) -- (6, 1.8);
\end{tikzpicture}
\legende{Figure 2 — Sommation temporelle. Trois stimulations successives d'une même afférence à \SI{100}{Hz} (intervalles \SI{10}{ms}). Les PPSE s'additionnent (courbe rose) et atteignent le seuil, déclenchant un potentiel d'action (pic noir).}
\end{popfigure}

\end{experience}

\begin{definition}[Sommation temporelle]
C'est l'\textbf{addition de potentiels postsynaptiques (PPSE ou PPSI) successifs} générés sur \textbf{une même synapse} (ou un même groupe de synapses activées par un même neurone présynaptique) lorsque les stimulations se succèdent à un intervalle de temps \textbf{inférieur à la durée de retour au repos} du potentiel précédent (constante de temps membranaire $\tau_m$).
\end{definition}

\begin{propriete}[Conditions de la sommation temporelle]
\begin{itemize}
    \item \textbf{Une seule voie afférente} stimulée de façon répétée.
    \item \textbf{Haute fréquence} de stimulation (généralement $> \SI{50}{Hz}$, intervalle $< \tau_m$).
    \item Les PPSE doivent avoir la \textbf{même polarité} (excitateurs s'additionnent entre eux, inhibiteurs aussi).
\end{itemize}
\emph{Mécanisme biophysique :} La membrane du neurone agit comme un \textbf{condenseur} (capacité $C_m$). Le courant synaptique ($I_{syn}$) charge ce condensateur. Si un second courant arrive avant que le condensateur ne se soit déchargé (via les fuites $K^+$, $Cl^-$), la tension ($V = Q/C$) monte davantage.
\end{propriete}

\courssub{La sommation spatiale : l'addition dans l'espace}

Un neurone typique (ex. motoneurone spinal) reçoit des milliers de synapses sur ses dendrites et son soma (jusqu'à \SI{10000}{à 100000} pour un motoneurone $\alpha$). Que se passe-t-il si plusieurs afférences \textbf{différentes} sont actives \textbf{simultanément} ?

\begin{experience}[Mise en évidence de la sommation spatiale]
\textbf{Protocole :} On place deux électrodes de stimulation sur \textbf{deux fibres afferentes distinctes} (A et B) qui font synapse sur \textbf{le même motoneurone}, mais à des endroits différents (ex. dendrites proximales vs distales).
\begin{enumerate}
    \item Stimulation seule de A $\rightarrow$ PPSE de \SI{5}{mV} (sous-seuil).
    \item Stimulation seule de B $\rightarrow$ PPSE de \SI{5}{mV} (sous-seuil).
    \item Stimulation \textbf{simultanée} de A et B $\rightarrow$ PPSE composite de \SI{10}{mV} $\rightarrow$ \textbf{Seuil atteint} $\rightarrow$ Potentiel d'action.
\end{enumerate}

\begin{popfigure}
\begin{tikzpicture}[scale=0.9, every node/.style={font=\small}]
    \draw[->] (0,0) -- (10,0) node[right] {$t$ (ms)};
    \draw[->] (0,-1) -- (0,4) node[above] {$V_m$ (mV)};
    \foreach \x/\label in {0/0, 2/2, 4/4, 6/6, 8/8, 10/10} \draw (\x, -0.1) -- (\x, 0.1) node[below] {\label};
    \foreach \y/\label in {0/-70, 0.5/-65, 1/-60, 1.5/-55, 2/-50} \draw (-0.1, \y) -- (0.1, \y) node[left] {\label};
    \draw[dashed, popmuted] (0,0) -- (10,0) node[right, pos=0.95] {$V_{repos}$};
    \draw[dashed, poporange] (0,1.5) -- (10,1.5) node[right, poporange] {Seuil};
    % Stim A alone
    \draw[popblue, thick, domain=2:10, samples=100] plot (\x, {0.75*exp(-(\x-2)/3) * (1-exp(min(3,-(\x-2)/0.5)))});
    \node[popblue] at (3, 1.2) {A seul};
    % Stim B alone (shifted right for clarity on graph, but conceptually same time)
    \draw[popgreen, thick, domain=5:10, samples=100, dashed] plot (\x, {0.75*exp(-(\x-5)/3) * (1-exp(min(3,-(\x-5)/0.5)))});
    \node[popgreen] at (7, 1.2) {B seul};
    % A + B Simultaneous (at t=2)
    \draw[poppink, ultra thick, domain=2:10, samples=100] plot (\x, {1.5*exp(-(\x-2)/3) * (1-exp(min(3,-(\x-2)/0.5)))});
    \node[poppink, align=center] at (5, 3.2) {A + B simultanés\\(Sommation spatiale)};
    \draw[poppink, ->] (5.2, 2.9) -- (4.5, 2.2);
    % AP
    \draw[popdark, very thick] (3.2, 1.5) -- (3.3, 3.5) -- (3.5, -0.5) -- (3.7, 0);
\end{tikzpicture}
\legende{Figure 3 — Sommation spatiale. Deux afférences A et B inefficaces seules (PPSE \SI{5}{mV} chacune, pointillés bleus/verts) deviennent efficaces si stimulées simultanément (courbe rose, PPSE \SI{10}{mV}) : le seuil est franchi.}
\end{popfigure}

\end{experience}

\begin{definition}[Sommation spatiale]
C'est l'\textbf{addition de potentiels postsynaptiques} (PPSE ou PPSI) générés \textbf{simultanément} sur \textbf{plusieurs synapses distinctes} d'un même neurone postsynaptique, par l'activation de \textbf{plusieurs neurones présynaptiques différents}.
\end{definition}

\begin{propriete}[Atténuation et localisation des synapses]
Les PPSE naissent sur les dendrites ou le soma. Ils se propagent passivement (courant local) vers le cône d'émergence. Leur amplitude \textbf{décroît avec la distance} (propriétés câblées du neurone : résistance axiale $r_i$, résistance membranaire $r_m$, capacité $c_m$).
\begin{itemize}
    \item Synapses \textbf{proximales} (soma, dendrites épaisses) $\rightarrow$ PPSE peu atténué au cône d'émergence $\rightarrow$ fort poids intégratif.
    \item Synapses \textbf{distales} (fines dendrites) $\rightarrow$ PPSE fortement atténué $\rightarrow$ faible poids intégratif.
\end{itemize}
C'est pourquoi la sommation spatiale est plus efficace si les synapses activées sont proches du cône d'émergence ou nombreuses.
\end{propriete}

\courssub{La sommation spatio-temporelle et l'intégration excitatrice/inhibitrice}

Dans la réalité physiologique, les deux modes opèrent conjointement : un neurone reçoit une \textbf{pluie de PPSE et de PPSI} sur ses milliers de synapses, avec des timings variés. C'est la \textbf{sommation spatio-temporelle} (ou \textbf{intégration synaptique}).

\begin{experience}[Intégration des signaux excitateurs et inhibiteurs (Coombs, Eccles, Fatt, 1955)]
\textbf{Protocole :} On stimule une voie excitatrice (E) et une voie inhibitrice (I) convergent sur le même motoneurone.
\begin{enumerate}
    \item E seul $\rightarrow$ Potentiel d'action.
    \item I seul $\rightarrow$ PPSI (hyperpolarisation, ex. \SI{-5}{mV}).
    \item E + I simultanés $\rightarrow$ Le PPSI \textbf{soustrait} sa tension au PPSE. Si $PPSE + PPSI < \text{Seuil}$, \textbf{pas de potentiel d'action}.
    \item Décalage temporel : Si I arrive \textbf{avant} E, il hyperpolarise la membrane, augmentant la "distance" au seuil $\rightarrow$ inhibition plus efficace. Si I arrive \textbf{après} le pic du PPSE, son effet est réduit.
\end{enumerate}

\begin{popfigure}
\begin{tikzpicture}[scale=1, every node/.style={font=\small}]
    % Neuron Schema
    % Dendrites
    \draw[popblue, thick] (-3, 2) .. controls (-1, 2.5) and (-1, 1.5) .. (0, 1.5);
    \draw[popblue, thick] (-3, -2) .. controls (-1, -1.5) and (-1, -2.5) .. (0, -1.5);
    \draw[popblue, thick] (-2.5, 2.5) .. controls (-0.5, 3) and (-0.5, 2) .. (0, 1.5);
    \draw[popblue, thick] (-2.5, -2.5) .. controls (-0.5, -3) and (-0.5, -2) .. (0, -1.5);
    % Soma
    \draw[popdark, thick, fill=popblueL!30] (0,0) ellipse (1.5 and 1.5);
    % Axon Hillock (Cone d'émergence)
    \draw[poporange, thick, fill=poporangeL!50] (1.5,0) ellipse (0.4 and 0.6);
    % Axon
    \draw[popdark, thick] (1.9, 0.1) -- (5, 0.1);
    \draw[popdark, thick] (1.9, -0.1) -- (5, -0.1);
    % Myelin sheaths
    \foreach \x in {2.2, 2.8, 3.4, 4.0, 4.6} {
        \draw[popgray, thick] (\x, 0.2) -- (\x+0.3, 0.2);
        \draw[popgray, thick] (\x, -0.2) -- (\x+0.3, -0.2);
        \draw[popgray, thick] (\x, 0.2) -- (\x, -0.2);
        \draw[popgray, thick] (\x+0.3, 0.2) -- (\x+0.3, -0.2);
    }
    % Synapses Excitatory (Green +)
    \node[popgreen, circle, draw, thick, inner sep=2pt] (E1) at (-2.5, 2.2) {$+$};
    \node[popgreen, circle, draw, thick, inner sep=2pt] (E2) at (-1.5, -1.8) {$+$};
    \node[popgreen, circle, draw, thick, inner sep=2pt] (E3) at (0.5, 1.2) {$+$};
    % Synapses Inhibitory (Red -)
    \node[poppink, circle, draw, thick, inner sep=2pt] (I1) at (-2, 0.5) {$-$};
    \node[poppink, circle, draw, thick, inner sep=2pt] (I2) at (0.8, -0.8) {$-$};
    % Presynaptic fibers
    \draw[popmuted, ->, thick] (-3.5, 2.2) -- (E1);
    \draw[popmuted, ->, thick] (-3.5, -1.8) -- (E2);
    \draw[popmuted, ->, thick] (-3.5, 0.5) -- (I1);
    % Labels
    \node[align=center] at (-3.5, 3) {Afférences\\excitatrices};
    \node[align=center] at (-3.5, -2.8) {Afférences\\excitatrices};
    \node[align=center] at (-3.5, 0.5) {Afférences\\inhibitrices};
    \node[poporange, align=center, font=\bfseries\small] at (1.5, -1.2) {Cône d'émergence\\(Zone de décision)};
    \draw[poporange, ->] (1.5, -1.0) -- (1.5, -0.6);
    \node[align=center] at (5.5, 0) {Axone\\(Propagation\\tout-ou-rien)};
    % Summation formula
    \node[draw, rounded corners, fill=popcream, align=left, font=\small] at (-4.5, -4) {
        \textbf{Intégration au cône d'émergence :}\\
        $V_{totale} = V_{repos} + \sum PPSE_i + \sum PPSI_j$\\
        (avec $PPSI_j < 0$)\\
        Si $V_{totale} \ge V_{seuil} \rightarrow$ Potentiel d'action\\
        Si $V_{totale} < V_{seuil} \rightarrow$ Pas d'influx
    };
\end{tikzpicture}
\legende{Figure 4 — Schéma d'intégration synaptique. Un neurone postsynaptique (motoneurone) reçoit des PPSE (boutons verts $+$) et des PPSI (boutons roses $-$). La somme algébrique de ces potentiels, propagée passivement jusqu'au cône d'émergence (orange), décide du déclenchement du potentiel d'action.}
\end{popfigure}

\end{experience}

\begin{propriete}[Rôle intégrateur du neurone : Centre de décision]
Le neurone est l'\textbf{unité fondamentale de traitement de l'information nerveuse}. Il réalise une \textbf{opération algébrique} (somme pondérée) sur ses entrées :
\[
V_{cône} = V_{repos} + \sum_{i=1}^{n} w_i \cdot PPSE_i + \sum_{j=1}^{m} w'_j \cdot PPSI_j
\]
où $w_i, w'_j$ sont des \textbf{coefficients de pondération} (efficacité synaptique, dépendant de la distance au cône, du nombre de récepteurs, de la quantité de neurotransmetteur) et où les $PPSI_j$ sont des valeurs \textbf{négatives} (hyperpolarisantes).
\begin{itemize}
    \item \textbf{Seuil atteint ($V_{cône} \ge V_{seuil}$) :} Ouverture massive des canaux \ce{Na+} voltage-dépendants $\rightarrow$ \textbf{Potentiel d'action} (Loi du \textbf{tout ou rien} : amplitude et forme constantes, indépendantes de l'intensité de la stimulation supraliminaire).
    \item \textbf{Seuil non atteint :} Retour passif au repos. L'information est \textbf{filtrée}, non transmise.
\end{itemize}
\end{propriete}

\begin{aretenir}[Synthèse : Les trois sommations]
\begin{center}
\begin{tabularx}{\linewidth}{|l|X|X|X|}\hline
\textbf{Type} & \textbf{Stimulation} & \textbf{Localisation synapses} & \textbf{Condition clé} \\ \hline
\textbf{Temporelle} & Une afférence, \textbf{répétée} (haute fréquence) & Une zone synaptique & Intervalle $< \tau_m$ (constante de temps) \\ \hline
\textbf{Spatiale} & \textbf{Plusieurs} afférences, \textbf{simultanées} & Zones synaptiques \textbf{distinctes} & Convergence sur même neurone \\ \hline
\textbf{Spatio-temporelle} & Plusieurs afférences, \textbf{répétées} & Zones distinctes & Combinaison des deux précédentes \\ \hline
\end{tabularx}
\end{center}
\textbf{Résultat :} Seule la tension au \textbf{cône d'émergence} compte. $PPSE$ et $PPSI$ s'additionnent \textbf{algébriquement}.
\end{aretenir}

\courssub{Exemple chiffré et exercice résolu}

\begin{exemplebox}[Calcul de seuils et sommation]
Un motoneurone spinal a un potentiel de repos $V_{repos} = \SI{-70}{mV}$ et un seuil d'excitabilité $V_{seuil} = \SI{-55}{mV}$.
\begin{enumerate}
    \item Une afférence A génère un PPSE de \SI{+12}{mV} au niveau du cône d'émergence. Une afférence B génère un PPSE de \SI{+8}{mV}. Une afférence inhibitrice C génère un PPSI de \SI{-6}{mV}.
    \item Déterminer si un potentiel d'action est déclenché dans les cas suivants :
    \begin{itemize}
        \item[a)] Stimulation de A seule.
        \item[b)] Stimulation de B seule.
        \item[c)] Stimulation simultanée de A et B.
        \item[d)] Stimulation simultanée de A, B et C.
    \end{itemize}
\end{enumerate}
\end{exemplebox}

\begin{exoresolu}[Analyse de l'exemple chiffré]
\textbf{Données :} $V_{repos} = \SI{-70}{mV}$ ; $V_{seuil} = \SI{-55}{mV}$ $\Rightarrow$ \textbf{Dépolarisation nécessaire} $\Delta V_{nec} = \SI{15}{mV}$.
$PPSE_A = \SI{+12}{mV}$ ; $PPSE_B = \SI{+8}{mV}$ ; $PPSI_C = \SI{-6}{mV}$.

\tcblower

\textbf{Calculs :}
\begin{align*}
\text{a) A seul :} \quad & V_{max} = -70 + 12 = \SI{-58}{mV} \quad (< -55) \quad \rightarrow \textbf{Pas de PA (sous-seuil)}. \\
\text{b) B seul :} \quad & V_{max} = -70 + 8 = \SI{-62}{mV} \quad (< -55) \quad \rightarrow \textbf{Pas de PA (sous-seuil)}. \\
\text{c) A + B (Sommation spatiale) :} \quad & V_{max} = -70 + 12 + 8 = \SI{-50}{mV} \quad (> -55) \quad \rightarrow \textbf{PA DÉCLENCHÉ}. \\
\text{d) A + B + C (Intégration E/I) :} \quad & V_{max} = -70 + 12 + 8 - 6 = \SI{-56}{mV} \quad (< -55) \quad \rightarrow \textbf{Pas de PA (inhibition l'emporte)}.
\end{align*}

\textbf{Conclusion :} La sommation spatiale (c) permet de dépasser le seuil. L'ajout d'une inhibition (d) ramène le potentiel sous le seuil : le neurone "décide" de ne pas transmettre. Cela illustre le \textbf{poids relatif} des entrées excitatrices vs inhibitrices.
\end{exoresolu}

\begin{methode}[Identifier le type de sommation sur un document (Bac)]
Face à un tracé de potentiels de membrane (oscillogramme) ou un schéma de circuit neuronal :
\begin{enumerate}
    \item \textbf{Compte le nombre de sites de stimulation (afférences) :}
    \begin{itemize}
        \item \textbf{1 seul site} stimulé de façon répétée $\rightarrow$ \textbf{Sommation temporelle}.
        \item \textbf{Plusieurs sites} stimulés simultanément $\rightarrow$ \textbf{Sommation spatiale}.
        \item \textbf{Plusieurs sites} stimulés de façon répétée $\rightarrow$ \textbf{Sommation spatio-temporelle}.
    \end{itemize}
    \item \textbf{Observe l'intervalle entre stimulations (si temporelle) :} Doit être court ($< \SI{20}{ms}$ typiquement) pour que les PPSE se chevauchent.
    \item \textbf{Vérifie la polarité :} Seuls les PPS de même signe (tous excitateurs ou tous inhibiteurs) s'additionnent constructivement. Des signes opposés font une \textbf{somme algébrique} (intégration E/I).
    \item \textbf{Repère le cône d'émergence :} C'est le point de mesure virtuel où la décision est prise.
\end{enumerate}
\end{methode}

\courssub{Exercice d'entraînement}

\begin{exercice}[Analyse d'un enregistrement expérimental]
On enregistre l'activité électrique d'un neurone postsynaptique recevant des afférences de deux neurones présynaptiques distincts, $E_1$ (excitateurs) et $I_1$ (inhibiteur). Le potentiel de repos est \SI{-65}{mV}, le seuil est \SI{-50}{mV}.
On obtient les tracés suivants (schématisés) :

\textbf{Figure 5 — Tracés expérimentaux pour l'exercice.}
\begin{center}
\begin{tikzpicture}[scale=0.7, every node/.style={font=\small}]
    % Trace 1: E1 alone
    \draw[->] (0,0) -- (8,0) node[right] {$t$};
    \draw[->] (0,-1) -- (0,3) node[above] {$V_m$};
    \draw[dashed, popmuted] (0,0) -- (8,0);
    \draw[dashed, poporange] (0,1.5) -- (8,1.5) node[right] {Seuil};
    \draw[popblue, thick, domain=1:8, samples=100] plot (\x, {1.2*exp(-(\x-1)/3)*(1-exp(min(3,-(\x-1)/0.5)))});
    \node at (4, -1.5) {Tracé 1 : Stim $E_1$ seul};
    
    % Trace 2: I1 alone
    \begin{scope}[xshift=10cm]
    \draw[->] (0,0) -- (8,0) node[right] {$t$};
    \draw[->] (0,-1) -- (0,3) node[above] {$V_m$};
    \draw[dashed, popmuted] (0,0) -- (8,0);
    \draw[dashed, poporange] (0,1.5) -- (8,1.5) node[right] {Seuil};
    \draw[poppink, thick, domain=1:8, samples=100] plot (\x, {-0.8*exp(-(\x-1)/3)*(1-exp(min(3,-(\x-1)/0.5)))});
    \node at (4, -1.5) {Tracé 2 : Stim $I_1$ seul};
    \end{scope}
    
    % Trace 3: E1 + I1 simultaneous
    \begin{scope}[yshift=-6cm]
    \draw[->] (0,0) -- (8,0) node[right] {$t$};
    \draw[->] (0,-1) -- (0,3) node[above] {$V_m$};
    \draw[dashed, popmuted] (0,0) -- (8,0);
    \draw[dashed, poporange] (0,1.5) -- (8,1.5) node[right] {Seuil};
    \draw[poppurple, ultra thick, domain=1:8, samples=100] plot (\x, {0.4*exp(-(\x-1)/3)*(1-exp(min(3,-(\x-1)/0.5)))});
    \node at (4, -1.5) {Tracé 3 : Stim $E_1 + I_1$ simultanés};
    \end{scope}
    
    % Trace 4: E1 high freq
    \begin{scope}[xshift=10cm, yshift=-6cm]
    \draw[->] (0,0) -- (8,0) node[right] {$t$};
    \draw[->] (0,-1) -- (0,3) node[above] {$V_m$};
    \draw[dashed, popmuted] (0,0) -- (8,0);
    \draw[dashed, poporange] (0,1.5) -- (8,1.5) node[right] {Seuil};
    \draw[popblue, thick, domain=1:8, samples=100] plot (\x, {1.2*exp(-(\x-1)/3)*(1-exp(min(3,-(\x-1)/0.5)))});
    \draw[popblue, thick, domain=2.5:8, samples=100] plot (\x, {1.2*exp(-(\x-2.5)/3)*(1-exp(-(\x-2.5)/0.5))});
    \draw[popblue, thick, domain=4:8, samples=100] plot (\x, {1.2*exp(-(\x-4)/3)*(1-exp(min(3,-(\x-4)/0.5)))});
    % Sum envelope approx
    \draw[poppink, ultra thick, domain=1:5.5, samples=100] plot (\x, {
        ( (\x>=1) ? 1.2*exp(-(\x-1)/3)*(1-exp(min(3,-(\x-1)/0.5))) : 0 ) +
        ( (\x>=2.5) ? 1.2*exp(-(\x-2.5)/3)*(1-exp(-(\x-2.5)/0.5)) : 0 ) +
        ( (\x>=4) ? 1.2*exp(-(\x-4)/3)*(1-exp(min(3,-(\x-4)/0.5))) : 0 )
    });
    \draw[popdark, very thick] (3.5, 1.5) -- (3.6, 2.8) -- (3.8, -0.5) -- (4.0, 0);
    \node at (4, -1.5) {Tracé 4 : Stim $E_1$ à \SI{400}{Hz}};
    \end{scope}
\end{tikzpicture}
\end{center}

\textbf{Questions :}
\begin{enumerate}
    \item Calculer l'amplitude (en mV) du PPSE généré par $E_1$ seul et du PPSI généré par $I_1$ seul (échelle : 1 unité verticale = \SI{10}{mV}).
    \item Quel type de sommation est illustré par le Tracé 3 ? Justifier. Quel est le résultat fonctionnel ?
    \item Quel type de sommation est illustré par le Tracé 4 ? Justifier. Quel est le résultat fonctionnel ?
    \item Expliquer pourquoi le Tracé 1 ne déclenche pas de potentiel d'action alors que le Tracé 4 en déclenche un, bien que ce soit la même afférence $E_1$ qui est stimulée.
\end{enumerate}
\end{exercice}

\begin{corrige}[Exercice d'entraînement]
\textbf{1. Amplitudes :}
L'échelle verticale a pour origine le repos (\SI{-65}{mV}). Le seuil est à \SI{-50}{mV} (donc \SI{15}{mV} au-dessus du repos). Sur le graphique, le seuil est à 1,5 unités. Donc \textbf{1 unité = \SI{10}{mV}}.
\begin{itemize}
    \item Tracé 1 ($E_1$ seul) : Pic à $\approx 1,2$ unités $\rightarrow$ \textbf{PPSE = \SI{+12}{mV}}. Potentiel max = \SI{-53}{mV} (sous-seuil).
    \item Tracé 2 ($I_1$ seul) : Pic négatif à $\approx -0,8$ unités $\rightarrow$ \textbf{PPSI = \SI{-8}{mV}}. Potentiel min = \SI{-73}{mV}.
\end{itemize}

\textbf{2. Tracé 3 :}
\begin{itemize}
    \item \textbf{Type :} \textbf{Sommation spatiale} (deux afférences \textbf{différentes} $E_1$ et $I_1$, stimulées \textbf{simultanément}).
    \item \textbf{Calcul :} $V_{max} = -65 + 12 - 8 = \SI{-61}{mV}$ (ou graphiquement $1,2 - 0,8 = 0,4$ unités = \SI{4}{mV} au-dessus repos $\rightarrow$ \SI{-61}{mV}).
    \item \textbf{Résultat :} \textbf{Pas de potentiel d'action}. L'inhibition annule l'excitation (somme algébrique).
\end{itemize}

\textbf{3. Tracé 4 :}
\begin{itemize}
    \item \textbf{Type :} \textbf{Sommation temporelle} (une \textbf{seule} afférence $E_1$, stimulée \textbf{répétitivement} à haute fréquence : 3 stimuli à \SI{400}{Hz}, intervalle \SI{2,5}{ms}).
    \item \textbf{Mécanisme :} Le 2ème et 3ème PPSE arrivent avant la fin du précédent. Ils s'additionnent (enveloppe rose). Le sommet dépasse le seuil (\SI{1,5}{unité}).
    \item \textbf{Résultat :} \textbf{Déclenchement d'un potentiel d'action} (pic noir).
\end{itemize}

\textbf{4. Explication Tracé 1 vs Tracé 4 :}
Dans le Tracé 1, un \textbf{PPSE unique} (\SI{+12}{mV}) amène le potentiel à \SI{-53}{mV}, \textbf{inférieur au seuil} (\SI{-50}{mV}). La dépolarisation est sous-seuil.
Dans le Tracé 4, la \textbf{haute fréquence} (\SI{400}{Hz}) provoque une \textbf{sommation temporelle} : les charges capacitives successives s'accumulent sur la membrane (constante de temps $\tau_m > \SI{2,5}{ms}$). La dépolarisation cumulée atteint \SI{-46}{mV} (\textbf{supérieur au seuil}), ouvrant les canaux \ce{Na+} voltage-dépendants du cône d'émergence $\rightarrow$ \textbf{Loi du tout ou rien} : un potentiel d'action de amplitude standard est émis.
\end{corrige}

\begin{savaistu}[Le neurone, une machine à calculer ?]
Le concept de \textbf{sommation spatio-temporelle} fait du neurone bien plus qu'un simple "fil électrique". C'est une \textbf{unité de calcul analogique} (certains disent "unité logique floue").
\begin{itemize}
    \item **Pondération synaptique (Poids $w$) :\textbf{ L'efficacité d'une synapse n'est pas fixe. Elle change avec l'apprentissage (potentialisation à long terme - LTP, dépression à long terme - LTD). C'est la base cellulaire de la }mémoire**.
    \item \textbf{Intégration dendritique active :} Les dendrites ne sont pas passives ! Elles possèdent des canaux voltage-dépendants (\ce{Na+}, \ce{Ca^{2+}}, \ce{NMDA}) qui peuvent amplifier localement les PPSE (pics dendritiques), rendant le calcul non-linéaire.
    \item \textbf{Codage de l'information :} Le neurone code l'intensité du stimulus par la \textbf{fréquence} de ses potentiels d'action (codage en fréquence), lui-même déterminé par la \textbf{fréquence et le nombre} des afférences actives (sommation).
\end{itemize}
Au Cameroun, la compréhension de ces mécanismes est cruciale pour la \textbf{neuropharmacologie} : les anticonvulsivants (ex. phénobarbital, valproate) utilisés contre l'épilepsie au CHU de Yaoundé ou à l'Hôpital Laquintinie de Douala augmentent l'inhibition (GABA) ou diminuent l'excitation (Glutamate), déplaçant l'équilibre de la sommation vers le "non-déclenchement".
\end{savaistu}

\coursec{V. Effets des substances sur la transmission synaptique}

Après avoir analysé comment le neurone intègre les messages nerveux par les phénomènes de sommation spatiale, temporelle et spatio-temporelle (section IV), nous abordons maintenant une question majeure de santé publique et de pharmacologie : \textbf{comment des molécules exogènes — médicaments, drogues, toxines, polluants — peuvent-elles perturber, potentialiser ou bloquer la transmission synaptique ?} Comprendre ces mécanismes, c'est comprendre à la fois la base rationnelle des traitements neurologiques et psychiatriques, et les dangers réels des addictions et des intoxications, y compris dans notre contexte camerounais (agriculture, automédication, usage détourné du tramadol).

\courssub{Cadre général : sites d'action et classification des substances}

\begin{definition}[Agoniste et antagoniste]
Une substance exogène qui modifie la transmission synaptique chimique agit toujours sur \textbf{une étape précise} du cycle du neuromédiateur : sa synthèse, son stockage vésiculaire, sa libération (exocytose), sa fixation sur les récepteurs postsynaptiques, sa recapture par le transporteur présynaptique ou sa dégradation enzymatique dans la fente synaptique.

On distingue deux grandes classes fonctionnelles :
\begin{itemize}
    \item \textbf{Un \cle{agoniste}} : substance qui \textbf{mime} l'action du neuromédiateur endogène (agoniste direct : se fixe sur le récepteur et l'active) ou qui \textbf{renforce} son action (agoniste indirect : augmente sa concentration dans la fente en favorisant sa libération ou en bloquant sa recapture/dégradation).
    \item \textbf{Un \cle{antagoniste}} : substance qui \textbf{bloque} l'action du neuromédiateur. Le plus souvent, c'est un \emph{antagoniste compétitif} : il se fixe sur le site de liaison du récepteur sans l'activer, empêchant ainsi le médiateur naturel de s'y fixer. L'effet est réversible et dépend des concentrations relatives.
\end{itemize}
\end{definition}

\begin{methode}[Analyser l'effet d'une substance sur la transmission synaptique]
Pour expliquer rigoureusement l'effet d'une substance (médicament, drogue, toxine), suis toujours ces trois étapes :
\begin{enumerate}
    \item \textbf{Identifier la cible moléculaire précise} : enzyme de synthèse ou de dégradation, protéine de recapture (transporteur), récepteur postsynaptique (préciser le sous-type : nicotinique, muscarinique, GABA\textsubscript{A}, dopaminergique D\textsubscript{2}, opioïde $\mu$, etc.), protéine de libération (complexe SNARE), canaux ioniques présynaptiques (Ca\textsuperscript{2+}, K\textsuperscript{+}).
    \item \textbf{Préciser le mécanisme d'action} : blocage compétitif ou non compétitif, mimétisme (agoniste), inhibition enzymatique (réversible ou irréversible), blocage du transporteur de recapture, modification de la libération (facilitation ou inhibition).
    \item \textbf{Déduire la conséquence sur le signal postsynaptique} : augmentation ou diminution de la fréquence/amplitude des PSP (potentiels postsynaptiques), changement de l'excitabilité du neurone ou de la cellule effectrice (muscle, glande), effets systémiques observables (paralysie, convulsions, sédation, euphorie, troubles autonomes).
\end{enumerate}
\end{methode}

\begin{popfigure}
\begin{tikzpicture}[scale=0.9, every node/.style={font=\small}]
    % Presynaptic terminal
    \draw[popdark, thick, fill=popblueL!30] (0,0) ellipse (2.5cm and 1.2cm);
    \node[popdark] at (0,1.5) {\textbf{Terminaison présynaptique}};
    
    % Mitochondria
    \foreach \x in {-1.8,-0.5,0.8} {
        \draw[poporange, fill=poporangeL!50] (\x,-0.3) ellipse (0.3cm and 0.15cm);
    }
    
    % Vesicles
    \foreach \i in {1,...,12} {
        \pgfmathsetmacro{\vx}{-2.2 + rand*1.8}
        \pgfmathsetmacro{\vy}{-0.1 + rand*0.6}
        \draw[popblue, fill=popblue!30] (\vx,\vy) circle (0.12cm);
    }
    
    % Voltage-gated Ca2+ channels
    \draw[poppurple, thick] (-1.5,1.1) rectangle (-1.0,1.25);
    \node[poppurple, anchor=north] at (-1.25,1.05) {Canaux Ca\textsuperscript{2+} VG};
    
    % Synaptic cleft
    \draw[popmuted, thick] (-2.5,-1.3) -- (2.5,-1.3);
    \draw[popmuted, thick] (-2.5,-1.8) -- (2.5,-1.8);
    \node[popmuted, anchor=north] at (0,-1.35) {\textbf{Fente synaptique (\SI{20}{\nano\meter})}};
    
    % Acetylcholine molecules in cleft
    \foreach \i in {1,...,8} {
        \pgfmathsetmacro{\cx}{-2 + rand*4}
        \pgfmathsetmacro{\cy}{-1.55 + rand*0.2}
        \node[popgreen, rotate=rand*45] at (\cx,\cy) {\ce{ACh}};
    }
    
    % Postsynaptic membrane
    \draw[popdark, thick, fill=popgreenL!20] (-2.5,-2.5) rectangle (2.5,-3.2);
    \node[popdark] at (0,-3.5) {\textbf{Membrane postsynaptique (plaque motrice / dendrite)}};
    
    % Receptors
    \foreach \x in {-1.8,-0.6,0.6,1.8} {
        \draw[popgreen, thick, fill=popgreen!20] (\x,-2.5) rectangle (\x+0.4,-2.9);
        \node[popgreen] at (\x+0.2,-2.7) {Récepteur};
        \node[popgreen, font=\tiny] at (\x+0.2,-2.95) {nAChR};
    }
    
    % Ion channels opening
    \draw[->, popdark, thick] (0.8,-2.9) -- (0.8,-3.3) node[below] {Na\textsuperscript{+}/K\textsuperscript{+} influx};
    
    % Acetylcholinesterase
    \draw[poporange, thick] (-0.2,-1.55) -- (-0.2,-1.8) node[midway, right] {AChE};
    \node[poporange, anchor=north, font=\tiny] at (-0.2,-1.85) {Dégradation};
    
    % ANNOTATIONS: Sites d'action des substances (synapse cholinergique uniquement)
    % 1. Curare - postsynaptic receptor
    \draw[popred, thick, ->] (-2.5,-2.7) -- (-3.5,-2.7);
    \node[popred, anchor=east, align=right] at (-3.6,-2.7) {\textbf{CURARE}\\(Antagoniste nAChR)\\(Blocage fixation ACh)};
    
    % 2. Nicotine - postsynaptic receptor
    \draw[popred, thick, ->] (2.5,-2.7) -- (3.5,-2.7);
    \node[popred, anchor=west, align=left] at (3.6,-2.7) {\textbf{NICOTINE}\\(Agoniste nAChR)\\(Mimétisme ACh)};
    
    % 3. Acetylcholinesterase inhibitors - enzyme in cleft
    \draw[popblue, thick, ->] (-0.2,-1.4) -- (-0.2,-0.5);
    \node[popblue, anchor=south, align=center] at (-0.2,-0.4){\textbf{Inhibiteurs AChE}\\(Organophosphorés, Gaz innervants)\\(Blocage dégradation ACh)};
    
    % 4. Botulinum toxin - vesicle release (SNARE)
    \draw[poppurple, thick, ->] (-1.5,0.2) -- (-1.5,1.0);
    \node[poppurple, anchor=south, align=center] at (-1.5,1.1){\textbf{Toxine botulique}\\(Blocage libération vésiculaire\\(clivage SNARE))};
\end{tikzpicture}
\legende{Schéma d'une synapse chimique (type plaque motrice cholinergique) annoté avec les sites d'action principaux des substances agissant sur l'acétylcholine. Les couleurs distinguent les cibles : \textcolor{popred}{récepteurs postsynaptiques nAChR}, \textcolor{popblue}{enzyme de dégradation (AChE)}, \textcolor{poppurple}{machinerie de libération vésiculaire}. La cocaïne, le tramadol et l'alcool agissent sur d'autres types de synapses (monoaminergiques, GABAergiques/glutamatergiques).}
\end{popfigure}

\courssub{Substances agissant sur la transmission cholinergique (acétylcholine)}

L'acétylcholine (ACh) est le médiateur de la plaque motrice (synapse neuro-musculaire), du système nerveux parasympathique, de certains neurones sympathiques (glandes sudoripares) et de voies centrales (mémoire, éveil). C'est une cible historique de la pharmacologie et de la toxicologie.

\textbf{Le curare : antagoniste compétitif des récepteurs nicotiniques de la plaque motrice.}\par
Le curare est un mélange d'alcaloïdes (principalement la \emph{d-tubocurarine}) extrait de lianes amazoniennes (\emph{Chondrodendron tomentosum}, \emph{Strychnos toxifera}), utilisé traditionnellement comme poison de flèche par les chasseurs amérindiens.

\begin{experience}[Preuve expérimentale du site d'action du curare (préparation gastrocnémien-sciatique de grenouille)]
\textbf{Matériel :} Grenouille, muscle gastrocnémien avec son nerf sciatique attaché, bac à perfusion oxygéné (solution de Ringer), électrodes de stimulation, capteur de force (myographe), curare (solution à \SI{1}{\micro\molar}).
\textbf{Protocole :}
\begin{enumerate}
    \item Stimuler le nerf sciatique par des impulsions électriques supramaximales (\SI{0.5}{\milli\second}, \SI{1}{\hertz}) $\rightarrow$ observer les contractions musculaires twitch régulières.
    \item Perfuser le muscle avec la solution de curare pendant \SI{10}{\minute}.
    \item Stimuler à nouveau le nerf sciatique aux mêmes paramètres $\rightarrow$ \textbf{abolition totale de la contraction}.
    \item Stimuler \textbf{directement le muscle} (électrodes sur le ventre musculaire) aux mêmes paramètres $\rightarrow$ \textbf{contraction normale conservée}.
    \item Laver longuement à la solution de Ringer $\rightarrow$ récupération progressive de la réponse à la stimulation nerveuse.
\end{enumerate}
\textbf{Observations :} Le curare abolit la réponse à la stimulation \emph{nerveuse} mais respecte la réponse à la stimulation \emph{directe} du muscle.
\textbf{Interprétation :} Le curare n'altère ni l'excitabilité de l'axone (potentiel d'action propagé), ni la contractilité du muscle (mécanisme actine-myosine, couplage excitation-contraction). Il bloque spécifiquement la \textbf{transmission de l'influx nerveux au muscle}, c'est-à-dire au niveau de la \textbf{synapse neuro-musculaire (plaque motrice)}.
\textbf{Conclusion :} Le curare est un \textbf{antagoniste compétitif réversible des récepteurs nicotiniques à l'acétylcholine (nAChR)} de la plaque motrice. Il occupe le site de fixation de l'ACh sans ouvrir le canal ionique intrinsèque (Na\textsuperscript{+}/K\textsuperscript{+}), empêchant la dépolarisation de la plaque motrice (PPSE) et donc le déclenchement du potentiel d'action musculaire.
\end{experience}

\begin{propriete}[Effet du curare sur la transmission neuromusculaire]
Le curare paralyse les muscles squelettiques en bloquant les récepteurs nAChR de la plaque motrice. La paralysie est \textbf{réversible} (lavage, élimination rénale) et \textbf{compétitive} : une augmentation de la concentration d'ACh dans la fente (ex. par un inhibiteur de l'acétylcholinestérase comme la néostigmine) peut déplacer l'équilibre et vaincre le blocage. C'est la base de l'antidote en anesthésie.
\end{propriete}

\begin{popfigure}
\begin{tikzpicture}
\begin{axis}[
    width=\linewidth, height=7cm,
    xlabel={Temps (s)}, ylabel={Force de contraction (mN)},
    xmin=0, xmax=30, ymin=0, ymax=1.1,
    xtick={5,10,15,20,25,30}, ytick={0,0.5,1.0},
    legend pos=north east,
    grid=major, grid style={popline},
]
% Normal twitches
\addplot[popblue, thick, samples=300, domain=0:30] {
    (mod(x,1)<0.1) * (exp(-(mod(x,1)-0.05)*50) * 1.0)
};
\addlegendentry{Contrôle (Ringer)}

% Curare: nerve stimulation -> flat line
\addplot[popred, thick, samples=300, domain=0:30] {
    (x<10) * (mod(x,1)<0.1) * (exp(-(mod(x,1)-0.05)*50) * 1.0)
    + (x>=10 && x<20) * 0
    + (x>=20) * (mod(x,1)<0.1) * (exp(-(mod(x,1)-0.05)*50) * 0.9) % recovery after wash
};
\addlegendentry{+ Curare (stim. nerve)}

% Direct muscle stimulation during curare
\addplot[popgreen, dashed, thick, samples=300, domain=10:20] {
    (mod(x,1)<0.1) * (exp(-(mod(x,1)-0.05)*50) * 1.0)
};
\addlegendentry{Stim. directe muscle (pendant curare)}

% Annotations
\node[popred, anchor=north] at (axis cs:15,1.05) {Perfusion curare};
\node[popgreen, anchor=north] at (axis cs:15,0.9) {Lavage Ringer};
\end{axis}
\end{tikzpicture}
\legende{Enregistrement myographique (force vs temps) sur muscle gastrocnémien de grenouille isolé. \textbf{Bleu :} contrôle, stimulation nerveuse \SI{1}{\hertz} $\rightarrow$ twitchs réguliers. \textbf{Rouge :} perfusion de curare (\SI{10}{\minute}, zone grisée) $\rightarrow$ abolition complète de la réponse à la stimulation nerveuse (ligne plate), récupération partielle après lavage. \textbf{Vert pointillé :} stimulation électrique \emph{directe} du muscle pendant la perfusion de curare $\rightarrow$ contractions normales, prouvant l'intégrité du muscle et du couplage excitation-contraction.}
\end{popfigure}

\begin{attention}[Piège fréquent : Curare et Acétylcholinestérase]
\textbf{Erreur :} « Le curare détruit l'acétylcholine » ou « Le curare inhibe l'acétylcholinestérase ».
\textbf{Correction :} Le curare \textbf{ne touche pas à l'acétylcholine} ni à l'enzyme qui la dégrade. Il se fixe \textbf{uniquement sur le récepteur postsynaptique} (nAChR) à la place de l'ACh, sans l'activer. L'ACh est libérée normalement, diffuse normalement, est dégradée normalement par l'AChE, mais \textbf{ne peut plus ouvrir le canal ionique}. C'est un blocage \emph{postsynaptique} pur.
\end{attention}

\textbf{L'atropine : antagoniste des récepteurs muscariniques.}\par
L'atropine, alcaloïde de la belladone (\emph{Atropa belladonna}), bloque sélectivement les \textbf{récepteurs muscariniques} (mAChR, métabotropes, couplés aux protéines G) sans agir sur les récepteurs nicotiniques (canaux ioniques). Elle empêche les effets parasympathiques de l'ACh : \textbf{mydriase} (dilatation pupillaire, blocage du sphincter de l'iris), \textbf{tachycardie} (blocage du frein vagal sur le nœud sinusal), \textbf{sécheresse buccale} (blocage salivaire), \textbf{ralentissement du transit}, \textbf{rétention urinaire}. Usage médical : \textbf{ophtalmologie} (mydriatique pour examen du fond d'œil), \textbf{antidote} des intoxications par organophosphorés (voir ci-dessous), préanesthésie (antisialorrhoïque).

\textbf{Les inhibiteurs de l'acétylcholinestérase (AChE) : excès d'acétylcholine.}\par
L'acétylcholinestérase (AChE) hydrolyse l'ACh en choline + acetate dans la fente synaptique, terminant le signal. Son inhibition entraîne une \textbf{accumulation d'ACh} et une \textbf{surstimulation prolongée} de \emph{tous} les récepteurs cholinergiques (nicotiniques \textbf{et} muscariniques).

\begin{itemize}
    \item \textbf{Inhibiteurs réversibles (carbamates) :} Néostigmine, pyridostigmine (traitement de la \textbf{myasthénie grave}, maladie auto-immune anti-récepteurs nAChR), physostigmine (glaucome). Usage thérapeutique contrôlé.
    \item \textbf{Inhibiteurs \emph{irréversibles} (organophosphorés) :} \textbf{Insecticides} (parathion, malathion, chlorpyrifos — très utilisés dans les plantations de cacao et de coton au Cameroun), \textbf{gaz innervants} (sarin, VX, tabun — armes chimiques). Ils phosphorylent le site actif de l'AChE de façon quasi définitive (vieillissement). La récupération nécessite la \textbf{resynthèse de l'enzyme} (jours à semaines).
\end{itemize}

\begin{exemplebox}[Intoxication par organophosphorés : contexte agricole camerounais]
Dans les zones de production cacaoyère (Sud, Centre, Est) ou cotonnière (Nord), les applicateurs non formés ou mal équipés (absence de gants, masques, combinaisons) manipulent des concentrés d'insecticides organophosphorés. L'absorption cutanée et inhalée provoque un \textbf{syndrome cholinergique aigu} : \textbf{SLUDGE} (Salivation, Larmoiement, Urination, Diarrhée, Gastro-intestinal cramps, Émétisme) + \textbf{myosis} (pupilles ponctiformes), \textbf{bradycardie}, \textbf{bronchorrhée}, \textbf{fasciculations} musculaires (surstimulation nAChR), \textbf{faiblesse} puis \textbf{paralysie respiratoire} (défaillance centre respiratoire + paralysie muscles intercostaux/diaphragme). \textbf{Urgence vitale :} Atropine (antagoniste muscarinique, dose massive itérative) + Oximes (pralidoxime, réactive l'AChE phosphorylée si donné tôt) + Assistance respiratoire. \textbf{Prévention :} Formation, EPI obligatoires, substitution par des insecticides moins toxiques (néonicotinoïdes, régulateurs de croissance), respect des délais de rentrée.
\end{exemplebox}

\textbf{La nicotine : agoniste des récepteurs nicotiniques.}\par
Alcaloïde du tabac (\emph{Nicotiana tabacum}), la nicotine est un \textbf{agoniste direct} des récepteurs nAChR (sous-unités $\alpha$4$\beta$2 surtout dans le système de récompense, $\alpha$7 dans l'hippocampe, $\alpha$3$\beta$4 dans les ganglions autonomes).
\begin{itemize}
    \item \textbf{Effets aigus :} Stimulation puis désensibilisation des nAChR. Libération de \textbf{dopamine} dans le noyau accumbens (circuit de récompense) $\rightarrow$ plaisir, renforcement. Libération d'adrénaline (médullosurrénale) $\rightarrow$ tachycardie, hypertension, hyperglycémie. Stimulation cognitive transitoire (attention, mémoire de travail).
    \item \textbf{Effets chroniques :} \textbf{Tolérance} (régulation à la baisse des récepteurs, désensibilisation), \textbf{dépendance} physique et psychique forte (critères DSM-5). Sevrage : irritabilité, anxiété, trouble de l'attention, fringale.
    \item \textbf{Contexte :} Tabagisme actif/passif : 1\textsuperscript{re} cause de mortalité évitable (cancers, BPCO, maladies cardiovasculaires). Au Cameroun, prévalence tabagique masculine élevée ; lutte antitabac (loi 2006, paquet neutre, taxation).
\end{itemize}

\courssub{Substances agissant sur les systèmes monoaminergiques et opioïdes}

\textbf{La cocaïne : blocage de la recapture de la dopamine et de la noradrénaline.}\par
Alcaloïde de \emph{Erythroxylum coca}. La cocaïne se lie au \textbf{transporteur de la dopamine (DAT)} et au \textbf{transporteur de la noradrénaline (NET)} sur la membrane présynaptique, les bloquant de façon compétitive.
\begin{itemize}
    \item \textbf{Conséquence :} Accumulation massive de DA et NA dans la fente synaptique $\rightarrow$ surstimulation prolongée des récepteurs D\textsubscript{1}/D\textsubscript{2} (voie mésolimbique : euphorie, mégalomanie, hyperactivité) et $\alpha$/$\beta$-adrénergiques (sympathomimétique : tachycardie, hypertension, vasoconstriction, mydriase, hyperthermie, risque d'infarctus/AVC).
    \item \textbf{Dépendance :} Très forte (renforcement par pic dopaminergique rapide). \textbf{Crack} (base libre fumable) : passage hémato-encéphalique en \SI{<10}{\second}, pic intense, chute brutale $\rightarrow$ craving extrême.
\end{itemize}

\textbf{Le tramadol : opioïde atypique à double mécanisme.}\par
Analgésique de palier II (OMS), détourné massivement au Cameroun et en Afrique de l'Ouest comme « drogue du pauvre » ou « coupe » pour travailleurs (motos-taxis, cultivateurs, étudiants).
\begin{itemize}
    \item \textbf{Mécanisme 1 :} Agoniste faible des \textbf{récepteurs opioïdes $\mu$} (morphiniques) $\rightarrow$ analgésie, sédation, dépression respiratoire (dose-dépendante), constipation, euphorie.
    \item \textbf{Mécanisme 2 :} Inhibiteur de la \textbf{recapture de la sérotonine (SERT)} et de la \textbf{noradrénaline (NET)} $\rightarrow$ effet antidépresseur/anxiolytique, mais risque de \textbf{syndrome sérotoninergique} (hyperthermie, rigidité, clonus, agitation) en association avec d'autres sérotoninergiques (ISRS, tramadol + antidépresseurs).
    \item \textbf{Dangers du détournement :} Surdose $\rightarrow$ \textbf{dépression respiratoire fatale} (réversible par naloxone), \textbf{convulsions} (seuil épileptogène abaissé par action sérotoninergique/noradrénergique), dépendance sévère (double sevrage opioïde + monoaminergique). \textbf{Interdit de vente libre au Cameroun depuis 2017}, mais circuit illicite persistant.
\end{itemize}

\courssub{L'alcool (éthanol) : potentialisation GABAergique et inhibition glutamatergique}

L'éthanol est une petite molécule amphiphile qui traverse librement la barrière hémato-encéphalique. Il n'a pas de récepteur spécifique mais module allostériquement plusieurs canaux ioniques :
\begin{itemize}
    \item \textbf{Potentialisation des récepteurs GABA\textsubscript{A}} (canal Cl\textsuperscript{-} inhibiteur) : augmente la fréquence d'ouverture du canal $\rightarrow$ \textbf{hyperpolarisation} accrue $\rightarrow$ sédation, anxiolyse, myorelaxation, amnésie antérograde, \textbf{dépression respiratoire} (dose élevée), anticonvulsivant.
    \item \textbf{Inhibition des récepteurs NMDA (glutamate)} : blocage du canal cationique (Na\textsuperscript{+}/Ca\textsuperscript{2+}) $\rightarrow$ trouble de la plasticité synaptique (LTP), amnésie, ataxie, neurotoxicité fœtale (syndrome d'alcoolisation fœtale).
    \item \textbf{Libération de dopamine} (voie mésolimbique, via inhibition GABAergique des neurones inhibiteurs de la VTA) $\rightarrow$ euphorie, désinhibition, renforcement.
\end{itemize}
\textbf{Effets aigus dose-dépendants :} \SI{0,2}{\gram/\liter} (désinhibition) $\rightarrow$ \SI{0,5}{\gram/\liter} (troubles coordination, réflexes) $\rightarrow$ \SI{1,0}{\gram/\liter} (ivresse manifeste) $\rightarrow$ \SI{2,0}{\gram/\liter} (coma éthylique) $\rightarrow$ \SI{>3,5}{\gram/\liter} (risque vital : dépression respiratoire, hypothermie, inhalation vomissements). \textbf{Sevrage aigu :} syndrome de sevrage (tremblements, anxiété, convulsions, delirium tremens) par \textbf{hyperexcitabilité} (régulation à la hausse NMDA, à la baisse GABA\textsubscript{A}).

\courssub{Sensibilisation : dépendance, tolérance, sevrage et prévention}

\begin{definition}[Dépendance, tolérance, sevrage]
\begin{itemize}
    \item \textbf{Dépendance (addiction) :} Pathologie cérébrale chronique, récidivante, définie par la \textbf{perte de contrôle} de la consommation, la \textbf{compulsion} (recherche obsessionnelle), la \textbf{continuation malgré les conséquences négatives} (sanitaires, scolaires, familiales, juridiques). Base neurobiologique : remodelage durable du circuit de la récompense (dopamine, glutamate, CRF, dynorphine) et du contrôle exécutif (cortex préfrontal).
    \item \textbf{Tolérance :} Nécessité d'augmenter les doses pour obtenir le même effet (tolérance pharmacocinétique : induction enzymatique CYP450 ; tolérance pharmacodynamique : désensibilisation/réduction du nombre de récepteurs, adaptation homéostatique).
    \item \textbf{Sevrage :} Syndrome physique/psychique à l'arrêt ou à la réduction de la substance, signe d'adaptation neurobiologique. Opposé aux effets aigus (ex. sevrage alcool/benzos : hyperactivité, convulsions ; sevrage opioïdes : douleurs, diarrhée, anxiété ; sevrage cocaïne/nicotine : dysphorie, fatigue, fringale).
\end{itemize}
\end{definition}

\begin{aretenir}[Messages clés de prévention (contexte Terminale / Jeunesse camerounaise)]
\begin{itemize}
    \item \textbf{Le cerveau en maturation (jusqu'à $\approx$ 25 ans) est hyper-vulnérable} : l'alcool, le cannabis, le tramadol, la cocaïne perturbent la myélinisation, l'élagage synaptique, la maturation préfrontale $\rightarrow$ troubles cognitifs durables, risque accru de troubles psychiatriques (schizophrénie, dépression), échec scolaire.
    \item \textbf{Automédication = danger :} Tramadol, benzodiazépines (diazépam, clonazépam), codéine, antihistaminiques sédatifs pris sans ordonnance $\rightarrow$ dépendance, surdose, interactions mortelles (ex. tramadol + alcool = dépression respiratoire majorée).
    \item \textbf{Médicaments = outils thérapeutiques, pas récréatifs :} La morphine soulage le cancéreux, le méthylphénidate permet à l'enfant TDAH de scolariser, la buprénorphine/méthadone (traitement de substitution aux opioïdes) sont \textbf{efficaces et sûrs} \emph{uniquement} prescrits, dosés, surveillés par un médecin.
    \item \textbf{Environnement :} Lutte contre les pesticides illicites (organophosphorés interdits : parathion, méthamidophos), port des EPI obligatoire, surveillance des résidus dans l'alimentation.
    \item \textbf{Ressources :} Centre de prise en charge des addictions (CPA) à l'Hôpital Central / Jamot (Yaoundé), Hôpital Laquintinie (Douala), associations (CNAD, Blue Cross), numéro vert \textbf{1510} (urgence SAMU) / \textbf{119} (enfance en danger).
\end{itemize}
\end{aretenir}

\begin{savaistu}[Du poison au médicament : l'histoire du curare]
Le curare, poison de chasse amazonien (XVII\textsuperscript{e} s.), a fasciné les physiologistes. \textbf{Claude Bernard} (1842-1856) démontre sur la grenouille qu'il paralyse sans affecter le muscle ni le nerf, localisant l'action à la « plaque motrice ». En 1942, \textbf{Griffith et Johnson} l'utilisent pour la première fois comme \textbf{myorelaxant en chirurgie} (intubation, chirurgie abdominale/thoracique sans respiration artificielle impossible avant). Aujourd'hui, ses dérivés de synthèse (atracurium, cisatracurium, rocuronium, vecuronium) sont des \textbf{médicaments essentiels} (liste OMS) en anesthésie-réanimation, à action plus courte, prévisible, sans effets histaminergiques majeurs. La néostigmine (inhibiteur AChE réversible) en est l'antidote de routine. \textbf{Exemple parfait de la toxicologie devenue pharmacologie rationnelle.}
\end{savaistu}

\begin{popfigure}
\begin{center}
\begin{tabularx}{\linewidth}{|>{\bfseries}l|X|X|X|}\hline
\rowcolor{popblueL}
\textbf{Substance} & \textbf{Cible synaptique précise} & \textbf{Mécanisme d'action} & \textbf{Effet physiologique principal / Conséquence} \\\hline
\textbf{Curare} & Récepteur nAChR (plaque motrice, ganglions) & Antagoniste compétitif (blocage fixation ACh) & Paralysie flasque musculaire (arrêt respiratoire), blocage transmission ganglionnaire \\\hline
\textbf{Atropine} & Récepteurs muscariniques (mAChR) & Antagoniste compétitif & Mydriase, tachycardie, sécheresse buccale, ralentissement transit, rétention urinaire \\\hline
\textbf{Inhibiteurs AChE (Organophosphorés, Névrogaz)} & Acétylcholinestérase (fente synaptique) & Inhibition irréversible (phosphorylation) $\rightarrow$ accumulation ACh & Syndrome cholinergique : SLUDGE, myosis, bradycardie, fasciculations, paralysie respiratoire, convulsions \\\hline
\textbf{Nicotine} & Récepteurs nAChR ($\alpha$4$\beta$2, $\alpha$7, ganglions) & Agoniste direct (puis désensibilisation) & Libération DA (récompense), stimulation puis blocage ganglionnaire, tachycardie, HTA, dépendance forte \\\hline
\textbf{Cocaïne} & Transporteurs DAT (DA), NET (NA) & Blocage recapture $\rightarrow$ accumulation fente & Euphorie, mégalomanie, sympathomimétique (HTA, tachycardie, vasoconstriction), risque AVC/IDM, dépendance sévère \\\hline
\textbf{Tramadol} & Récepteurs opioïdes $\mu$ + Transporteurs SERT/NET & Agoniste $\mu$ + Inhibition recapture 5-HT/NA & Analgésie, sédation, euphorie, risque convulsions, syndrome sérotoninergique, dépression respiratoire (surdose), dépendance \\\hline
\textbf{Alcool (Éthanol)} & Récepteurs GABA\textsubscript{A} (potentialisation) + NMDA (inhibition) & Modulation allostérique canaux ioniques & Sédation, désinhibition, ataxie, amnésie, coma éthylique, dépression respiratoire, sevrage convulsif \\\hline
\end{tabularx}
\end{center}
\legende{Tableau de synthèse : substances, cibles synaptiques, mécanismes et effets majeurs. À connaître par cœur pour le Bac (colonne « Cible » et « Mécanisme » sont les plus discriminantes).}
\end{popfigure}

\begin{exoresolu}[Analyse d'une intoxication aiguë]
\textbf{Énoncé :} Un applicateur de produits phytosanitaires de 32 ans, travaillant sans protection dans une plantation de cacao au Centre-Cameroun, est admis aux urgences pour : salivation abondante, larmoiement, bronchorrhée, diarrhée, vomissements, myosis bilatéral, bradycardie à 45 bpm, fasciculations généralisées, confusion, difficulté respiratoire. Le produit utilisé est un insecticide à base de \textbf{chlorpyrifos-ethyl} (organophosphoré).\\
\textbf{Questions :}
\begin{enumerate}
    \item Identifier la cible moléculaire du chlorpyrifos et le type d'inhibition enzymatique.
    \item Expliquer le mécanisme physiopathologique reliant cette cible aux signes cliniques observés (classer en effets muscariniques, nicotiniques, centraux).
    \item Proposer le traitement antidotique spécifique d'urgence et son mécanisme d'action.
    \item Pourquoi l'atropine seule ne suffit-elle pas à traiter la paralysie respiratoire ?
\end{enumerate}
\tcblower
\textbf{Solution détaillée :}
\begin{enumerate}
    \item \textbf{Cible :} Acétylcholinestérase (AChE) synaptique (et plasmatique). \textbf{Type :} Inhibition \textbf{irréversible} (phosphorylation du sérine du site actif, « vieillissement » rapide pour le chlorpyrifos-oxon, métabolite actif).
    \item \textbf{Mécanisme :} L'AChE ne dégrade plus l'ACh $\rightarrow$ accumulation d'ACh dans \emph{toutes} les fentes cholinergiques $\rightarrow$ surstimulation prolongée des récepteurs.
    \begin{itemize}
        \item \textbf{Effets muscariniques (parasympathomimétiques) :} SLUDGE (Salivation, Larmoiement, Urination, Diarrhée, Gastro-intestinal cramps, Émétisme) + Bronchorrhée + Bradycardie + Myosis. (Récepteurs mAChR sur effecteurs parasympathiques).
        \item \textbf{Effets nicotiniques :} Fasciculations (plaques motrices), puis faiblesse/paralysie musculaire (dépolarisation prolongée $\rightarrow$ inactivation canaux $Na_v^+$ $\rightarrow$ blocage de la transmission), hypertension/tachycardie transitoires (ganglions sympathiques).
        \item \textbf{Effets centraux :} Confusion, agitation, convulsions, dépression respiratoire centrale (accumulation ACh dans SNC).
    \end{itemize}
    \item \textbf{Traitement d'urgence :}
    \begin{itemize}
        \item \textbf{Atropine sulfate IV :} \SI{1-2}{\milli\gram} (adulte) toutes les \SI{5-10}{\minute} jusqu'à disparition des signes muscariniques (sécheresse buccale, FR > 20/min, FC > 80/min). \textbf{Mécanisme :} Antagoniste compétitif des récepteurs muscariniques $\rightarrow$ bloque les effets de l'excès d'ACh sur les effecteurs parasympathiques.
        \item \textbf{Pralidoxime (PAM) IV :} \SI{1-2}{\gram} en perfusion lente, renouvelable. \textbf{Mécanisme :} Oxime $\rightarrow$ réactive l'AChE phosphorylée \textbf{avant vieillissement} (décrochage du groupement phosphorylé) $\rightarrow$ restaure la dégradation de l'ACh. \textbf{Efficace si donné précocement (< 24-48h).}
        \item \textbf{Diazépam IV :} Prévention/traitement des convulsions centrales (GABA\textsubscript{A}).
        \item \textbf{Assistance respiratoire :} Intubation + ventilation mécanique si paralysie muscles respiratoires / dépression centrale.
    \end{itemize}
    \item \textbf{L'atropine ne bloque que les récepteurs muscariniques.} La paralysie respiratoire a deux composantes : (1) \textbf{Centrale} (dépression bulbaire) $\rightarrow$ atropine ne passe pas bien la BHE, diazépam/ventilation nécessaires. (2) \textbf{Périphérique} (blocage de la plaque motrice par dépolarisation prolongée / inactivation canaux $Na_v^+$) $\rightarrow$ récepteurs \textbf{nicotiniques}, insensibles à l'atropine. Seule la \textbf{réactivation de l'AChE par la pralidoxime} (ou la resynthèse enzymatique sur semaines) peut lever ce blocage nicotinique. D'où l'association vitale Atropine + Pralidoxime + Ventilation.
\end{enumerate}
\end{exoresolu}

\begin{exercice}[Entraînement Bac : Mécanismes d'action]
\textbf{1.} Une patiente de 65 ans, traitée pour myasthénie grave par pyridostigmine (inhibiteur réversible de l'AChE), présente une aggravation brutale : ptôsis, diplopie, dysphagie, faiblesse membres, détresse respiratoire. Le diagnostic hésite entre \textbf{crise myasthénique} (aggravation maladie, anticorps anti-récepteurs) et \textbf{surdosage cholinergique} (excès de pyridostigmine). On réalise un \textbf{test à l'edrophonium} (inhibiteur AChE ultra-court, \SI{2}{\minute}) : injection IV de \SI{2}{\milli\gram} $\rightarrow$ \textbf{amélioration spectaculaire transitoire} de la force.
\begin{itemize}
    \item[a)] Conclure sur le diagnostic (crise ou surdosage) et justifier par le mécanisme.
    \item[b)] Quel serait l'effet de l'edrophonium en cas de surdosage ? Pourquoi ?
\end{itemize}

\textbf{2.} Schématiser (schéma fonctionnel fléché) le mécanisme d'action de la cocaïne sur une synapse dopaminergique du noyau accumbens (terminée présynaptique : vésicules, DAT, mitochondries ; fente : DA ; terminée postsynaptique : récepteurs D\textsubscript{1}, D\textsubscript{2}). Indiquer le site d'action, le sens de la flèche de concentration [DA] fente, et la conséquence sur l'activité du neurone postsynaptique (voie directe / indirecte).

\textbf{3.} Expliquer pourquoi l'association \textbf{Tramadol + Fluoxétine (ISRS)} est formellement contre-indiquée (risque de syndrome sérotoninergique). Citer 3 signes cliniques majeurs de ce syndrome.
\end{exercice}

\begin{corrige}[Exercice n°1-3]
\textbf{1.} \textbf{Diagnostic : Crise myasthénique (sous-dosage relatif / aggravation maladie).} \textbf{Justification :} La pyridostigmine inhibe l'AChE $\rightarrow$ augmente [ACh] fente $\rightarrow$ compense la perte de récepteurs nAChR. En crise, les récepteurs fonctionnels sont trop peu nombreux $\rightarrow$ [ACh] efficace insuffisante. L'edrophonium (inhibiteur AChE supplémentaire) augmente encore [ACh] $\rightarrow$ sature les récepteurs restants $\rightarrow$ améliore la transmission $\rightarrow$ force augmente. \textbf{En cas de surdosage cholinergique :} [ACh] déjà excessive $\rightarrow$ dépolarisation prolongée plaque motrice $\rightarrow$ inactivation canaux $Na_v^+$ $\rightarrow$ blocage de la transmission (paralysie par dépoliarisation). L'edrophonium aggrave le blocage $\rightarrow$ \textbf{aggravation de la faiblesse}. Le test est donc discriminant.
\textbf{2.} Schéma : Terminaison présynaptique $\xrightarrow{\text{libération}}$ DA dans fente $\xrightarrow{\text{DAT (bloqué par cocaïne)}}$ Recapture inhibée $\rightarrow$ \textbf{[DA] fente $\uparrow\uparrow$} $\rightarrow$ Fixation massive sur D\textsubscript{1} (voie directe, excitation, $\uparrow$ cAMP) et D\textsubscript{2} (voie indirecte, inhibition, $\downarrow$ cAMP) $\rightarrow$ \textbf{Activation nette du circuit de récompense} (euphorie, renforcement).
\textbf{3.} Tramadol inhibe recapture sérotonine (SERT) + Fluoxétine inhibe recapture sérotonine (SERT) $\rightarrow$ \textbf{Accumulation massive sérotonine fente} $\rightarrow$ Surstimulation récepteurs 5-HT\textsubscript{1A}, 5-HT\textsubscript{2A} $\rightarrow$ \textbf{Syndrome sérotoninergique}. \textbf{Triade majeure :} (1) \textbf{Hyperthermie} ($>$\SI{38,5}{\celsius}, souvent $>$\SI{40}{\celsius}), (2) \textbf{Rigidité musculaire} (hypertonie, clonus, hyperréflexie), (3) \textbf{Troubles neurovégétatifs/mentaux} (agitation, confusion, tachycardie, hypertension, mydriase, diarrhée, sudation). Urgence vitale : arrêt des sérotoninergiques, cyproheptadine (antagoniste 5-HT\textsubscript{2A}), refroidissement, curarisation/ventilation si sévère.
\end{corrige}

\courssub{Synthèse et sensibilisation : une responsabilité partagée}

La transmission synaptique, par sa complexité moléculaire (récepteurs, transporteurs, enzymes), offre une \textbf{cible pharmacologique d'une précision extraordinaire} — base de la neurologie, de la psychiatrie, de l'anesthésie, de la douleur — mais aussi une \textbf{vulnérabilité majeure} aux xénobiotiques. Chaque substance étudiée (curare, atropine, organophosphorés, nicotine, cocaïne, tramadol, alcool) illustre un principe universel : \textbf{modifier une étape du cycle du médiateur change le message nerveux, parfois de façon irréversible ou fatale.}

En tant que futurs scientifiques, professionnels de santé, agronomes, décideurs ou simples citoyens camerounais, vous devez :
\begin{itemize}
    \item \textbf{Connaître les mécanismes} pour prescrire, dispenser, manipuler, légiférer en toute rigueur.
    \item \textbf{Reconnaître les syndromes toxiques} (cholinergique, opioïde, sérotoninergique, sevrage alcoolique) pour alerter et orienter vers les structures adaptées (CPA, urgences, centre antipoison).
    \item \textbf{Promouvoir la prévention} : lutte contre le tabagisme, l'alcoolisation fœtale, le détournement du tramadol, l'usage non protégé des pesticides ; éducation à la santé dès le lycée ; lutte contre la stigmatisation des usagers (la dépendance est une maladie, pas un vice).
    \item \textbf{Valoriser le bon usage du médicament} : la morphine soulage le cancéreux, le méthylphénidate permet à l'enfant TDAH de scolariser, la buprénorphine sort de la rue l'héroïnomane. \textbf{La science bien comprise protège ; l'ignorance tue.}
\end{itemize}

\begin{aretenir}[Check-list révision Bac - Section V]
\begin{itemize}
    \item \cle{Agoniste} (mime/renforce) vs \cle{Antagoniste} (bloque) : définition, exemples pour chaque substance.
    \item \textbf{Cible précise} pour chaque substance : Récepteur (sous-type), Enzyme, Transporteur.
    \item \textbf{Mécanisme} : Blocage compétitif, Mimétisme, Inhibition enzymatique (réversible/irréversible), Blocage recapture.
    \item \textbf{Expérience clé :} Curare sur grenouille (stim. nerve vs stim. directe) $\rightarrow$ preuve site postsynaptique.
    \item \textbf{Signes cliniques types} : Syndrome cholinergique (SLUDGE + myosis + fasciculations), Syndrome opioïde (mydriase $\rightarrow$ myosis, dépression respiratoire), Syndrome sérotoninergique (Hyperthermie + Rigidité + Troubles neurovégétatifs), Ivresse/Sevrage alcoolique.
    \item \textbf{Antidotes :} Atropine (muscarinique), Pralidoxime (réactive AChE), Naloxone (antagoniste opioïde $\mu$), Diazépam (convulsions sevrage/GABA), Cyproheptadine (sérotoninergique).
    \item \textbf{Contexte Cameroun :} Pesticides organophosphorés (cacao/coton), Tramadol détourné, Tabagisme/Alcool jeunes, Accès aux soins addictions.
\end{itemize}
\end{aretenir}

\coursec{VI. Synthèse et sensibilisation}

Dans la section précédente, tu as découvert, substance par substance, comment des molécules aussi diverses que le curare, la nicotine, la cocaïne, le tramadol ou certains insecticides viennent perturber la transmission synaptique. Il est temps de prendre du recul : toutes ces observations obéissent à une même logique. Cette dernière section a trois objectifs : \textbf{organiser} tes connaissances en un bilan comparatif clair, \textbf{relier} chaque perturbation synaptique aux troubles observés chez l'organisme entier, et te donner les arguments scientifiques d'une véritable \textbf{sensibilisation}, celle qu'un jury de Baccalauréat attend de toi.

\courssub{Bilan comparatif : synapses chimiques et synapses électriques}

Commençons par l'essentiel. Au terme de cette leçon, tu dois être capable de comparer sans hésiter les deux grands modes de transmission synaptique.

\begin{center}
\begin{tabularx}{\linewidth}{|l|X|X|}
\hline
\rowcolor{popblueL} \textbf{Critère} & \textbf{Synapse chimique} & \textbf{Synapse électrique} \\
\hline
Structure & Fente synaptique de \SI{20}{\nano\metre} à \SI{40}{\nano\metre} ; vésicules synaptiques, récepteurs membranaires postsynaptiques & Jonctions communicantes (\emph{gap junctions}) : canaux protéiques reliant directement les cytoplasmes ; espace d'environ \SI{3}{\nano\metre} \\
\hline
Médiateur & Neuromédiateur (acétylcholine, dopamine, GABA\ldots) libéré par exocytose & Aucun : passage direct du courant (ions) d'une cellule à l'autre \\
\hline
Sens de transmission & Unidirectionnel (pré $\rightarrow$ post) & Souvent bidirectionnel \\
\hline
Vitesse & Retard synaptique d'environ \SI{0,5}{\milli\second} & Quasi instantanée \\
\hline
Plasticité & Modulable : excitatrice ou inhibitrice, sommable, fatigable, cible des substances & Peu modulable, très fiable \\
\hline
Localisation & Majoritaire dans le système nerveux des Vertébrés (synapses neuro-neuroniques, plaque motrice, synapse neuro-glandulaire) & Certains centres nerveux, muscle cardiaque, muscle lisse \\
\hline
\end{tabularx}
\end{center}

\begin{aretenir}[L'idée directrice]
La \cle{synapse chimique} est le type dominant chez l'Homme. Sa richesse fonctionnelle (médiateur, récepteurs spécifiques, enzymes de dégradation, pompes de recapture) en fait un point de contrôle remarquable de la communication nerveuse\ldots mais aussi une \cle{cible privilégiée} de toutes les substances psychoactives, toxines et médicaments. Chaque étape de son fonctionnement peut être bloquée, mimée ou amplifiée par une molécule étrangère.
\end{aretenir}

\begin{propriete}[Principe général d'action des substances psychoactives]
Toute substance psychoactive agit en modifiant \textbf{au moins une étape} de la transmission synaptique :
\begin{enumerate}
\item la \textbf{libération} du neuromédiateur (exocytose bloquée ou provoquée) ;
\item la \textbf{fixation} sur les récepteurs postsynaptiques (agoniste qui mime le médiateur, ou antagoniste qui bloque le récepteur) ;
\item la \textbf{recapture} du médiateur par l'élément présynaptique ;
\item sa \textbf{dégradation} enzymatique dans la fente synaptique.
\end{enumerate}
Dans tous les cas, la conséquence est une modification — excès ou déficit — de l'activité du neurone postsynaptique.
\end{propriete}

\courssub{De la perturbation synaptique aux troubles de l'organisme}

Le raisonnement auquel tu dois t'entraîner est une chaîne causale en trois maillons : \emph{substance} $\rightarrow$ \emph{perturbation synaptique} $\rightarrow$ \emph{trouble observable}. Construisons-le sur quelques cas emblématiques.

\textbf{Paralysie flasque : le curare et le botulisme.} Le curare, poison des flèches de chasse d'Amazonie, est un antagoniste compétitif des récepteurs nicotiniques de l'acétylcholine au niveau de la \cle{plaque motrice}. Le médiateur est libéré normalement mais ne peut plus se fixer : le muscle ne reçoit plus d'ordre de contraction. L'animal — ou l'Homme intoxiqué — présente une paralysie dite \emph{flasque} (muscles relâchés), qui gagne les muscles respiratoires et peut entraîner la mort par asphyxie.

\begin{savaistu}[Complément utile au Bac : le botulisme]
Le botulisme est une intoxication due à une toxine sécrétée par la bactérie \emph{Clostridium botulinum}, qui se développe dans les conserves mal stérilisées. La \cle{toxine botulique} pénètre dans les terminaisons nerveuses cholinergiques et \textbf{détruit les protéines nécessaires à l'exocytose} : l'acétylcholine n'est plus libérée dans la fente synaptique. Résultat : une paralysie flasque progressive (troubles de la vue, de la déglutition, puis paralysie respiratoire). À dose infime et contrôlée, cette même toxine (le « Botox ») est utilisée en médecine et en esthétique pour relâcher localement certains muscles. Ce mécanisme illustre parfaitement le premier maillon de la propriété générale : le blocage de la \textbf{libération} du médiateur.
\end{savaistu}

\textbf{Contractures et tétanie : les insecticides.} De nombreux insecticides (organophosphorés, carbamates) sont des \textbf{inhibiteurs de l'acétylcholinestérase}, l'enzyme qui détruit normalement l'acétylcholine dans la fente synaptique. Le médiateur s'accumule alors et stimule en permanence le muscle : contractions involontaires, fasciculations, puis tétanie (contraction soutenue) pouvant bloquer la respiration. C'est le mécanisme inverse du curare, pour un danger tout aussi mortel — d'où les précautions d'usage des pesticides dans les exploitations agricoles (gants, masque, lavage des mains).

\begin{savaistu}[Complément utile au Bac : la toxine tétanique]
La toxine sécrétée par \emph{Clostridium tetani} (tétanos, contracté par des plaies souillées de terre) remonte le long des axones moteurs jusqu'à la moelle épinière, où elle \textbf{bloque la libération des médiateurs inhibiteurs} (glycine, GABA) par les interneurones inhibiteurs. Privés de frein, les motoneurones déchargent sans contrôle : d'où les contractures douloureuses généralisées (trismus, opisthotonos). La vaccination antitétanique, obligatoire et rappelée lors des consultations prénatales au Cameroun, prévient efficacement cette maladie.
\end{savaistu}

\textbf{Dépendance : nicotine, cocaïne, tramadol.} Ces substances détournent le \cle{circuit de la récompense}, un ensemble de neurones du cerveau (notamment l'aire tegmentale ventrale et le noyau accumbens) dont le neuromédiateur clé est la \cle{dopamine}. Naturellement, ce circuit est activé par les comportements utiles à la survie (manger, boire, se reproduire) et produit une sensation de plaisir qui incite à les renouveler. La plupart des drogues provoquent une libération ou une accumulation \emph{anormale et massive} de dopamine dans les synapses de ce circuit : la cocaïne bloque la recapture de la dopamine, la nicotine active des récepteurs qui stimulent les neurones dopaminergiques, les opioïdes lèvent un frein inhibiteur sur ces mêmes neurones. Le plaisir artificiel, bien plus intense que les plaisirs naturels, « reprogramme » le circuit : la consommation devient une priorité biologique. C'est la \cle{dépendance}.

\textbf{Dépression respiratoire : opioïdes et alcool.} Les opioïdes (morphine, héroïne, tramadol à forte dose) se fixent sur les récepteurs opioïdes des neurones des centres respiratoires du tronc cérébral et provoquent leur hyperpolarisation : l'activité de ces centres s'effondre. L'alcool, à forte concentration, potentialise quant à lui la transmission inhibitrice à GABA et réduit l'excitabilité neuronale globale. Dans les deux cas, le rythme respiratoire ralentit, puis peut s'arrêter : c'est la cause principale des décès par surdose.

\begin{exemplebox}[Le tramadol, un danger bien réel chez les jeunes]
Le tramadol est un \textbf{antalgique opioïde} délivré uniquement sur ordonnance, indiqué dans les douleurs intenses (chirurgie, cancer). Détourné en automédication — parfois pour « tenir » au travail, au sport ou en boîte de nuit — il provoque à forte dose une \textbf{dépression respiratoire}, des convulsions, des comas et des décès, particulièrement chez les jeunes. Sa consommation répétée installe en quelques semaines une dépendance physique et psychique très difficile à vaincre. Ce n'est pas un « produit de courage » : c'est un médicament dangereux hors du cadre médical.
\end{exemplebox}

\courssub{Comprendre la dépendance : une maladie du circuit synaptique}

\begin{attention}[Tolérance et sevrage : le piège biologique]
La dépendance s'installe parce que le cerveau \textbf{s'adapte} à la présence répétée de la substance : c'est la \cle{tolérance}. Par exemple, face à un excès chronique de dopamine, les neurones postsynaptiques réduisent le nombre ou la sensibilité de leurs récepteurs. Conséquences :
\begin{itemize}
\item il faut des \textbf{doses croissantes} pour obtenir le même effet ;
\item sans la substance, l'activité synaptique devient \textbf{insuffisante} : l'arrêt brutal provoque le \cle{sevrage} (angoisse, douleurs, tremblements, insomnie), qui pousse à reprendre la drogue.
\end{itemize}
La dépendance est donc une \textbf{maladie du fonctionnement synaptique}, avec des modifications mesurables des récepteurs et des médiateurs — et non un simple « manque de volonté ». C'est pourquoi elle se soigne (accompagnement médical et psychologique) et ne se juge pas moralement.
\end{attention}

\begin{savaistu}[La dopamine, « molécule du plaisir »]
La dopamine est surnommée la \emph{molécule du plaisir} (plus exactement : de la \textbf{motivation} et de la récompense). Des expériences de mesure par microdialyse chez l'animal ont montré que la plupart des drogues d'abus (cocaïne, amphétamines, nicotine, alcool, opioïdes) augmentent fortement — parfois de 2 à 10 fois — la concentration de dopamine dans les synapses du noyau accumbens, bien au-delà de ce que produit un plaisir naturel. Le circuit de la récompense, conçu par l'évolution pour nous guider vers les comportements vitaux, se trouve ainsi \textbf{détourné} au profit de la substance.
\end{savaistu}

\begin{popfigure}
\begin{center}
\begin{tikzpicture}[
  region/.style={draw=popdark, thick, rounded corners=6pt, fill=popblueL, minimum width=3.4cm, minimum height=1.2cm, align=center, font=\small\bfseries},
  synapse/.style={draw=poporange, thick, fill=poporangeL, rounded corners=3pt, align=center, font=\footnotesize, minimum width=2.6cm},
  note/.style={font=\footnotesize, align=center, text=popmuted},
  flux/.style={-{Stealth[length=3mm]}, very thick, popblue}
]
\node[region, align=center] (atv) at (0,0){Aire tegmentale\\ventrale (ATV)};
\node[region, fill=poppinkL, align=center] (na) at (7.5,0){Noyau accumbens\\(plaisir, motivation)};
\node[region, fill=popgreenL, align=center] (cpf) at (7.5,3.2){Cortex préfrontal\\(décision, contrôle)};
\node[synapse, align=center] (syn) at (3.75,-1.6){Synapses\\dopaminergiques};

\draw[flux] (atv) -- node[above, font=\footnotesize\itshape]{axones dopaminergiques} (na);
\draw[flux, popgreen] (cpf) -- node[right, font=\footnotesize\itshape]{contrôle} (na);
\draw[flux, poporange, dashed] (syn.north) -- ++(0,0.5) -| node[pos=0.25, above, font=\footnotesize\itshape, text=poporange]{dopamine} (na.south);

\node[note, align=center] at (0,-2.6){Plaisir naturel : libération\\modérée de dopamine};
\node[note, text=poporange, align=center] at (7.5,-2.6){Drogues : dopamine massive\\et anormale $\rightarrow$ dépendance};
\draw[-{Stealth}, popmuted] (0,-1.6) -- (syn.west);
\draw[-{Stealth}, poporange] (syn.east) -- (7.5,-1.6) -- (7.5,-2.2);
\end{tikzpicture}
\end{center}
\legende{Schéma simplifié du circuit de la récompense (complément, hors strict programme). Les neurones de l'aire tegmentale ventrale (ATV) projettent leurs axones vers le noyau accumbens, où ils libèrent de la dopamine dans les synapses. Le cortex préfrontal module ce circuit. Les drogues d'abus provoquent un afflux anormal de dopamine, ce qui « reprogramme » le circuit et installe la dépendance.}
\end{popfigure}

\courssub{Carte mentale de synthèse de la leçon}

\begin{popfigure}
\begin{center}
\begin{tikzpicture}[
  root/.style={draw=popdark, very thick, fill=popblue, text=white, rounded corners=6pt, align=center, font=\bfseries, minimum width=3cm},
  branche/.style={draw=popdark, thick, rounded corners=5pt, align=center, font=\small\bfseries, minimum width=3.2cm},
  feuille/.style={align=center, font=\footnotesize, text width=3.6cm},
  lien/.style={-{Stealth[length=2.5mm]}, thick, popmuted}
]
\node[root] (s) at (0,0) {LA SYNAPSE};

\node[branche, fill=popblueL] (st) at (-5.5,2.6) {Structure};
\node[feuille] at (-5.5,4.6) {élément présynaptique (vésicules), fente synaptique, membrane postsynaptique (récepteurs)};
\draw[lien] (s) -- (st);

\node[branche, fill=popgreenL] (tr) at (5.5,2.6) {Transmission};
\node[feuille] at (5.5,4.8) {chimique (neuromédiateur, exocytose, PPSE/PSI) ou électrique (gap junctions) ; expérience de Loewi};
\draw[lien] (s) -- (tr);

\node[branche, fill=poppurpleL] (so) at (-5.5,-2.6) {Sommation};
\node[feuille] at (-5.5,-4.8) {spatiale, temporelle, spatio-temporelle : le neurone intègre les PPSE et PSI et « décide » de décharger};
\draw[lien] (s) -- (so);

\node[branche, fill=poporangeL] (su) at (5.5,-2.6) {Substances};
\node[feuille] at (5.5,-5.2) {curare : bloque les récepteurs ; toxine botulique : bloque la libération ; insecticides : bloquent la dégradation ; cocaïne : bloque la recapture ; opioïdes : dépression respiratoire};
\draw[lien] (s) -- (su);
\end{tikzpicture}
\end{center}
\legende{Carte mentale de synthèse : quatre idées maîtresses organisent toute la leçon — la structure de la synapse, la nature de la transmission, la sommation des messages et l'action des substances sur chaque étape.}
\end{popfigure}

\courssub{Sensibilisation et conduite à tenir}

Les drogues ne créent rien de nouveau dans le cerveau : elles \textbf{détournent des mécanismes naturels} de communication nerveuse. C'est précisément ce qui les rend si dangereuses : elles s'insinuent dans le fonctionnement le plus intime de notre système nerveux et le modifient durablement. La conduite à tenir, pour toi et ton entourage, découle directement de la science :

\begin{itemize}
\item \textbf{Refuser} toute substance psychoactive non prescrite (tramadol de rue, cannabis, cocaïne, alcool en excès), même « pour essayer » : la dépendance peut s'installer dès les premières consommations ;
\item \textbf{Respecter strictement les prescriptions médicales} : un médicament psychoactif (antalgique, anxiolytique) n'est sûr qu'à la dose, à la durée et pour l'indication prévues par le médecin ;
\item \textbf{Se protéger lors de l'usage des pesticides} (inhibiteurs de cholinestérases) : gants, masque, vêtements couvrants, lavage soigneux après épandage, stockage hors de portée des enfants ;
\item \textbf{Informer son entourage} : expliquer scientifiquement les mécanismes (comme tu viens de l'apprendre) est bien plus convaincant qu'un simple interdit ;
\item \textbf{Orienter vers les structures de santé} toute personne dépendante : la dépendance est une maladie qui se soigne, pas une faute morale.
\end{itemize}

\begin{methode}[Réussir une question de sensibilisation au Bac]
Une question du type « À l'aide de vos connaissances, sensibilisez un jeune sur les dangers du tramadol » exige un \textbf{argumentaire scientifique}, pas un discours moral. Structure ta réponse en trois temps :
\begin{enumerate}
\item \textbf{Rappeler le fonctionnement normal} : la synapse, le neuromédiateur, son rôle ;
\item \textbf{Expliquer le mécanisme de la substance} : quelle étape est perturbée, et quelle en est la conséquence sur l'organisme (dépression respiratoire, dépendance par le circuit dopaminergique\ldots) ;
\item \textbf{Conclure par la conduite à tenir} : refus, respect des prescriptions, recours au personnel de santé.
\end{enumerate}
Un exposé qui ne dit que « la drogue c'est mal » sans aucun mécanisme synaptique ne rapporte quasiment aucun point.
\end{methode}

\courssub{Trace écrite et exercice de synthèse}

\begin{aretenir}[Trace écrite — résumé de la leçon]
La synapse est la zone de contact fonctionnel entre un neurone et une cellule cible (neurone, muscle, glande). La transmission y est \textbf{chimique} (libération d'un neuromédiateur par exocytose, fixation sur des récepteurs postsynaptiques, puis recapture ou dégradation) ou, plus rarement, \textbf{électrique} (passage direct du courant par des jonctions communicantes). Le neurone postsynaptique \textbf{intègre} les messages excitateurs (PPSE) et inhibiteurs (PSI) par sommation spatiale et temporelle avant de décharger ou non un potentiel d'action. De nombreuses substances (drogues, médicaments, toxines, pesticides) perturbent une étape de la transmission — libération, fixation, recapture, dégradation — et provoquent paralysies, tétanies, dépressions respiratoires ou dépendances. La dépendance est une maladie du circuit synaptique de la récompense (dopamine) : elle impose prévention, respect des prescriptions et prise en charge médicale.
\end{aretenir}

\begin{center}
\begin{tabularx}{\linewidth}{|l|X|X|X|}
\hline
\rowcolor{poporangeL} \textbf{Substance} & \textbf{Cible synaptique} & \textbf{Effet sur la transmission} & \textbf{Trouble observé} \\
\hline
Curare & Récepteurs nicotiniques (plaque motrice) & Blocage de la fixation de l'ACh & Paralysie flasque \\
\hline
Toxine botulique & Terminaison présynaptique & Blocage de l'exocytose de l'ACh & Paralysie flasque \\
\hline
Toxine tétanique & Interneurones inhibiteurs (moelle) & Blocage de la libération de glycine/GABA & Contractures (tétanie) \\
\hline
Insecticides (organophosphorés) & Acétylcholinestérase & Accumulation de l'ACh dans la fente & Tétanie, asphyxie \\
\hline
Cocaïne & Pompe de recapture & Accumulation de dopamine & Excitation, dépendance \\
\hline
Nicotine & Récepteurs nicotiniques cérébraux & Activation des neurones à dopamine & Dépendance \\
\hline
Tramadol, opioïdes & Récepteurs opioïdes (tronc cérébral) & Dépression des centres respiratoires & Coma, arrêt respiratoire \\
\hline
Alcool (forte dose) & Récepteurs GABA & Excès d'inhibition & Ivresse, dépression respiratoire \\
\hline
\end{tabularx}
\end{center}

\begin{exoresolu}[Sensibilisation argumentée (type Bac)]
\textbf{Énoncé.} Ton jeune frère, élève en classe de Première, te dit que « prendre du tramadol avant un match de football, ce n'est pas grave, ça donne de la force ». En t'appuyant sur tes connaissances sur la transmission synaptique, rédige un texte de sensibilisation de 10 à 15 lignes pour le convaincre d'y renoncer.
\tcblower
\textbf{Solution détaillée.} Le tramadol est un antalgique opioïde réservé, sur ordonnance, aux douleurs intenses. Il agit au niveau des synapses du système nerveux central en se fixant sur les récepteurs opioïdes, notamment ceux des centres respiratoires du tronc cérébral. À dose excessive, il ralentit l'activité de ces centres : la respiration devient superficielle puis peut s'arrêter — c'est la dépression respiratoire, cause de comas et de décès chez les jeunes qui en abusent. Par ailleurs, le tramadol augmente l'activité du circuit de la récompense, dont le neuromédiateur est la dopamine : le cerveau s'adapte à sa présence (tolérance), si bien qu'il faut des doses croissantes pour le même effet, et que l'arrêt provoque un sevrage douloureux. C'est ainsi que s'installe la dépendance, une véritable maladie du fonctionnement synaptique. Loin de « donner de la force », le tramadol masque la douleur — signal d'alarme utile — et expose à des blessures aggravées pendant l'effort. La conduite à tenir est claire : ne jamais consommer de médicament sans prescription, refuser les produits proposés par l'entourage, et signaler tout cas de dépendance à un personnel de santé. Ta santé et ta vie valent infiniment plus qu'un match.
\end{exoresolu}

\begin{exercice}[Pour t'entraîner]
\begin{enumerate}
\item Complète le tableau de synthèse des substances en ajoutant, pour chacune, l'étape de la transmission perturbée (libération, fixation, recapture ou dégradation).
\item Explique pourquoi l'injection d'acétylcholinestérase ne sauverait pas un patient intoxiqué au curare, alors qu'un inhibiteur de cette enzyme peut soulager certains patients atteints de myasthénie (maladie où les récepteurs nicotiniques sont détruits).
\item Rédige un court message de sensibilisation (5 lignes) sur les précautions à prendre lors de l'épandage d'insecticides dans une plantation de cacao.
\end{enumerate}
\end{exercice}

\begin{corrige}[Exercice de synthèse]
\begin{enumerate}
\item Curare : fixation ; toxine botulique : libération ; toxine tétanique : libération (de médiateurs inhibiteurs) ; insecticides : dégradation ; cocaïne : recapture ; nicotine : fixation (agoniste) ; opioïdes et alcool : fixation sur récepteurs spécifiques avec excès d'inhibition.
\item Dans l'intoxication au curare, les récepteurs sont \emph{occupés} par le curare : même si l'acétylcholine s'accumulait, elle ne pourrait pas se fixer davantage ; ajouter de l'enzyme ne ferait qu'aggraver le déficit. Dans la myasthénie, les récepteurs sont \emph{moins nombreux} mais fonctionnels : inhiber l'acétylcholinestérase prolonge la présence de l'ACh dans la fente, ce qui augmente la probabilité de fixation sur les récepteurs restants et améliore la contraction musculaire.
\item Exemple attendu : « Les insecticides organophosphorés bloquent l'enzyme qui détruit l'acétylcholine dans nos synapses : le médiateur s'accumule et provoque des contractions incontrôlées pouvant arrêter la respiration. Portez gants, masque et vêtements couvrants, ne mangez ni ne buvez pendant l'épandage, lavez-vous soigneusement après, et conservez les produits hors de portée des enfants. »
\end{enumerate}
\end{corrige}

Tu disposes désormais d'une vision complète de la communication synaptique, de ses perturbations et de leurs conséquences. Retiens le fil conducteur de toute cette leçon : \textbf{comprendre le mécanisme, c'est déjà se protéger} — et c'est aussi la marque d'un excellent candidat au Baccalauréat.

\newpage

\coursec{Activité d'intégration}

\coursec{Sensibilisation sur les effets de certaines substances sur la transmission synaptique}

\courssub{Objectifs de l'activité}

À l'issue de cette activité d'intégration, tu seras capable de :

\begin{itemize}
\item identifier et légender une synapse chimique sur une électronographie (savoir-faire : schématiser l'ultrastructure) ;
\item interpréter l'expérience historique de Loewi et en déduire la nature chimique de la transmission synaptique ;
\item distinguer, sur des enregistrements de potentiels postsynaptiques, une sommation temporelle d'une sommation spatiale ;
\item localiser la cible synaptique d'une substance (drogue, médicament, toxine) et expliquer les symptômes d'une intoxication réelle ;
\item produire un argumentaire de sensibilisation fondé exclusivement sur des preuves scientifiques, à destination de jeunes élèves.
\end{itemize}

Cette activité mobilise l'ensemble des savoirs de la leçon dans une \cle{famille de situations} : les perturbations de la transmission synaptique.

\courssub{Situation}

\textbf{Enquête scientifique : pourquoi le tramadol et les pesticides intoxiquent-ils les jeunes Camerounais ?}

\medskip

Au mois de mars, le service des urgences de l'Hôpital de District de Bafoussam reçoit, à quelques heures d'intervalle, deux patients présentant des troubles neurologiques :

\begin{itemize}
\item \textbf{Cas n°1.} Rodrigue, élève de Première âgé de 17 ans, est amené inconscient par ses camarades. Il a consommé, lors d'une fête, des comprimés de \cle{tramadol} (un antalgique opioïde détourné de son usage médical) achetés au marché. À l'examen : respiration très ralentie (6 mouvements par minute au lieu de 14 à 18), pupilles contractées (myosis), somnolence profonde.
\item \textbf{Cas n°2.} M. Kamdem, planteur de 45 ans de la région de Penja, a pulvérisé sans équipement de protection un insecticide \cle{organophosphoré} sur sa plantation de bananiers. Il présente des sueurs abondantes, des contractions musculaires involontaires (fasciculations), des crampes et une salivation excessive (hypersalivation).
\end{itemize}

Le médecin de garde, souhaitant sensibiliser les jeunes, confie à ta classe de Terminale D un dossier scientifique de quatre documents. Ta mission : expliquer rigoureusement les symptômes des deux patients et rédiger un message de prévention.

\medskip

\textbf{Document 1 — Électronographie d'une synapse.} Une micrographie électronique (grossissement $\times \num{50000}$) montre deux éléments cellulaires séparés par un espace étroit d'environ \SI{20}{nm}. L'élément supérieur est un renflement contenant de nombreuses vésicules sphériques de \SI{50}{nm} de diamètre et des mitochondries ; la membrane de l'élément inférieur est épaissie en regard de cet espace.

\begin{popfigure}
\begin{center}
\begin{tikzpicture}[scale=0.9, every node/.style={font=\small}]
% Presynaptic terminal
\draw[popdark, thick, fill=popblueL!30] (0,0) .. controls (0,1.5) and (3,1.5) .. (3,0) .. controls (3,-1.5) and (0,-1.5) .. cycle;
% Mitochondria
\draw[poporange, fill=poporange!20] (0.5,0.5) ellipse (0.3 and 0.15);
\draw[poporange, fill=poporange!20] (1.5,-0.2) ellipse (0.25 and 0.12);
\draw[poporange, fill=poporange!20] (2.2,0.6) ellipse (0.2 and 0.1);
% Vesicles
\foreach \x/\y in {0.8/0.2, 1.2/-0.4, 1.8/0.1, 2.5/-0.3, 0.6/-0.6, 2.0/0.5, 2.7/0.0}
    \draw[popblue, thick, fill=popblue!30] (\x,\y) circle (0.07);
% Synaptic cleft
\draw[popmuted, thick] (0,-1.5) -- (3,-1.5);
\draw[popmuted, thick] (0,-1.7) -- (3,-1.7);
\draw[popmuted, thin, densely dashed] (0,-1.5) -- (0,-1.7);
\draw[popmuted, thin, densely dashed] (3,-1.5) -- (3,-1.7);
\node[popmuted, align=center] at (1.5,-1.6) {Fente synaptique $\approx$ 20 nm};
% Postsynaptic membrane
\draw[popdark, thick, fill=popgreenL!30] (0,-1.7) .. controls (0,-2.2) and (3,-2.2) .. (3,-1.7) .. controls (3,-1.9) and (0,-1.9) .. cycle;
% Receptors
\foreach \x in {0.5, 1.2, 1.9, 2.6}
    \draw[popgreen, thick] (\x,-1.7) -- (\x,-2.0) node[midway, right=1pt] {\tiny Récepteur};
% Labels
\node[popdark, align=center] at (1.5, 1.8) {Bouton synaptique (présynaptique)};
\node[popdark, align=center] at (1.5, -2.5) {Membrane postsynaptique};
\draw[->, popdark] (0.5,0.5) -- ++(-0.8,0.5) node[left] {Mitochondrie};
\draw[->, popdark] (1.2,-0.4) -- ++(-0.8,-0.5) node[left] {Vésicule synaptique ($\approx$ 50 nm)};
\draw[->, popdark] (2.6,-1.8) -- ++(0.8,-0.5) node[right] {Récepteur aux neurotransmetteurs};
\end{tikzpicture}
\end{center}
\legende{Schéma d'interprétation de l'électronographie (Document 1) : ultrastructure d'une synapse chimique. On distingue le bouton synaptique vésiculé, la fente synaptique et la membrane postsynaptique à récepteurs.}
\end{popfigure}

\medskip

\textbf{Document 2 — L'expérience de Loewi (1921).} Deux cœurs de grenouille isolés, A et B, battent spontanément dans du liquide de Ringer. Le nerf vague (nerf X) du cœur A est stimulé électriquement : la fréquence cardiaque de A chute de 60 à 40 battements par minute. Le liquide ayant baigné le cœur A est alors recueilli et versé sur le cœur B, dont le nerf vague n'a \emph{pas} été stimulé : la fréquence de B chute également de 62 à 42 battements par minute.

\begin{popfigure}
\begin{center}
\begin{tikzpicture}[scale=0.85, every node/.style={font=\small}]
% Heart A
\draw[popdark, thick, fill=poppinkL] (0,0) ellipse (1.5 and 1);
\node[popdark] at (0,0.5) {Cœur A};
\draw[poppurple, thick, decorate, decoration={snake, amplitude=1pt, segment length=4pt}] (-2,0) -- (-0.5,0) node[midway, above] {Nerf vague};
\draw[->, poppurple, thick] (-2.5,0.3) -- (-2,0.3) node[left] {Stimulation électrique};
% Bath A
\draw[popblue, thick, fill=popblueL!20] (-2.5,-1.5) rectangle (1.5,-0.5);
\node[popblue] at (-0.5,-1.0) {Bain de Ringer A};
% Transfer arrow
\draw[->, popdark, thick] (2.5,-1.0) -- (4.5,-1.0) node[midway, above] {Transfert du liquide};
% Heart B
\draw[popdark, thick, fill=poppinkL] (7,0) ellipse (1.5 and 1);
\node[popdark] at (7,0.5) {Cœur B};
\draw[poppurple, thick] (5.5,0) -- (8.5,0); % Vagus not stimulated
\node[poppurple] at (7,-0.5) {Nerf vague (non stimulé)};
% Bath B
\draw[popblue, thick, fill=popblueL!20] (4.5,-1.5) rectangle (8.5,-0.5);
\node[popblue] at (6.5,-1.0) {Bain de Ringer B + Liquide A};
% Frequency labels
\node[popdark, align=center] at (0,-2.5) {FC initiale : 60 bpm\\FC finale : 40 bpm};
\node[popdark, align=center] at (7,-2.5) {FC initiale : 62 bpm\\FC finale : 42 bpm};
\end{tikzpicture}
\end{center}
\legende{Schéma principe de l'expérience de Loewi (Document 2) : la substance libérée par le nerf vague (Vagusstoff = acétylcholine) diffuse dans le bain et ralentit le cœur B sans innervation directe.}
\end{popfigure}

\medskip

\textbf{Document 3 — Enregistrements de potentiels postsynaptiques.} Un neurone postsynaptique reçoit deux afférences excitatrices $a$ et $b$. Son potentiel de membrane de repos vaut $\SI{-70}{mV}$ et son seuil d'excitation $\SI{-55}{mV}$.

\begin{itemize}
\item \textbf{Expérience 1 (Sommation temporelle) :} une stimulation unique de $a$ provoque un PPSE (potentiel postsynaptique excitateur) de $+\SI{5}{mV}$, insuffisant. Trois stimulations successives de $a$, à \SI{2}{ms} d'intervalle, provoquent des dépolarisations successives de $+5$, $+10$ puis $+\SI{15}{mV}$ : le seuil est atteint et un potentiel d'action (PA) naît.
\item \textbf{Expérience 2 (Sommation spatiale) :} une stimulation unique de $a$ ($+\SI{5}{mV}$) ou de $b$ ($+\SI{5}{mV}$) reste sous-liminaire ; la stimulation \emph{simultanée} de $a$ et $b$ provoque une dépolarisation de $+\SI{10}{mV}$... qui reste sous le seuil ; mais avec quatre afférences identiques stimulées ensemble ($a$, $b$, $c$, $d$ : $4 \times \SI{5}{mV} = +\SI{20}{mV}$), le seuil est franchi et un PA apparaît.
\end{itemize}

\begin{popfigure}
\begin{center}
\begin{tikzpicture}[scale=0.7, every node/.style={font=\small}]
% Axes style
\def\axisstyle{->, >=Stealth, popdark}
\def\gridstyle{popline!30}
% --- Temporal Summation ---
\begin{scope}[xshift=0cm]
\draw[->, >=Stealth, popdark] (0,0) -- (7,0) node[right] {$t$ (ms)};
\draw[->, >=Stealth, popdark] (0,0) -- (0,4.5) node[above] {$V_m$ (mV)};
% Grid lines
\foreach \y/\lab in {0/-70, 1/-60, 2/-50, 3/-40, 4/-30} {
    \draw[\gridstyle] (0,\y) -- (7,\y);
    \node[left] at (0,\y) {\tiny $\lab$};
}
\foreach \x in {1,2,3,4,5,6} \draw[\gridstyle] (\x,0) -- (\x,4.5);
% Threshold line
\draw[poporange, dashed, thick] (0,1.5) -- (7,1.5) node[right, poporange] {\tiny Seuil (-55 mV)};
% Resting potential
\draw[popblue, dashed] (0,0) -- (7,0) node[right, popblue] {\tiny Repos (-70 mV)};
% Presynaptic spikes (schematic)
\foreach \t in {0.5, 2.5, 4.5} {
    \draw[poppurple, thick] (\t, -0.3) -- (\t, -0.1);
    \draw[poppurple, thick] (\t, -0.2) -- (\t+0.1, -0.2);
}
\node[poppurple] at (3.5, -0.6) {Afférence $a$ (3 spikes)};
% EPSPs summation
\draw[popcurve, thick, domain=0.5:7, samples=100, smooth] plot (\x, {1.5*(1-exp(-(\x-0.5)/0.8))});
\draw[popcurve, thick, domain=2.5:7, samples=100, smooth] plot (\x, {1.5*(1-exp(-(\x-2.5)/0.8)) + 1.5*exp(-(\x-2.5)/0.8)});
\draw[popcurve, thick, domain=4.5:7, samples=100, smooth] plot (\x, {1.5*(1-exp(-(\x-4.5)/0.8)) + 1.5*exp(-(\x-2.5)/0.8) + 1.5*exp(-(\x-0.5)/0.8)});
% Simplified visual summation for clarity
\draw[popcurve, very thick, domain=0.5:7, samples=200, smooth, popred] plot (\x, {
    ( (\x>=0.5) ? 1.5*(1-exp(-(\x-0.5)/0.8))*exp(-(\x-0.5)/1.5) : 0 ) +
    ( (\x>=2.5) ? 1.5*(1-exp(-(\x-2.5)/0.8))*exp(-(\x-2.5)/1.5) : 0 ) +
    ( (\x>=4.5) ? 1.5*(1-exp(-(\x-4.5)/0.8))*exp(-(\x-4.5)/1.5) : 0 )
});
\node[popred, align=center] at (3.5, 3.8) {Sommation\\temporelle};
\node[popred] at (5.5, 2.8) {PA !};
\end{scope}

% --- Spatial Summation ---
\begin{scope}[xshift=9cm]
\draw[->, >=Stealth, popdark] (0,0) -- (7,0) node[right] {$t$ (ms)};
\draw[->, >=Stealth, popdark] (0,0) -- (0,4.5) node[above] {$V_m$ (mV)};
\foreach \y/\lab in {0/-70, 1/-60, 2/-50, 3/-40, 4/-30} {
    \draw[\gridstyle] (0,\y) -- (7,\y);
    \node[left] at (0,\y) {\tiny $\lab$};
}
\foreach \x in {1,2,3,4,5,6} \draw[\gridstyle] (\x,0) -- (\x,4.5);
\draw[poporange, dashed, thick] (0,1.5) -- (7,1.5) node[right, poporange] {\tiny Seuil (-55 mV)};
\draw[popblue, dashed] (0,0) -- (7,0) node[right, popblue] {\tiny Repos (-70 mV)};
% 4 simultaneous spikes
\foreach \t in {2} {
    \draw[poppurple, thick] (\t, -0.3) -- (\t, -0.1);
    \draw[poppurple, thick] (\t, -0.2) -- (\t+0.1, -0.2);
    \draw[poppurple, thick] (\t+0.3, -0.3) -- (\t+0.3, -0.1);
    \draw[poppurple, thick] (\t+0.3, -0.2) -- (\t+0.4, -0.2);
    \draw[poppurple, thick] (\t+0.6, -0.3) -- (\t+0.6, -0.1);
    \draw[poppurple, thick] (\t+0.6, -0.2) -- (\t+0.7, -0.2);
    \draw[poppurple, thick] (\t+0.9, -0.3) -- (\t+0.9, -0.1);
    \draw[poppurple, thick] (\t+0.9, -0.2) -- (\t+1.0, -0.2);
}
\node[poppurple, align=center] at (3.5, -0.6) {Afférences $a,b,c,d$\\(simultanées)};
% Large EPSP
\draw[popcurve, very thick, domain=2:7, samples=200, smooth, popred] plot (\x, { 3.0*(1-exp(-(\x-2)/0.5))*exp(-(\x-2)/2.0) });
\node[popred, align=center] at (3.5, 3.8) {Sommation\\spatiale};
\node[popred] at (4.5, 3.2) {PA !};
\end{scope}
\end{tikzpicture}
\end{center}
\legende{Modélisation des sommations (Document 3). \textbf{Gauche :} Sommation temporelle (3 spikes sur une même afférence $a$ à 2 ms d'intervalle). \textbf{Droite :} Sommation spatiale (4 afférences $a,b,c,d$ stimulées simultanément). Le seuil (-55 mV) n'est franchi que par la somme algébrique des PPSE.}
\end{popfigure}

\medskip

\textbf{Document 4 — Cibles synaptiques des substances.}

\begin{itemize}
\item Le tramadol active les récepteurs opioïdes des synapses des centres respiratoires du bulbe rachidien : il \emph{inhibe} la libération des médiateurs excitateurs (effet présynaptique) et hyperpolarise les neurones postsynaptiques en ouvrant des canaux $\text{K}^+$ (effet postsynaptique, PPSI). La courbe ci-dessous montre la fréquence respiratoire en fonction de la dose.
\item Les organophosphorés \emph{inhibent l'acétylcholinestérase} (AChE), l'enzyme qui hydrolyse normalement l'acétylcholine (ACh) dans la fente synaptique en moins de \SI{1}{ms}. La concentration synaptique d'ACh passe ainsi de \num{0,1} à \SI{0,8}{mmol/L} (valeurs indicatives) : la stimulation des récepteurs cholinergiques (nicotiniques et muscariniques) devient excessive et prolongée.
\end{itemize}

\begin{popfigure}
\begin{center}
\begin{tikzpicture}
\begin{axis}[width=0.75\linewidth,height=5.5cm,
xlabel={Dose de tramadol (mg)},ylabel={Fréquence respiratoire (mouvements/min)},
xmin=0,xmax=400,ymin=0,ymax=20,
grid=both,grid style={popline!40},
tick label style={font=\small}, label style={font=\small},
legend style={at={(0.98,0.98)}, anchor=north east, font=\small, fill=white, fill opacity=0.8}]
\addplot[popcurve, domain=0:400, samples=200] {16*exp(-x/120)+1.5};
\addplot[mark=*, poporange, only marks] coordinates {(0,16) (100,8.7) (200,4.5) (300,3.1)};
\addlegendentry{Modèle exponentiel}
\addlegendentry{Données exp.}
\draw[popred, dashed, thick] (200,0) -- (200,20) node[above, pos=0.5, rotate=90, popred] {\tiny Seuil critique $\approx$ 200 mg};
\end{axis}
\end{tikzpicture}
\end{center}
\legende{Effet de la dose de tramadol sur la fréquence respiratoire : l'inhibition synaptique des centres bulbaires entraîne une dépression respiratoire dose-dépendante, potentiellement mortelle (arrêt respiratoire) au-delà de \SI{200}{mg}.}
\end{popfigure}

\courssub{Consignes}

Par groupes de quatre, tu répondras aux questions suivantes, en rédigeant des réponses argumentées :

\begin{enumerate}
\item \textbf{Document 1.} Identifie la structure observée, nomme et légende ses quatre éléments essentiels. Précise le type anatomique de synapse que pourrait représenter cette électronographie.
\item \textbf{Document 2.} Formule l'hypothèse que Loewi cherchait à tester. Interprète les deux étapes de l'expérience et conclus sur la nature de la transmission synaptique. Quel nom donne-t-on aujourd'hui à la substance libérée par le nerf vague ?
\item \textbf{Document 3.} Nomme le type de sommation réalisé dans l'expérience 1 et dans l'expérience 2, en justifiant par les données chiffrées. Quelle propriété fondamentale du neurone postsynaptique ces expériences illustrent-elles ?
\item \textbf{Document 4.} Pour chaque patient, localise précisément la cible synaptique de la substance (pré- ou postsynaptique, enzyme, récepteur) et explique la chaîne causale reliant le mécanisme moléculaire aux symptômes observés.
\item \textbf{Synthèse.} Rédige un message de sensibilisation (15 à 20 lignes, ou le plan d'une affiche) destiné aux élèves de Seconde, en t'appuyant \emph{uniquement} sur des arguments scientifiques issus des documents et de ton cours.
\end{enumerate}

\courssub{Ressources à mobiliser}

\begin{aretenir}[Boîte à outils de la leçon]
\begin{itemize}
\item \textbf{Ultrastructure d'une synapse chimique :} bouton synaptique (vésicules de neurotransmetteur $\approx \SI{50}{nm}$, mitochondries), membrane présynaptique (canaux $\text{Ca}^{2+}$ voltage-dépendants), fente synaptique ($\approx \SI{20}{nm}$), membrane postsynaptique porteuse de récepteurs spécifiques (canaux ioniques ligand-dépendants).
\item \textbf{Types anatomiques :} synapses neuro-neuroniques (SNC, SNA), neuro-musculaires (plaque motrice, toujours excitatrice, ACh/récepteur nicotinique), neuro-glandulaires.
\item \textbf{Types physiologiques :}
      \begin{itemize}
      \item \cle{Transmission chimique} (majoritaire) : délai synaptique ($\approx \SI{0,5}{ms}$), unidirectionnelle, modulable (médiateur, récepteur, enzyme). Démontrée par Loewi (ACh = Vagusstoff).
      \item \cle{Transmission électrique} (jonctions communicantes / \emph{gap junctions}) : continuité cytoplasmique par connexons (connexines), bidirectionnelle, quasi instantanée, pas de médiateur. Synchronisation (cœur, réflexes d'échappement).
      \end{itemize}
\item \textbf{Expérience de Loewi (1921) :} démontre la nature \cle{chimique} de la transmission (libération d'un médiateur diffusible).
\item \textbf{Intégration neuronale :} sommation algébrique des PPSE (dépolarisants, ex. glutamate, ACh) et PPSI (hyperpolarisants, ex. GABA, glycine).
      \begin{itemize}
      \item \cle{Sommation temporelle} : même afférence, stimulations rapprochées ($<\tau$ membranaire).
      \item \cle{Sommation spatiale} : afférences différentes, stimulations simultanées.
      \item \cle{Sommation spatio-temporelle} : combinaison des deux.
      \end{itemize}
      Rôle intégrateur : le neurone est un \emph{coïncidence detector} ; PA déclenché si $V_m \ge V_{\text{seuil}}$ ($\approx \SI{-55}{mV}$).
\item \textbf{Modes d'action des substances exogènes sur la synapse chimique :}
      \begin{itemize}
      \item \textbf{Agonistes / Mimétiques :} se fixent sur le récepteur et l'activent (ex. nicotine sur récepteur nicotinique, tramadol/opioïdes sur récepteur $\mu$).
      \item \textbf{Antagonistes / Bloquants :} se fixent sans activer, empêchent la fixation du médiateur endogène (ex. curare sur récepteur nicotinique, naloxone sur récepteur opioïde).
      \item \textbf{Action sur la libération :} blocage (toxine botulique, tétanique) ou facilitation (venin araignée noire, amphétamines).
      \item \textbf{Action sur la dégradation / recapture :} inhibition de l'enzyme (organophosphorés sur AChE, physostigmine) ou du transporteur de recapture (cocaïne/antidépresseurs ISRS sur SERT/DAT).
      \end{itemize}
\end{itemize}
\end{aretenir}

\courssub{Démarche et corrigé commenté}

\begin{exoresolu}[Corrigé commenté de l'enquête]
\textbf{Énoncé de l'activité.} À partir de l'étude de deux cas cliniques réels (intoxication au tramadol et aux organophosphorés) et de quatre documents scientifiques (ultrastructure, expérience historique, électrophysiologie, pharmacologie), identifier les cibles synaptiques, expliquer les mécanismes physiopathologiques et produire un message de prévention fondé sur la preuve scientifique.
\tcblower

\textbf{Question 1 — Identification de la synapse (Document 1).}

\emph{Démarche :} on repère les indices structuraux clés : renflement vésiculé, espace intercellulaire étroit, spécialisation membraneuse postsynaptique.

L'électronographie montre une \cle{synapse chimique}. Sa légende comporte quatre éléments essentiels :
\begin{enumerate}
\item le \textbf{bouton synaptique} (élément supérieur renflé, terme axonal), contenant les \textbf{vésicules synaptiques} ($\approx \SI{50}{nm}$ de diamètre) remplies de neurotransmetteur et des \textbf{mitochondries} (fournissent l'ATP pour le recyclage vésiculaire et la pompe $\text{Ca}^{2+}$) ;
\item la \textbf{membrane présynaptique}, riche en canaux $\text{Ca}^{2+}$ voltage-dépendants, siège de l'exocytose déclenchée par l'arrivée du potentiel d'action ;
\item la \textbf{fente synaptique} ($\approx \SI{20}{nm}$), espace extracellulaire de diffusion du médiateur chimique ;
\item la \textbf{membrane postsynaptique}, épaissie (densité postsynaptique) car riche en \textbf{récepteurs} spécifiques (canaux ioniques ligand-dépendants) et en protéines d'échafaudage.
\end{enumerate}
Selon la nature de l'élément postsynaptique (non visible ici), il peut s'agir d'une synapse neuro-neuronique (interneurone), neuro-musculaire (plaque motrice) ou neuro-glandulaire. L'aspect très épaissi de la membrane réceptrice est compatible avec une plaque motrice ou une synapse neuro-neuronique à forte densité réceptrice.

\medskip

\textbf{Question 2 — Interprétation de l'expérience de Loewi.}

\emph{Hypothèse testée :} la transmission de l'influx nerveux à la synapse fait intervenir une substance chimique diffusible (hypothèse chimique) et non un courant électrique direct (hypothèse électrique).

\emph{Étape 1 (Témoin de l'effet nerveux) :} la stimulation électrique du nerf vague du cœur A ralentit sa fréquence (60 $\to$ 40 bpm). Cela confirme l'effet inhibiteur physiologique du nerf vague sur le nœud sinusal.

\emph{Étape 2 (Test décisif du milieu) :} le liquide de perfusion du cœur A, prélevé \emph{après} stimulation, est appliqué sur le cœur B \emph{sans stimulation de son nerf vague}. Le cœur B ralentit aussi (62 $\to$ 42 bpm).

\emph{Interprétation :} aucune connexion nerveuse n'existe entre A et B. Seul le liquide a été transféré. La stimulation du nerf vague a donc provoqué la libération, dans le bain, d'une \textbf{substance chimique soluble} (le « Vagusstoff ») capable de reproduire l'effet du nerf à distance.

\emph{Conclusion :} la transmission synaptique (ici neuro-cardiaque) est de nature \cle{chimique}. La substance identifiée ultérieurement par Dale et Loewi est l'\textbf{acétylcholine} (ACh).

\medskip

\textbf{Question 3 — Types de sommation et rôle intégrateur.}

\emph{Expérience 1 :} la même afférence $a$ est stimulée trois fois à \SI{2}{ms} d'intervalle. Les PPSE unitaires ($+\SI{5}{mV}$) s'additionnent dans le temps car l'intervalle est inférieur à la constante de temps membranaire ($\tau_m \approx \SI{10}{-20}{ms}$) : dépolarisation cumulée $+5 \to +10 \to +\SI{15}{mV}$. Le seuil ($\SI{-55}{mV}$) est franchi $\to$ \cle{sommation temporelle}.

\emph{Expérience 2 :} des afférences \emph{différentes} ($a$, $b$, $c$, $d$) sont stimulées \emph{simultanément}. Leurs PPSE unitaires ($+\SI{5}{mV}$ chacun) s'additionnent dans l'espace (sur la membrane postsynaptique) : $4 \times \SI{5}{mV} = +\SI{20}{mV}$. Le seuil est franchi $\to$ \cle{sommation spatiale}.

\emph{Propriété fondamentale illustrée :} le \textbf{rôle intégrateur du neurone postsynaptique}. Il effectue une \emph{somme algébrique} (intégration) des potentiels postsynaptiques excitateurs (PPSE) et inhibiteurs (PPSI) reçus sur l'ensemble de son arbre dendritique et de son corps cellulaire. Un potentiel d'action n'est déclenché au cône d'émergence (zone de déclenchement, seuil bas) que si la dépolarisation nette atteint le seuil d'excitation. C'est la base du codage neuronal (fréquence des PA $\propto$ intensité de la somme des entrées).

\medskip

\textbf{Question 4 — Mécanismes d'intoxication : de la cible moléculaire aux symptômes.}

\emph{Cas n°1 — Tramadol (Rodrigue, 17 ans).}
\begin{itemize}
\item \textbf{Cible synaptique :} \textbf{Double}. (1) \textbf{Présynaptique} : activation des récepteurs opioïdes $\mu$ (Gi/o) sur les terminaisons excitatrices (glutamatergiques) des centres respiratoires (bulbe rachidien) $\to$ inhibition de l'entrée $\text{Ca}^{2+}$ $\to$ \textbf{diminution de la libération de glutamate}. (2) \textbf{Postsynaptique} : activation des récepteurs opioïdes $\mu$ sur les neurones respiratoires bulbaire $\to$ ouverture canaux $\text{K}^+$ (rectificateur entrant) $\to$ \textbf{hyperpolarisation (PPSI)}.
\item \textbf{Chaîne causale :} Renforcement de l'inhibition (PPSI) + affaiblissement de l'excitation (moins de PPSE) $\to$ impossibilité d'atteindre le seuil de déclenchement des PA dans les neurones du centre respiratoire (groupe respiratoire ventral/pré-Bötzinger) $\to$ \textbf{arrêt de la commande motrice phrénique} $\to$ \textbf{bradypnée sévère} (6 mov/min, courbe : $\approx \SI{4,5}{mov/min}$ à \SI{200}{mg}) $\to$ hypoxie, hypercapnie, coma, risque de \textbf{arrêt respiratoire mortel}. Myosis (récepteurs $\mu$ noyau d'Edinger-Westphal).
\end{itemize}

\emph{Cas n°2 — Organophosphoré (M. Kamdem, 45 ans).}
\begin{itemize}
\item \textbf{Cible synaptique :} \textbf{Enzymatique (fente synaptique)}. Inhibition \textbf{irréversible} (phosphorylation du site actif sérine) de l'\textbf{acétylcholinestérase (AChE)}.
\item \textbf{Chaîne causale :} Blocage de l'hydrolyse de l'ACh ($\text{ACh} \xrightarrow{\text{AChE}} \text{Choline} + \text{Acétate}$) $\to$ accumulation massive d'ACh dans la fente (conc. $\times 8$, de \num{0,1} à \SI{0,8}{mmol/L}) $\to$ stimulation \textbf{prolongée et excessive} des récepteurs cholinergiques :
      \begin{itemize}
      \item \textbf{Récepteurs nicotiniques (plaques motrices, ganglions) :} dépolarisation prolongée $\to$ inactivation canaux $\text{Na}^+$ $\to$ \textbf{bloc de la transmission} (paralysie flasque secondaire) mais d'abord \textbf{fasciculations, crampes} (décharges répétées non contrôlées).
      \item \textbf{Récepteurs muscariniques (effeteurs parasympathiques : glandes, muscle lisse, cœur) :} surstimulation $\to$ \textbf{syndrome cholinergique} : hypersalivation, sueurs, larmoiements, bronchorrhée, diarrhée, bradycardie, myosis (SLUDGE : Salivation, Lacrimation, Urination, Diarrhée, Gastro-intestinal, Emesis).
      \end{itemize}
\end{itemize}

\medskip

\textbf{Question 5 — Message de sensibilisation (exemple de production attendue).}

\begin{exemplebox}[Message de prévention pour élèves de Seconde]
\textbf{TRAMADOL \& PESTICIDES : VOS SYNAPSES EN DANGER — LA SCIENCE VOUS ALERTE}

Saviez-vous qu'un comprimé de tramadol « pour tenir » ou une pulvérisation sans masque peut arrêter votre respiration ou bloquer vos muscles ? Ce ne sont pas des rumeurs, c'est de la physiologie.

\textbf{1. Le tramadol éteint votre centre respiratoire.} En se fixant sur les récepteurs opioïdes de vos neurones bulbaire, il ouvre des canaux potassiques (hyperpolarisation) et bloque la libération de glutamate. Résultat : vos neurones respiratoires ne peuvent plus déclencher de potentiels d'action. À \SI{200}{mg}, la fréquence respiratoire chute sous \SI{5}{mov/min} (courbe Doc 4). \textbf{Vous arrêtez de respirer sans vous en rendre compte.}

\textbf{2. Les organophosphorés empoisonnent votre « frein » enzymatique.} Ils détruisent l'acétylcholinestérase, l'enzyme qui nettoie l'acétylcholine en \SI{<1}{ms}. L'ACh s'accumule (\num{0,1} $\to$ \SI{0,8}{mmol/L}) et surstimule vos récepteurs : muscles (crampes, fasciculations), glandes (sueurs, salivation), cœur (bradycardie). C'est le syndrome cholinergique, urgence vitale.

\textbf{La synapse n'est pas une poubelle.} Toute molécule qui imite, bloque ou prolonge un neurotransmetteur perturbe le code nerveux. \textbf{Refusez le tramadol hors ordonnance. Exigez gants, masque, lunettes pour tout pesticide.} Protégez votre système nerveux : c'est votre vie qui est en jeu.
\end{exemplebox}

\end{exoresolu}

\courssub{Bilan de l'activité}

Cette enquête t'a permis de mettre en évidence que :

\begin{itemize}
\item la \textbf{structure} de la synapse (vésicules, fente, récepteurs, enzymes de dégradation) conditionne son \textbf{fonctionnement} : toute substance qui agit sur l'un de ces éléments perturbe la communication nerveuse ;
\item la nature \cle{chimique} de la transmission, démontrée par Loewi, explique pourquoi des molécules étrangères (drogues, médicaments détournés, toxines environnementales) peuvent imiter, bloquer ou prolonger l'action des neurotransmetteurs endogènes ;
\item le \textbf{rôle intégrateur} du neurone (sommations temporelle, spatiale, spatio-temporelle) fait de la synapse un véritable centre de décision cellulaire, vulnérable aux perturbations quantitatives (dose) et qualitatives (nature de la substance) ;
\item les intoxications au tramadol (fléau sanitaire jeune au Cameroun) et aux pesticides (risque agricole majeur), s'expliquent par des mécanismes synaptiques précis et distincts (récepteurs opioïdes vs enzyme AChE) : la prévention efficace repose sur la \textbf{compréhension scientifique rigoureuse}, non sur la peur irrationnelle.
\end{itemize}

Tu es désormais capable d'agir comme un véritable relais de santé publique dans ton environnement : c'est là l'essence du \cle{savoir-être} visé par le programme officiel.

\begin{savaistu}[Contexte camerounais \& Santé publique]
\textbf{Tramadol au Cameroun :} Le tramadol (souvent \SI{100}{mg} ou \SI{200}{mg}, détourné de la filière médicale) est la drogue la plus consommée par les jeunes (enquête MINSANTE/ONUDD 2022). Son usage récréatif (« Tramol », « Tramadol Night ») expose à une dépression respiratoire mortelle, d'autant plus que la marge thérapeutique est étroite et la tolérance s'installe vite, poussant aux surdoses.
\textbf{Pesticides organophosphorés :} L'agriculture camerounaise (cacao, banane, coton, maraîchage) utilise massivement des OP (chlorpyrifos, malathion, parathion). L'absence d'EPI (équipements de protection individuelle) et la méconnaissance des délais de rentrée entraînent des intoxications aiguës fréquentes (centres antipoison de Yaoundé/Douala) et des neuropathies périphériques chroniques.
\textbf{Numéro d'urgence :} Centre Antipoison du Cameroun (CAPC) - Hôpital Central de Yaoundé : \textbf{22 22 04 44} / \textbf{22 22 04 45}.
\end{savaistu}

\begin{attention}[Pièges fréquents au Bac - Vigilance rédactionnelle]
\begin{itemize}
\item \textbf{Confusion PPSE/PPSI :} Un PPSE est une \emph{dépolarisation} (vers 0 mV, ex. $+5$ mV). Un PPSI est une \emph{hyperpolarisation} (vers -90 mV, ex. $-5$ mV). Ne pas écrire « PPSE inhibiteur ».
\item \textbf{Sommation temporelle vs spatiale :} Temporelle = \textbf{une} afférence, stimulations \textbf{rapprochées dans le temps}. Spatiale = \textbf{plusieurs} afférences, stimulation \textbf{simultanée}. Justifier toujours par les données (intervalle, nombre d'afférences).
\item \textbf{Cible du tramadol :} Ne pas dire « il bloque les récepteurs ». Il les \textbf{active} (agoniste) $\to$ effet inhibiteur sur le neurone (PPSI). Préciser « récepteurs opioïdes $\mu$ ».
\item \textbf{Organophosphorés :} Ils n'agissent pas sur le récepteur mais sur l'\textbf{enzyme} (AChE). Conséquence : excès de médiateur endogène (ACh), pas mimétisme direct.
\item \textbf{Unités :} Fréquence en \textbf{mov/min} ou \textbf{bpm}, potentiel en \textbf{mV}, temps en \textbf{ms}, concentration en \textbf{mmol/L} (ou $\mu$mol/L). Respecter les chiffres significatifs.
\item \textbf{Synapse électrique :} Ne pas l'oublier si le sujet la demande (jonction communicante, connexines, bidirectionnelle, pas de délai, synchronisation).
\end{itemize}
\end{attention}

\newpage

\coursec{Méthodes et exercices résolus}

\coursec{I. Structure et types anatomiques des synapses}

\begin{definition}[Synapse : définition et ultrastructure]
Une \cle{synapse} est une zone de contact fonctionnel spécialisée entre un élément \emph{présynaptique} (terminaison axonale le plus souvent) et un élément \emph{postsynaptique} (dendrite, corps cellulaire, fibre musculaire ou cellule glandulaire), assurant la transmission de l'influx nerveux. L'\cle{ultrastructure} d'une synapse chimique type comprend :
\begin{itemize}
\item \textbf{Côté présynaptique :} un \cle{bouton synaptique} (renflement terminal) contenant de nombreuses \cle{vésicules synaptiques} ($\approx$ \SI{40}{nm} de diamètre) stockant le neuromédiateur, des \cle{mitochondries} abondantes (fournissant l'ATP pour l'exocytose et le recyclage vésiculaire), des canaux \ce{Ca^{2+}} voltage-dépendants et des zones de libération (\emph{active zones}) riches en protéines SNARE.
\item \textbf{Espace intercellulaire :} la \cle{fente synaptique} (largeur \SI{15}{nm} à \SI{20}{nm}), traversée par des fibres de la matrice extracellulaire et contenant des enzymes de dégradation du médiateur (ex. acétylcholinestérase à la plaque motrice).
\item \textbf{Côté postsynaptique :} une membrane épaissie (\cle{densité postsynaptique}) portant des \cle{récepteurs spécifiques} (canaux ioniques chimiodépendants ou récepteurs métabotropes) et des protéines de scaffolding.
\end{itemize}
\end{definition}

\begin{propriete}[Types anatomiques de synapses]
Selon la nature de l'élément postsynaptique, on distingue trois types anatomiques :
\begin{center}
\begin{tabularx}{\linewidth}{|l|X|X|X|}\hline
\rowcolor{popblueL}
\textbf{Type} & \textbf{Élément présynaptique} & \textbf{Élément postsynaptique} & \textbf{Exemple / Localisation} \\ \hline
\cle{Neuro-neuronique} & Axone (bouton terminal) & Dendrite (synapse axo-dendritique) ou corps cellulaire (axo-somatique) ou axone (axo-axonique) & Circuits du SNC (cortex, moelle épinière) \\ \hline
\cle{Neuro-musculaire} (Plaque motrice) & Axone moteur (terminaison en plaque) & Fibre musculaire striée (plaque motrice) & Jonction neuromusculaire squelettique \\ \hline
\cle{Neuro-glandulaire} & Axone (système nerveux autonome) & Cellule glandulaire (exocrine ou endocrine) & Glandes sudoripares, médullosurrénale \\ \hline
\end{tabularx}
\end{center}
\end{propriete}

\courssub{Schématiser la structure et l'ultrastructure d'une synapse}

\begin{methode}[Schématiser une synapse chimique en 6 étapes rigoureuses]
\begin{enumerate}
\item \textbf{Identifier et flécher le sens de transmission} : la synapse est unidirectionnelle (présynaptique $\rightarrow$ postsynaptique). Tracer une flèche épaisse dès le début.
\item \textbf{Dessiner le bouton synaptique} (élément présynaptique) : forme arrondie/ovoïde. Représenter la membrane présynaptique. À l'intérieur : vésicules synaptiques (petits cercles réguliers, $\approx$ \SI{40}{nm}) et mitochondries (allongées, avec crêtes).
\item \textbf{Dessiner la fente synaptique} : espace parallèle aux membranes, largeur \SI{15}{nm}--\SI{20}{nm} (échelle respectée : fente $\ll$ bouton).
\item \textbf{Dessiner l'élément postsynaptique} : membrane postsynaptique (souvent plus épaisse au niveau de la densité postsynaptique). Ajouter des récepteurs (protéines transmembranaires) côté fente.
\item \textbf{Légender avec le vocabulaire exact du programme} : bouton synaptique, vésicules synaptiques, membrane présynaptique, fente synaptique, récepteurs postsynaptiques (canaux chimiodépendants), membrane postsynaptique, mitochondrie, matrice extracellulaire.
\item \textbf{Préciser l'échelle} : indiquer \SI{100}{nm} ou \SI{200}{nm} en barre d'échelle ; valorise la copie au Bac.
\end{enumerate}
\textbf{Réflexe Bac} : un schéma sans titre, sans légende, sans flèche de sens, ou avec des proportions aberrantes (vésicules grosses comme des mitochondries) perd des points systématiquement.
\end{methode}

\begin{exoresolu}[Schéma annoté d'une synapse chimique neuro-neuronique]
À partir de l'observation en microscopie électronique d'un contact axo-somatique montrant un renflement vésiculeux riche en mitochondries face à un corps cellulaire, séparés par un espace fin de $\approx$ \SI{20}{nm}, réaliser le schéma ultrastructurel annoté en indiquant le sens de l'influx.
\tcblower
\textbf{Raisonnement.} Le renflement contenant vésicules et mitochondries est le bouton synaptique (présynaptique). L'espace fin est la fente synaptique. Le corps cellulaire est l'élément postsynaptique. L'influx va du bouton vers le corps cellulaire.

\begin{popfigure}
\begin{center}
\begin{tikzpicture}[scale=1.1, every node/.style={font=\small}]
% Axone arrivant
\draw[thick, popblue] (-4,2.2) -- (-2.5,2.2);
\node[popblue, left] at (-4,2.2) {Axone};
% Bouton synaptique
\draw[fill=popblue!15, thick, popblue] (0,2.2) ellipse (2.0 and 1.2);
\node[popblue] at (0,3.7) {\textbf{Bouton synaptique (présynaptique)}};
% Vésicules
\foreach \x/\y in {-1.3/2.6,-0.5/2.8,0.4/2.5,1.2/2.7,-0.9/2.0,0.1/1.9,0.9/2.0,-1.5/1.9}{
  \draw[fill=popgold!80, popgold] (\x,\y) circle (0.12);}
% Mitochondries
\draw[fill=poporange!60, poporange] (-1.5,3.1) ellipse (0.4 and 0.18);
\draw[fill=poporange!60, poporange] (1.5,3.1) ellipse (0.4 and 0.18);
\draw[poporange] (-1.5,3.1) -- (-1.3,3.25) (-1.5,3.1) -- (-1.7,3.25); % crêtes
% Membrane présynaptique (courbe basale)
\draw[very thick, popblue] (-1.8,1.1) .. controls (0,0.8) .. (1.8,1.1);
% Fente synaptique
\fill[popcream] (-1.8,0.6) rectangle (1.8,1.1);
% Membrane postsynaptique (densité)
\draw[very thick, popgreen] (-2.4,0.6) .. controls (0,0.3) .. (2.4,0.6);
% Récepteurs (côté fente)
\foreach \x in {-1.0,-0.4,0.2,0.8,1.4}{
  \draw[fill=popgreen!60, popgreen] (\x,0.42) rectangle ++(0.16,0.22);}
% Élément postsynaptique (corps cellulaire)
\draw[fill=popgreen!10, popgreen, thick] (-3.0,-1.5) rectangle (3.0,0.35);
\node[popgreen] at (0,-0.8) {\textbf{Élément postsynaptique (corps cellulaire / dendrite)}};
% Noyau (optionnel)
\draw[fill=popgreen!30, popgreen] (-0.5,-1.0) circle (0.35);
% Flèche sens transmission
\draw[-{Stealth[length=5mm]}, ultra thick, popink] (0,1.7) -- (0,0.85) node[midway, right=2mm, popink, font=\bfseries] {Influx};
% Légendes fléchées (propres)
\draw[-{Stealth[length=2mm]}, popmuted] (-1.0,2.6) -- (-3.2,3.0) node[left, black] {1. Vésicules synaptiques};
\draw[-{Stealth[length=2mm]}, popmuted] (1.5,3.1) -- (3.5,3.4) node[right, black] {2. Mitochondries};
\draw[-{Stealth[length=2mm]}, popmuted] (1.8,1.1) -- (3.5,1.2) node[right, black] {3. Membrane présynaptique};
\draw[-{Stealth[length=2mm]}, popmuted] (0,0.85) -- (-3.5,0.9) node[left, black] {4. Fente synaptique ($\approx$ \SI{20}{nm})};
\draw[-{Stealth[length=2mm]}, popmuted] (0.6,0.5) -- (3.5,0.1) node[right, black] {5. Récepteurs postsynaptiques};
\draw[-{Stealth[length=2mm]}, popmuted] (-2.4,0.6) -- (-3.8,0.2) node[left, black] {6. Membrane postsynaptique};
% Échelle
\draw[|-|, thick, black] (3.2, -1.5) -- (3.2, -1.0) node[midway, right=1mm] {\SI{500}{nm}};
\end{tikzpicture}
\end{center}
\legende{Ultrastructure d'une synapse chimique neuro-neuronique (contact axo-somatique). Les vésicules (1) stockent le médiateur ; les mitochondries (2) fournissent l'ATP. La fente (4) sépare la membrane présynaptique (3) de la membrane postsynaptique (6) portant les récepteurs (5). Flèche rose : sens de transmission.}
\end{popfigure}

\textbf{Justifications.} Les vésicules libèrent le neuromédiateur par exocytose \ce{Ca^{2+}}-dépendante ; les mitochondries assurent l'énergie (ATP) pour le cyclage vésiculaire et la recapture ; les récepteurs postsynaptiques assurent la spécificité et la transduction du signal chimique en signal électrique.
\textbf{Conclusion.} La structure observée est une synapse chimique de type neuro-neuronique ; la transmission est unidirectionnelle du bouton présynaptique vers l'élément postsynaptique.
\end{exoresolu}

\courssub{Identifier les synapses à transmission électrique et à transmission chimique sur des électronographies}

\begin{methode}[Distinguer synapse chimique et synapse électrique sur un document (MEB)]
\begin{enumerate}
\item \textbf{Mesurer la largeur de la fente} (échelle fournie) :
\begin{itemize}
\item Fente \textbf{large} (\SI{15}{nm} à \SI{20}{nm}) $\Rightarrow$ \cle{synapse chimique}.
\item Fente \textbf{très étroite} (\SI{2}{nm} à \SI{3}{nm}) $\Rightarrow$ \cle{synapse électrique} (jonction communicante / \emph{gap junction}).
\end{itemize}
\item \textbf{Rechercher les vésicules synaptiques} dans l'élément présynaptique :
\begin{itemize}
\item Présentes (nombreuses) $\Rightarrow$ synapse chimique.
\item Absentes $\Rightarrow$ synapse électrique.
\end{itemize}
\item \textbf{Observer les membranes au niveau du contact} :
\begin{itemize}
\item Synapse électrique : présence de \cle{connexons} (structures protéiques hexamériques formant un pore continu entre les deux cytoplasmes).
\item Synapse chimique : pas de continuité cytoplasmique ; densité postsynaptique visible côté receveur.
\end{itemize}
\item \textbf{Conclure sur le mode de transmission} :
\begin{itemize}
\item Chimique : libération de neuromédiateur, diffusion, fixation récepteur $\rightarrow$ \textbf{lente} (\SI{0.5}{ms}--\SI{1}{ms} délai), \textbf{unidirectionnelle}, modulable (sommation, plasticité).
\item Électrique : passage direct du courant ionique (\emph{local current flow}) via connexons $\rightarrow$ \textbf{quasi instantanée}, souvent \textbf{bidirectionnelle}, synchronisation rapide (ex. escape reflexes, rythme cardiaque).
\end{itemize}
\end{enumerate}
\textbf{Réflexe Bac} : Citer \emph{deux critères observables} (largeur fente + vésicules/connexons) pour justifier l'identification. Ne jamais écrire « on voit que c'est chimique » sans preuve visuelle.
\end{methode}

\begin{exoresolu}[Analyse de deux électronographies de synapses]
Le document 1 montre un contact interneuronal avec une fente de \SI{20}{nm} et de nombreuses vésicules claires dans l'élément présynaptique. Le document 2 montre un contact intercellulaire où les membranes sont distantes de \SI{3}{nm}, traversées par des particules protéiques régulières (connexons) ; aucune vésicule n'est visible.
\begin{enumerate}
\item Identifier le type de synapse pour chaque document en justifiant par deux critères.
\item Préciser le mode de transmission, la vitesse relative et la directionnalité pour chacune.
\end{enumerate}
\tcblower
\textbf{1. Identification justifiée.}
\begin{itemize}
\item \textbf{Document 1} : Fente large (\SI{20}{nm}) \textbf{et} présence de vésicules synaptiques $\Rightarrow$ \cle{Synapse chimique}.
\item \textbf{Document 2} : Fente très étroite (\SI{3}{nm}) \textbf{et} présence de connexons (ponts intercellulaires) / absence de vésicules $\Rightarrow$ \cle{Synapse électrique}.
\end{itemize}
\textbf{2. Caractéristiques de transmission.}
\begin{center}
\begin{tabularx}{\linewidth}{|l|X|X|X|}\hline
\rowcolor{popblueL}
\textbf{Critère} & \textbf{Document 1 (Chimique)} & \textbf{Document 2 (Électrique)} \\ \hline
Mode & Diffusion neuromédiateur (ex. glutamate, GABA) & Courant ionique direct via connexons \\ \hline
Vitesse & Lente (délai synaptique $\approx$ \SI{0,5--1}{ms}) & Quasi instantanée (pas de délai) \\ \hline
Directionnalité & Unidirectionnelle (présynaptique $\rightarrow$ postsynaptique) & Bidirectionnelle (souvent) \\ \hline
Modulation & Forte (sommation, facilitation, inhibition) & Faible (principalement résistance jonctionnelle) \\ \hline
\end{tabularx}
\end{center}
\textbf{Conclusion.} Le document 1 illustre une synapse chimique (transmission médiate, lente, unidirectionnelle, plastique) ; le document 2 une synapse électrique (couplage direct, rapide, synchronisant).
\end{exoresolu}

\coursec{II. Nature de la transmission synaptique : l'expérience de Loewi}

\begin{experience}[L'expérience de Loewi (1921) : preuve de la transmission chimique]
\textbf{Matériel :} Deux cœurs de grenouille isolés, maintenus en survie dans du liquide de Ringer oxygéné à température constante. Cœur A (avec nerf vague intact), Cœur B (dénervé). Système de perfusion permettant de transférer le liquide du bain A vers le bain B.
\textbf{Protocole :}
\begin{enumerate}
\item Laisser battre les deux cœurs spontanément (rythme de base).
\item \textbf{Stimuler électriquement le nerf vague du cœur A} $\rightarrow$ ralentissement marqué du cœur A.
\item Recueillir le liquide de perfusion ayant baigné le cœur A \emph{pendant la stimulation} et le verser immédiatement sur le cœur B (témoin : liquide de perfusion du cœur A \emph{sans stimulation} versé sur B').
\item Observer la fréquence du cœur B.
\end{enumerate}
\textbf{Résultats :} Le cœur B ralentit significativement au contact du liquide du cœur A stimulé. Le cœur B' (témoin) ne change pas de rythme.
\textbf{Interprétation :} La stimulation du nerf vague libère une substance chimique soluble (l'\cle{acétylcholine}, nommée « Vagusstoff » par Loewi) dans le liquide. Cette substance diffuse et agit sur les récepteurs du cœur B, reproduisant l'effet inhibiteur.
\textbf{Conclusion historique :} La transmission neuro-efféctrice (et par extension synaptique) est de nature \cle{chimique}. Loewi reçoit le Prix Nobel 1936 (partagé avec Dale).
\end{experience}

\begin{methode}[Analyser l'expérience de Loewi pour le Bac]
\begin{enumerate}
\item \textbf{Problématique} : Nature de la transmission nerveuse (électrique vs chimique) ?
\item \textbf{Dispositif} : Deux organes effecteurs isolés, un innervé (A), un dénerve (B), transfert de fluide.
\item \textbf{Variable indépendante} : Stimulation du nerf vague (présence/absence).
\item \textbf{Variable dépendante} : Fréquence cardiaque du cœur B.
\item \textbf{Témoin indispensable} : Liquide de perfusion d'un cœur non stimulé (ou cœur B non perfusé) $\rightarrow$ pas de ralentissement.
\item \textbf{Argument clé} : Le cœur B est \textbf{dénervé} et \textbf{non stimulé} ; son ralentissement exclut une propagation électrique (continuité cytoplasmique) et prouve l'action d'un facteur humoral (chimique) libéré par le nerf.
\item \textbf{Conclusion} : Transmission chimique par libération d'un neuromédiateur (acétylcholine).
\end{enumerate}
\end{methode}

\begin{exoresolu}[L'expérience de Loewi : analyse complète]
Otto Loewi isole deux cœurs de grenouille. Le cœur 1 garde son nerf vague ; le cœur 2 est dénerve. La stimulation du nerf vague du cœur 1 ralentit ce cœur. Le liquide de perfusion du cœur 1 (stimulé) est transféré sur le cœur 2 : le cœur 2 ralentit.
\begin{enumerate}
\item Formuler le problème scientifique posé.
\item Analyser les résultats en identifiant le témoin et en expliquant pourquoi la transmission électrique est exclue.
\item Conclure sur la nature de la transmission et nommer le médiateur.
\end{enumerate}
\tcblower
\textbf{1. Problème.} Comment l'influx nerveux est-il transmis du nerf vague au muscle cardiaque : par continuité électrique ou par libération d'une substance chimique ?

\textbf{2. Analyse.}
\begin{itemize}
\item \textbf{Fait 1} : Stimulation nerf vague $\rightarrow$ ralentissement cœur 1. Le nerf vague est inhibiteur.
\item \textbf{Fait 2} : Cœur 2 (dénervé, non stimulé) ralentit au contact du liquide du cœur 1 stimulé.
\item \textbf{Témoin} : Liquide du cœur 1 non stimulé $\rightarrow$ pas d'effet sur cœur 2 (ou cœur 2 non perfusé).
\item \textbf{Interprétation} : L'absence de connexion nerveuse avec le cœur 2 élimine l'hypothèse électrique. Le ralentissement est dû à une substance chimique libérée par le nerf vague dans le bain du cœur 1, diffusée dans le liquide, et agissant sur le cœur 2. Cette substance est l'\cle{acétylcholine}.
\end{itemize}
\textbf{3. Conclusion.} L'expérience démontre la \cle{transmission chimique} au niveau de la synapse neuro-cardiaque : le nerf libère l'acétylcholine qui se fixe sur des récepteurs muscariniques du myocarde, ouvrant des canaux \ce{K+} (hyperpolarisation, ralentissement).
\end{exoresolu}

\coursec{III. Fonctionnement des synapses chimiques et électriques}

\courssub{Transmission chimique : mécanisme détaillé}

\begin{methode}[Décrire les 6 étapes de la transmission synaptique chimique]
\begin{enumerate}
\item \textbf{Arrivée du potentiel d'action (PA)} au bouton synaptique $\rightarrow$ dépolarisation de la membrane présynaptique.
\item \textbf{Ouverture des canaux \ce{Ca^{2+}} voltage-dépendants} $\rightarrow$ entrée massive d'ions \ce{Ca^{2+}} dans le bouton (gradient électrochimique favorable).
\item \textbf{Exocytose \ce{Ca^{2+}}-dépendante} : Les protéines SNARE (synaptobrévine, syntaxine, SNAP-25) médiatent la fusion des vésicules avec la membrane présynaptique $\rightarrow$ libération du neuromédiateur dans la fente (quantal). Consomme de l'ATP.
\item \textbf{Diffusion et fixation} : Le médiateur traverse la fente par diffusion passive et se fixe sur ses \cle{récepteurs spécifiques} (canaux chimiodépendants ionotropes ou récepteurs métabotropes) de la membrane postsynaptique.
\item \textbf{Potentiel postsynaptique} :
\begin{itemize}
\item \cle{PPSE} (Excitateurs) : Ouverture canaux \ce{Na+} (et \ce{K+}) $\rightarrow$ dépolarisation locale (ex. Glutamate, Acétylcholine sur récepteur nicotinique). Rapproche du seuil.
\item \cle{PPSI} (Inhibiteurs) : Ouverture canaux \ce{Cl-} (entrée) ou \ce{K+} (sortie) $\rightarrow$ hyperpolarisation ou stabilisation (ex. GABA, Glycine, Acétylcholine sur récepteur muscarinique cardiaque). Éloigne du seuil.
\end{itemize}
\item \textbf{Inactivation du médiateur (fin du signal)} :
\begin{itemize}
\item \textbf{Hydrolyse enzymatique} dans la fente (ex. Acétylcholinestérase $\rightarrow$ Choline + Acétate ; la choline est recaptée).
\item \textbf{Recapture présynaptique} par transporteurs spécifiques (ex. SERT pour sérotonine, DAT pour dopamine, NET pour noradrénaline, GAT pour GABA).
\item \textbf{Diffusion} hors de la fente.
\end{itemize}
\end{enumerate}
\end{methode}

\courssub{Transmission électrique : mécanisme et spécificités}

\begin{propriete}[Synapse électrique (Jonction communicante / Gap Junction)]
\begin{itemize}
\item \textbf{Structure} : Connexons (hémicanaux) formés de 6 \cle{connexines} s'alignant entre deux cellules adjacentes. Pore central ($\approx$ \SI{1.5}{nm}) perméable aux ions et petites molécules (\textless \SI{1}{kDa} : \ce{Ca^{2+}}, \ce{cAMP}, \ce{IP3}).
\item \textbf{Fonctionnement} : Passage direct du courant ionique (\emph{local current flow}) d'une cellule à l'autre. Pas de médiateur, pas de récepteur, pas de délai synaptique.
\item \textbf{Propriétés} : Transmission \textbf{rapide} (quasi instantanée), souvent \textbf{bidirectionnelle} (sauf rectification), \textbf{faible plasticité}. Permet la \textbf{synchronisation} d'ensembles cellulaires.
\item \textbf{Localisation} : Cœur (cellules myocardiques, nœud sinusal), muscles lisses viscéraux, certaines synapses du SNC (réseaux inhibiteurs interneuroniques, rétine), astrocytes.
\end{itemize}
\end{propriete}

\courssub{Sommations et rôle intégrateur du neurone}

\begin{methode}[Identifier et expliquer les types de sommation]
\begin{enumerate}
\item \textbf{Analyser le protocole de stimulation} :
\begin{itemize}
\item \cle{Sommation temporelle} : Une \textbf{seule} afférence excitatrice stimulée \textbf{répétitivement} à haute fréquence (intervalle $<$ durée du PPSE). Addition successive des PPSE dans le temps.
\item \cle{Sommation spatiale} : \textbf{Plusieurs} afférences excitatrices distinctes stimulées \textbf{simultanément}. Addition simultanée des PPSE dans l'espace (sur la membrane postsynaptique).
\item \cle{Sommation spatio-temporelle} : Combinaison des deux (plusieurs afférences stimulées à des instants rapprochés).
\end{itemize}
\item \textbf{Lire l'enregistrement intracellulaire} (potentiel de membrane $V_m$) :
\begin{itemize}
\item PPSE = dépolarisation (ex. $-70 \rightarrow -60$ mV). PPSI = hyperpolarisation (ex. $-70 \rightarrow -75$ mV).
\item Addition \textbf{algébrique} : $V_m = V_{repos} + \sum PPSE + \sum PPSI$.
\end{itemize}
\item \textbf{Comparer au seuil de dépolarisation} ($V_{seuil} \approx \SI{-55}{mV}$ pour $V_{repos} = \SI{-70}{mV}$) :
\begin{itemize}
\item $V_m \ge V_{seuil}$ $\rightarrow$ Déclenchement PA (loi du tout ou rien).
\item $V_m < V_{seuil}$ $\rightarrow$ Pas de PA (potentiel gradué).
\end{itemize}
\item \textbf{Conclure sur le rôle intégrateur} : Le neurone postsynaptique calcule en temps réel la somme pondérée de ses entrées excitatrices et inhibitrices. Il agit comme un \cle{coïncidence detector} (sommation spatiale) et un \cle{intégrateur temporel} (sommation temporelle).
\end{enumerate}
\end{methode}

\begin{exoresolu}[Sommations au niveau d'un motoneurone spinal]
Un motoneurone M ($V_{repos} = \SI{-70}{mV}$, $V_{seuil} = \SI{-55}{mV}$) reçoit deux afférences excitatrices A et B (chaque PPSE unitaire = \SI{+10}{mV}).
\begin{itemize}
\item Exp 1 : Stimulation unique de A $\rightarrow$ PPSE \SI{+10}{mV}. Pas de PA.
\item Exp 2 : Double stimulation de A (intervalle \SI{2}{ms}) $\rightarrow$ PA déclenché.
\item Exp 3 : Stimulation simultanée de A et B $\rightarrow$ PA déclenché.
\end{itemize}
\begin{enumerate}
\item Expliquer l'absence de PA en Exp 1.
\item Identifier le type de sommation en Exp 2 et Exp 3, justifier par le calcul.
\end{enumerate}
\tcblower
\textbf{1. Expérience 1.} $V_m = -70 + 10 = \SI{-60}{mV}$. Or $\SI{-60}{mV} < \SI{-55}{mV}$ (plus négatif que le seuil). La dépolarisation est \textbf{sous-liminaire} : aucun canal \ce{Na+} voltage-dépendant ne s'ouvre massivement, pas de PA.

\textbf{2. Identification et calculs.}
\begin{itemize}
\item \textbf{Exp 2} : Une seule afférence (A) stimulée deux fois rapprochées. Le second PPSE s'additionne à la phase de dépolarisation résiduelle du premier. $V_m = -70 + 10 + 10 = \SI{-50}{mV}$. $\SI{-50}{mV} > \SI{-55}{mV}$ $\rightarrow$ \textbf{Seuil atteint}. \cle{Sommation temporelle}.
\item \textbf{Exp 3} : Deux afférences distinctes (A et B) stimulées en même temps. $V_m = -70 + 10 + 10 = \SI{-50}{mV}$. $\SI{-50}{mV} > \SI{-55}{mV}$ $\rightarrow$ \textbf{Seuil atteint}. \cle{Sommation spatiale}.
\end{itemize}
\textbf{Conclusion.} Le motoneurone intègre l'information : la sommation temporelle (fréquence codée) et la sommation spatiale (convergence codée) permettent de transformer des signaux gradués sous-seuil en un signal tout-ou-rien (PA) propagé vers le muscle.
\end{exoresolu}

\courssub{Exploiter un enregistrement de potentiel postsynaptique}

\begin{methode}[Lire et interpréter une courbe $V_m(t)$ postsynaptique]
\begin{enumerate}
\item \textbf{Repérer les valeurs de référence} : $V_{repos}$ (ligne de base), $V_{seuil}$ (souvent indiqué), échelle des temps (ms).
\item \textbf{Qualifier chaque déflexion} :
\begin{itemize}
\item Montée vers 0 mV (dépolarisation) $\rightarrow$ \cle{PPSE} (médiateur excitateur : Glutamate, ACh nicotinique).
\item Descente vers valeurs plus négatives (hyperpolarisation) $\rightarrow$ \cle{PPSI} (médiateur inhibiteur : GABA, Glycine).
\end{itemize}
\item \textbf{Tester le franchissement du seuil} : PA = pointe aiguë, amplitude $\approx$ \SI{100}{mV}, durée \SI{1--2}{ms}, loi du tout ou rien.
\item \textbf{Analyser les interactions} : Addition algébrique PPSE + PPSI. Une PPSI peut empêcher un PA (inhibition « veto » ou « shunting »).
\item \textbf{Rédiger l'interprétation} : Lier chaque événement de la courbe à une libération synaptique précise (quelle afférence, quel médiateur probable).
\end{enumerate}
\end{methode}

\begin{exoresolu}[Interprétation d'un enregistrement complexe]
Enregistrement sur un neurone ($V_{repos} = \SI{-70}{mV}$, $V_{seuil} = \SI{-55}{mV}$) :
\begin{itemize}
\item $t_1$ : Dépolarisation \SI{+8}{mV} (PPSE).
\item $t_2$ (\SI{3}{ms} après $t_1$) : Second PPSE \SI{+8}{mV} s'additionnant au premier $\rightarrow$ PA.
\item $t_3$ : Stimulation afférence C $\rightarrow$ Hyperpolarisation \SI{-5}{mV} (PPSI) + Stimulation afférence D $\rightarrow$ PPSE \SI{+8}{mV} simultanées $\rightarrow$ Pas de PA.
\end{itemize}
\begin{enumerate}
\item Nommer les potentiels à $t_1$, $t_2$, $t_3$.
\item Expliquer la survenue du PA à $t_2$ et son absence à $t_3$ par calcul.
\end{enumerate}
\tcblower
\textbf{1. Nature des potentiels.}
$t_1$ : \cle{PPSE} (dépolarisation sous-seuil). $t_2$ : \cle{Sommation temporelle} de deux PPSE $\rightarrow$ \cle{Potentiel d'action}. $t_3$ : \cle{PPSI} (hyperpolarisation) + PPSE.

\textbf{2. Calculs et explication.}
\begin{itemize}
\item \textbf{À $t_2$} : Le premier PPSE (\SI{+8}{mV}, $V_m = \SI{-62}{mV}$) n'est pas terminé. Le second PPSE (\SI{+8}{mV}) s'y somme : $V_m = -70 + 8 + 8 = \SI{-54}{mV}$. $\SI{-54}{mV} > \SI{-55}{mV}$ $\rightarrow$ \textbf{Seuil franchi} $\rightarrow$ PA.
\item \textbf{À $t_3$} : Addition algébrique PPSE + PPSI : $V_m = -70 + 8 + (-5) = \SI{-67}{mV}$. $\SI{-67}{mV} < \SI{-55}{mV}$ $\rightarrow$ \textbf{Seuil non atteint}. L'inhibition (entrée \ce{Cl-} ou sortie \ce{K+}) « annule » l'excitation (entrée \ce{Na+}). C'est l'\cle{intégration synaptique} : le neurone « décide » de ne pas feu.
\end{itemize}
\end{exoresolu}

\coursec{IV. Effets des substances sur la transmission synaptique}

\begin{definition}[Classification pharmacologique des substances synaptiques]
Les substances exogènes (drogues, médicaments, toxines, pesticides) perturbent la transmission en ciblant une étape précise :
\begin{center}
\begin{tabularx}{\linewidth}{|l|X|X|l|}\hline
\rowcolor{popblueL}
\textbf{Cible moléculaire} & \textbf{Type d'action} & \textbf{Mécanisme moléculaire} & \textbf{Exemples classiques} \\ \hline
\cle{Récepteur postsynaptique} & \textbf{Agoniste / Mimétique} & Fixe le récepteur $\rightarrow$ l'active (mime le médiateur endogène) & Nicotine (nAChR), Muscarine (mAChR), Morphine/Héroïne ($\mu$-opioïde), Benzodiazépines (GABA$_A$), THC (CB1) \\ \hline
& \textbf{Antagoniste / Bloqueur} & Fixe le récepteur \textbf{sans l'activer} $\rightarrow$ empêche la fixation du médiateur endogène & \textbf{Curare} (nAChR compétitif), Atropine (mAChR), Naloxone (opioïde), Antipsychotiques (D2) \\ \hline
\cle{Enzyme de dégradation (fente)} & \textbf{Inhibiteur enzymatique} & Bloque l'hydrolyse $\rightarrow$ accumulation du médiateur dans la fente $\rightarrow$ surstimulation prolongée & \textbf{Organophosphorés} (inhibent AChE), Physostigmine, Névralgiques (Sarin, VX) \\ \hline
\cle{Machinerie de libération (présynaptique)} & \textbf{Bloqueur de l'exocytose} & Clive protéines SNARE / bloque canaux \ce{Ca^{2+}} $\rightarrow$ pas de libération vésiculaire & \textbf{Toxine botulique} (BoNT), Toxine tétanique (TeNT), $\omega$-Conotoxine \\ \hline
\cle{Transporteur de recapture (présynaptique)} & \textbf{Bloqueur de recapture} & Empêche le nettoyage du médiateur $\rightarrow$ prolongation/action accrue dans la fente & \textbf{Cocaïne} (DAT/NET/SERT), Antidépresseurs ISRS (SERT), Tricycliques (NET/SERT), Amphétamines (inverseur DAT) \\ \hline
\cle{Synthèse / Stockage (présynaptique)} & \textbf{Inhibiteur de synthèse / Vésiculaire} & Déplete les réserves de médiateur & $\alpha$-Méthyl-p-tyrosine (catecholamines), Réserpine (VMAT2) \\ \hline
\end{tabularx}
\end{center}
\end{definition}

\begin{methode}[Rédiger la chaîne causale complète (Réflexe Bac obligatoire)]
Pour toute substance, écrire : 
\textbf{Substance} $\xrightarrow{\text{cible}}$ \textbf{Perturbation moléculaire} $\xrightarrow{\text{conséquence}}$ \textbf{Modification du potentiel postsynaptique (PPSE/PPSI/PA)} $\xrightarrow{\text{effet}}$ \textbf{Phénotype physiologique / Symptôme}.
\begin{itemize}
\item \textit{Exemple Curare} : Curare $\xrightarrow{\text{récepteur nAChR}}$ Antagoniste compétitif $\xrightarrow{}$ Blocage fixation ACh $\xrightarrow{}$ Aucun PPSE $\xrightarrow{}$ \textbf{Paralysie flasque} (arrêt respiratoire si diaphragme).
\item \textit{Exemple Organophosphoré} : Insecticide $\xrightarrow{\text{AChE}}$ Inhibition irréversible $\xrightarrow{}$ Accumulation ACh $\xrightarrow{}$ PPSE prolongés/répétés $\xrightarrow{}$ \textbf{Tétanie, convulsions, mort par asphyxie}.
\end{itemize}
\end{methode}

\begin{exoresolu}[Analyse mécanistique : Curare, Toxine botulique, Organophosphoré]
Préparation nerf-muscle (plaques motrices, médiateur = Acétylcholine). Stimulation nerf $\rightarrow$ contraction normale.
\begin{itemize}
\item \textbf{Exp A (Curare)} : Stimulation nerf $\rightarrow$ Pas de contraction. ACh exogène sur muscle $\rightarrow$ Pas de contraction.
\item \textbf{Exp B (Toxine botulique)} : Stimulation nerf $\rightarrow$ Pas de contraction. ACh exogène sur muscle $\rightarrow$ Contraction normale.
\item \textbf{Exp C (Organophosphoré)} : Stimulation unique $\rightarrow$ Contractions prolongées, répétées (tétanie).
\end{itemize}
Expliquer le mécanisme précis de chaque substance.
\tcblower
\textbf{Expérience A — Curare.}
\begin{itemize}
\item \textbf{Observation clé} : L'ACh exogène est inefficace $\Rightarrow$ Cible \textbf{postsynaptique} (récepteur).
\item \textbf{Mécanisme} : Le curare est un \cle{antagoniste compétitif réversible} des récepteurs nicotiniques (\ce{nAChR}) de la plaque motrice. Il occupe le site de fixation de l'ACh sans ouvrir le canal \ce{Na+}/\ce{K+}.
\item \textbf{Conséquence} : Aucun PPSE, aucun PA musculaire $\rightarrow$ \textbf{Paralysie flasque}. Utilisation : poison de flèche (curare), myorelaxant en anesthésie (intubation).
\end{itemize}

\textbf{Expérience B — Toxine botulique (BoNT).}
\begin{itemize}
\item \textbf{Observation clé} : L'ACh exogène fonctionne $\Rightarrow$ Récepteurs intacts. Blocage \textbf{présynaptique}.
\item \textbf{Mécanisme} : La BoNT (protéine \SI{150}{kDa}) est internalisée dans le bouton $\rightarrow$ activité protéolytique (zinc-dépendante) $\rightarrow$ \cle{clivage des protéines SNARE} (synaptobrévine/VAMP, SNAP-25, syntaxine). Fusion vésicules impossible $\rightarrow$ \textbf{Blocage total de l'exocytose}.
\item \textbf{Conséquence} : Aucune libération d'ACh $\rightarrow$ \textbf{Paralysie flasque}. Pathologie : Botulisme (conserves mal stérilisées, miel nourrissons). Toxine la plus puissante connue (DL50 $\approx$ \SI{1}{ng/kg}).
\end{itemize}

\textbf{Expérience C — Insecticide organophosphoré (ex. Parathion, Chlorpyrifos).}
\begin{itemize}
\item \textbf{Observation clé} : Contractions prolongées, répétées (tétanie) après un seul stimulus $\Rightarrow$ Excitation excessive prolongée.
\item \textbf{Mécanisme} : Inhibition \textbf{irréversible} (phosphorylation du sérine active) de l'\cle{acétylcholinestérase (AChE)} dans la fente synaptique. L'ACh libérée n'est plus hydrolysée $\rightarrow$ \textbf{accumulation massive} dans la fente $\rightarrow$ fixation répétée sur \ce{nAChR} $\rightarrow$ PPSE prolongés, dépolarisation maintenue $\rightarrow$ salves de PA musculaires.
\item \textbf{Conséquence} : \textbf{Tétanie spastique}, puis dépolarisation bloquante (fatigue), insuffisance respiratoire. Intoxication fréquente chez agriculteurs non protégés (coton, cacao au Cameroun). Traitement : Atropine (antagoniste mAChR périphérique) + Pralidoxime (réactiveur AChE) + Diazépam (convulsions).
\end{itemize}
\textbf{Synthèse.} Trois cibles, trois issues : Blocage réception (Curare), Blocage libération (BoNT), Blocage dégradation (OP). Tous paralysent mais par des mécanismes radicalement opposés sur le signal synaptique.
\end{exoresolu}

\begin{savaistu}[Contexte camerounais : Pesticides, Drogues et Santé Publique]
\begin{itemize}
\item \textbf{Organophosphorés \& Carbamates} : Largement utilisés dans les filières \textbf{cacaoyère} (capsides, mirides) et \textbf{cotonnière} (jassides, charançons) au Cameroun. L'intoxication aiguë (syndrome cholinergique : SLUDGE - Salivation, Larmoiement, Urination, Diarrhée, Gastro-intestinal, Émétion + myosis, bronchospasme, bradycardie, fasciculations, paralysie) est une urgence vitale. La formation des planteurs (port de l'EPI, respect des délais de rentrée, triple rinçage) est un enjeu majeur de santé au travail.
\item \textbf{Drogues d'abus} :
\begin{itemize}
\item \textbf{Cannabis} (feuilles « bangui »/« zamal ») : Agoniste partiel récepteurs \cle{CB1} (pré-synaptique, inhibe libération GABA/Glutamate/Dopamine) $\rightarrow$ effets psychotropes, mémoire, motricité, appétit.
\item \textbf{Tramadol} (détournement médicament antidouleur) : Agoniste $\mu$-opioïde + inhibiteur recapture sérotonine/noradrénaline $\rightarrow$ risque convulsions, dépression respiratoire, dépendance sévère. Fléau sanitaire majeur chez les jeunes.
\item \textbf{Alcool (éthanol)} : Potentiateur GABA$_A$ (agoniste allostérique) + Antagoniste NMDA (glutamate) $\rightarrow$ dépression SNC, désinhibition, neurotoxicité fœtale (SAF).
\end{itemize}
\item \textbf{Médicaments cibles synapses} : Antidépresseurs ISRS (fluoxétine : bloqueur SERT), Antiparkinsoniens (L-DOPA : précurseur dopamine), Myorelaxants (curarisants en chirurgie), Anticholinestérasiques réversibles (néostigmine, pyridostigmine : myasthénie, réversal curare).
\end{itemize}
\end{savaistu}

\coursec{V. Synthèse et Sensibilisation}

\begin{aretenir}[Tableau de synthèse : Synapse Chimique vs Électrique]
\begin{center}
\begin{tabularx}{\linewidth}{|l|X|X|}\hline
\rowcolor{popblueL}
\textbf{Critère} & \textbf{Synapse Chimique} & \textbf{Synapse Électrique} \\ \hline
Ultrastructure (MEB) & Fente \SI{15--20}{nm}, Vésicules ++, Densité post. & Fente \SI{2--3}{nm}, Connexons, Pas de vésicules \\ \hline
Médiateur & Oui (ACh, Glutamate, GABA, DA, 5-HT, Gly...) & Non (courant ionique direct) \\ \hline
Récepteurs & Oui (Canaux chimiodépendants / Métabotropes) & Non (Pore connexine) \\ \hline
Délai synaptique & \SI{0.5--1}{ms} (libération + diffusion + fixation) & Quasi nul (\textless \SI{0.1}{ms}) \\ \hline
Directionnalité & Unidirectionnelle (polarité vésicules/récepteurs) & Bidirectionnelle (sauf jonctions rectifiantes) \\ \hline
Sommation / Plasticité & Oui (spatiale, temporelle, LTP/LTD, facilitation) & Non (loi d'Ohm, résistance jonctionnelle) \\ \hline
Rôle principal & Traitement de l'information, intégration, apprentissage & Synchronisation rapide, couplage métabolique \\ \hline
Localisation majeure & SNC, Plaque motrice, SNA & Cœur, Muscle lisse, Rétine, Astrocytes, SNC (quelques) \\ \hline
\end{tabularx}
\end{center}
\end{aretenir}

\begin{aretenir}[Effets des substances : Résumé des mécanismes clés pour le Bac]
\begin{itemize}
\item \textbf{Agoniste} (Nicotine, Morphine, Diazépam) $\rightarrow$ \textbf{Mime} le médiateur $\rightarrow$ \textbf{Active} le récepteur $\rightarrow$ Effet \textbf{excitateurs ou inhibiteur} selon récepteur.
\item \textbf{Antagoniste} (Curare, Atropine, Naloxone, Halopéridol) $\rightarrow$ \textbf{Bloque} le récepteur \textbf{sans l'activer} $\rightarrow$ \textbf{Inhibe} l'action du médiateur endogène.
\item \textbf{Inhibiteur d'enzyme} (Organophosphorés, Physostigmine, Donepezil) $\rightarrow$ \textbf{Accumule} le médiateur dans la fente $\rightarrow$ \textbf{Prolonge/Amplifie} l'action (excitatrice ou inhibitrice).
\item \textbf{Bloqueur libération} (Toxine botulique, Tétanique, $\omega$-Conotoxine) $\rightarrow$ \textbf{Empêche} l'exocytose $\rightarrow$ \textbf{Silence} synaptique $\rightarrow$ Paralysie flasque (Botulique) ou spastique (Tétanique : blocage GABA/Glycine interneurones).
\item \textbf{Bloqueur recapture} (Cocaïne, ISRS, Tricycliques, Amphétamine) $\rightarrow$ \textbf{Prolonge} la présence du médiateur dans la fente $\rightarrow$ \textbf{Amplifie} le signal (Dopamine/NA/5-HT).
\end{itemize}
\end{aretenir}

\begin{attention}[Les 5 erreurs fatales au Bac (SVT Terminale D)]
\begin{enumerate}
\item \textbf{Confondre synapse électrique et chimique sur document MEB} : Ne jamais conclure sur un seul critère. \textbf{Obligatoire} : Citer \textbf{largeur de la fente} \textbf{ET} \textbf{présence/absence de vésicules (ou connexons)}. Fente \SI{20}{nm} + vésicules = Chimique. Fente \SI{3}{nm} + connexons = Électrique.
\item \textbf{Raconter l'expérience de Loewi sans rigueur} : Oublier que le cœur 2 est \textbf{dénervé} et \textbf{non stimulé}. Oublier le \textbf{témoin} (liquide non stimulé). Ne pas nommer l'\textbf{acétylcholine} (Vagusstoff). C'est la preuve de la transmission \textbf{chimique}.
\item \textbf{Inverser sommation spatiale et temporelle} : \textbf{Temporelle} = \emph{une même} afférence, stimulations \emph{rapprochées dans le temps}. \textbf{Spatiale} = \emph{plusieurs} afférences \emph{différentes}, stimulation \emph{simultanée}. Lire l'énoncé : « même fibre » vs « fibres différentes ».
\item \textbf{Erreur de signe dans les calculs de sommation} : $V_{repos} = \SI{-70}{mV}$, Seuil = $\SI{-55}{mV}$. PPSE \SI{+10}{mV} $\rightarrow$ $V_m = \SI{-60}{mV}$. \textbf{$\SI{-60}{mV} < \SI{-55}{mV}$} (plus négatif) $\rightarrow$ \textbf{Sous-seuil}. Toujours écrire le calcul et l'inégalité correcte. Ne pas écrire « -60 est plus grand que -55 ».
\item \textbf{Décrire l'effet d'une drogue sans mécanisme moléculaire} : « La cocaïne excite » = \textbf{0 point}. Attendu : « La cocaïne \textbf{bloque le transporteur de recapture de la dopamine (DAT)} $\rightarrow$ accumulation DA dans fente $\rightarrow$ surstimulation récepteurs D1/D2 voie de la récompense $\rightarrow$ euphorie, dépendance ». \textbf{Cible + Mécanisme + Conséquence synaptique + Effet final}.
\end{enumerate}
\end{attention}

\begin{savaistu}[Message de prévention : Protéger son cerveau et son environnement]
\begin{itemize}
\item \textbf{Le cerveau en construction} : Jusqu'à $\approx$ \SI{25}{ans}, le cortex préfrontal (contrôle inhibiteur, décision) et l'hippocampe (mémoire) sont en maturation synaptique intense (élagage, myélinisation). L'exposition aux substances psychoactives (alcool, cannabis, tramadol, cocaïne) \textbf{perturbe durablement} cette plasticité $\rightarrow$ risques accrus de troubles cognitifs, psychiatriques (schizophrénie, dépression), dépendance.
\item \textbf{Responsabilité collective} :
\begin{itemize}
\item \textbf{Agriculteurs / Applicateurs} : Port systématique de l'EPI (masque, gants, combinaison) lors manipulation pesticides (organophosphorés, pyréthrinoïdes). Respect des DAR (Délais Avant Récolte) et DAR (Délais Avant Rentrée). Stockage hors de portée enfants.
\item \textbf{Personnel de santé} : Prescription rationnelle (antidouleurs opioïdes, benzodiazépines, antidépresseurs). Surveillance interactions (ex. ISRS + Tramadol = Syndrome sérotoninergique).
\item \textbf{Jeunes / Étudiants} : Refuser la première prise. Chercher de l'aide (Centre d'Addictologie, CMP, infirmerie scolaire) si consommation installée. Le sevrage est un parcours médical, pas une faiblesse.
\end{itemize}
\item \textbf{Connaître pour agir} : Comprendre la biologie des synapses (cibles moléculaires, réversibilité, neuroplasticité) permet de faire des choix éclairés et de porter assistance (position latérale de sécurité, appel SAMU \textbf{112} / \textbf{119} au Cameroun) en cas d'intoxication aiguë.
\end{itemize}
\end{savaistu}

\begin{exercice}[Entraînement Bac : Intégration synaptique et pharmacologie]
\begin{enumerate}
\item Sur une électronographie, on observe une fente de \SI{2}{nm} et des connexons. Quel type de synapse ? Justifier. Préciser deux propriétés fonctionnelles majeures.
\item Une substance X appliquée sur une préparation neuromusculaire bloque la contraction nerveuse. L'ACh exogène déclenche encore une contraction. Où et comment agit X ? Donner un exemple de substance réelle.
\item Un neurone reçoit 3 PPSE de \SI{+6}{mV} simultanés (afférences distinctes) et 1 PPSI de \SI{-4}{mV}. $V_{repos} = \SI{-70}{mV}$, Seuil = $\SI{-55}{mV}$. Y a-t-il PA ? Justifier par calcul. Nommer le type de sommation.
\item Un agriculteur présente : myosis, salivation abondante, bronchospasme, bradycardie \SI{40}{bpm}, fasciculations musculaires, incontinence. Quel diagnostic ? Quel antidote en urgence ? Expliquer le mécanisme toxique au niveau synaptique.
\end{enumerate}
\end{exercice}

\begin{corrige}[Exercice n°9 - Intégration synaptique et pharmacologie]
\begin{enumerate}
\item \textbf{Synapse électrique (jonction communicante).} Justification : Fente très étroite (\SI{2}{nm}) \textbf{et} présence de connexons (ponts protéiques intercellulaires) / absence de vésicules synaptiques. Propriétés : Transmission quasi instantanée (pas de délai synaptique) et bidirectionnelle (courant ionique direct).
\item \textbf{Cible : Présynaptique (libération). Mécanisme : Blocage de l'exocytose vésiculaire.} L'ACh exogène fonctionne $\Rightarrow$ récepteurs postsynaptiques intacts. La stimulation nerveuse échoue $\Rightarrow$ pas de libération. Exemple : \textbf{Toxine botulique} (clivage protéines SNARE). Autre exemple : Toxine tétanique (dans SNC, bloque libération GABA/Glycine $\rightarrow$ tétanos spastique).
\item \textbf{Calcul :} $V_m = -70 + (3 \times 6) + (-4) = -70 + 18 - 4 = \SI{-56}{mV}$. $\SI{-56}{mV} < \SI{-55}{mV}$ $\rightarrow$ \textbf{Pas de Potentiel d'Action} (seuil non atteint). Type de sommation : \textbf{Sommation spatiale} (3 afférences distinctes stimulées simultanément) avec composante inhibitrice (sommation algébrique).
\item \textbf{Diagnostic : Intoxication aiguë par organophosphoré (insecticide).} \textbf{Antidotes urgence :} Atropine (IV, doses répétées jusqu'à sécheresse buccale/fréquence cardiaque $> 80$) + Pralidoxime (réactivateur AChE, IV lent) + Diazépam (si convulsions) + Oxygénation/Ventilation. \textbf{Mécanisme synaptique :} Inhibition irréversible de l'acétylcholinestérase (AChE) dans la fente synaptique (plaques motrices, synapses cholinergiques SNA, SNC) $\rightarrow$ accumulation massive d'acétylcholine $\rightarrow$ surstimulation continue récepteurs muscariniques (effets parasympathomimétiques : SLUDGE, bradycardie, myosis, bronchospasme) et nicotiniques (fasciculations, puis paralysie flasque par dépolarisation bloquante).
\end{enumerate}
\end{corrige}

\newpage

\coursec{Exercices}

\courssub{Je vérifie mes connaissances}

\begin{exercice}[QCM : structure et types de synapses]
Pour chaque question, une seule réponse est exacte. Recopie la lettre correspondante.
\begin{enumerate}
\item L'espace qui sépare l'élément présynaptique de l'élément postsynaptique dans une synapse chimique s'appelle :
\begin{enumerate}
\item[a)] la fente synaptique ; b)] la gaine de myéline ; c)] le bouton synaptique ; d)] la plaque motrice.
\end{enumerate}
\item Les vésicules synaptiques contiennent :
\begin{enumerate}
\item[a)] des ions calcium ; b)] des neuromédiateurs ; c)] des récepteurs membranaires ; d)] de l'ATP.
\end{enumerate}
\item La synapse neuro-musculaire est aussi appelée :
\begin{enumerate}
\item[a)] synapse axo-axonique ; b)] jonction gap ; c)] plaque motrice ; d)] synapse neuro-glandulaire.
\end{enumerate}
\item Dans une synapse à transmission électrique, le courant passe directement d'une cellule à l'autre à travers :
\begin{enumerate}
\item[a)] des vésicules ; b)] des canaux jonctionnels (connexons) ; c)] la fente synaptique ; d)] les mitochondries.
\end{enumerate}
\end{enumerate}
\end{exercice}

\begin{exercice}[Vrai ou faux justifié]
Pour chacune des affirmations suivantes, indique si elle est vraie ou fausse, puis justifie ta réponse en une ou deux phrases.
\begin{enumerate}
\item La transmission synaptique chimique est bidirectionnelle.
\item L'expérience de Loewi démontre la nature chimique de la transmission au niveau de la synapse entre le nerf vague et le cœur de grenouille.
\item L'entrée d'ions \ce{Ca^{2+}} dans le bouton présynaptique déclenche l'exocytose des vésicules de neuromédiateur.
\item La sommation temporelle résulte de la stimulation simultanée de deux afférences différentes.
\item Le curare agit en se fixant sur les récepteurs postsynaptiques de l'acétylcholine.
\end{enumerate}
\end{exercice}

\begin{exercice}[Définitions]
Définis en deux ou trois lignes chacun des termes suivants : \cle{synapse}, \cle{neuromédiateur}, \cle{plaque motrice}, \cle{sommation spatiale}, \cle{potentiel postsynaptique excitateur} (PPSE).
\end{exercice}

\begin{exercice}[Calcul direct : vitesse de transmission]
Un influx nerveux parcourt un axone présynaptique de longueur $L = \SI{1,2}{m}$ à la vitesse de $v = \SI{80}{m.s^{-1}}$, puis franchit une synapse chimique avec un retard synaptique de $\Delta t = \SI{0,5}{ms}$.
\begin{enumerate}
\item Calcule la durée du trajet le long de l'axone.
\item Calcule la durée totale de la transmission jusqu'à l'élément postsynaptique.
\item Exprime, en pourcentage, la part du temps total due au franchissement synaptique. Commente.
\end{enumerate}
\end{exercice}

\courssub{Je m'entraîne}

\begin{exercice}[Schématisation d'une synapse chimique]
Réalise un schéma légendé de l'ultrastructure d'une synapse chimique neuro-neuronique telle qu'on l'observe en microscopie électronique. On attend au minimum six légendes : bouton présynaptique, vésicules synaptiques, mitochondries, membrane présynaptique, fente synaptique, membrane postsynaptique avec récepteurs. Indique par une flèche le sens de transmission de l'influx.
\end{exercice}

\begin{exercice}[Identification de synapses sur électronographies]
On dispose de deux électronographies. Document A : on observe deux membranes cellulaires séparées par un espace étroit d'environ \SI{20}{nm}, avec des vésicules d'un seul côté. Document B : les membranes des deux cellules sont séparées par un espace d'environ \SI{3}{nm}, sans vésicule, et des canaux protéiques relient les deux cytoplasmes.
\begin{enumerate}
\item Identifie le type de transmission (chimique ou électrique) dans chaque document. Justifie.
\item Donne un avantage fonctionnel de chaque type de synapse.
\item Cite une localisation possible de chaque type dans l'organisme.
\end{enumerate}
\end{exercice}

\begin{exercice}[Sommations synaptiques]
Un motoneurone reçoit deux afférences excitatrices A et B. On enregistre le potentiel de membrane postsynaptique dans trois situations :
\begin{itemize}
\item Situation 1 : une stimulation unique de A provoque un PPSE de \SI{2}{mV}, insuffisant (seuil : \SI{10}{mV}).
\item Situation 2 : on stimule A cinq fois à intervalles rapprochés ; les PPSE s'additionnent et atteignent \SI{11}{mV} : un potentiel d'action naît.
\item Situation 3 : on stimule simultanément A et B une seule fois ; les PPSE de \SI{2}{mV} et \SI{3}{mV} restent sous le seuil, mais en stimulant trois fois A et B ensemble, on obtient \SI{15}{mV} et un potentiel d'action.
\end{itemize}
\begin{enumerate}
\item Identifie le type de sommation mis en jeu dans les situations 2 et 3. Justifie.
\item Quelle situation illustre une sommation spatio-temporelle ?
\item Déduis de ces résultats le rôle intégrateur du neurone postsynaptique.
\end{enumerate}
\end{exercice}

\begin{exercice}[Effets de substances sur la synapse]
Pour chacune des substances suivantes, précise son site d'action synaptique et l'effet produit sur la transmission : le curare, la tétrodotoxine (bloque les canaux \ce{Na+} voltage-dépendants de la membrane axonale et présynaptique), la cocaïne (bloque la recapture de la dopamine), une benzodiazépine (potentialise les récepteurs du GABA, neuromédiateur inhibiteur). Présente ta réponse sous forme d'un tableau à trois colonnes : substance, site d'action, conséquence sur la transmission.
\end{exercice}

\courssub{Je me prépare au Bac}

\begin{exercice}[Type Bac : l'expérience de Loewi]
Au cours d'une séance de TP au lycée, on reproduit le montage historique de Loewi (1921). Deux cœurs de grenouille, A et B, sont maintenus en survie dans du liquide de Ringer. Le cœur A garde son nerf vague ; le cœur B est dénervé. Le liquide de perfusion du cœur A est recueilli et envoyé sur le cœur B.

\textbf{Protocole et résultats :}
\begin{itemize}
\item Étape 1 : sans stimulation, les deux cœurs battent à la fréquence de 60 battements par minute.
\item Étape 2 : on stimule électriquement le nerf vague du cœur A ; sa fréquence cardiaque chute à 40 battements par minute.
\item Étape 3 : quelques secondes plus tard, sans aucune stimulation, la fréquence du cœur B chute également à 40 battements par minute.
\item Étape 4 : si le liquide de perfusion de A n'est pas transféré vers B, la fréquence du cœur B reste inchangée malgré la stimulation du nerf vague de A.
\end{itemize}

\begin{enumerate}
\item Dégage le problème posé par cette expérience et formule l'hypothèse que Loewi cherchait à tester.
\item Analyse les résultats des étapes 2, 3 et 4 en comparant les fréquences cardiaques.
\item Interprète ces résultats : que révèle l'étape 4 sur le rôle du liquide de perfusion ?
\item Conclus sur la nature (électrique ou chimique) de la transmission entre le nerf vague et le cœur. Nomme le neuromédiateur mis en jeu.
\item Explique, à l'échelle de l'ultrastructure synaptique, les étapes de la libération de ce neuromédiateur.
\end{enumerate}
\end{exercice}

\begin{exercice}[Type Bac : le curare et la plaque motrice]
Le curare est un poison utilisé traditionnellement par certains peuples de chasseurs pour enduire les pointes de flèches. On réalise les expériences suivantes sur un préparé nerf-muscle de grenouille :
\begin{itemize}
\item Expérience 1 : la stimulation du nerf moteur provoque la contraction du muscle.
\item Expérience 2 : après application de curare sur la préparation, la stimulation du nerf moteur ne provoque plus de contraction.
\item Expérience 3 : après application de curare, la stimulation électrique directe du muscle provoque une contraction normale.
\item Expérience 4 : après application de curare, l'injection d'acétylcholine au niveau de la jonction nerf-muscle ne provoque aucune contraction.
\item Expérience 5 : l'enregistrement montre que, malgré le curare, l'influx parcourt normalement l'axone moteur et les vésicules présynaptiques libèrent normalement l'acétylcholine.
\end{itemize}

\begin{enumerate}
\item Définis la plaque motrice et schématise-la avec quatre légendes au moins.
\item À partir des expériences 2 et 3, montre que le curare n'agit ni sur le muscle ni sur la conduction nerveuse.
\item À partir des expériences 4 et 5, localise précisément le site d'action du curare et propose un mécanisme moléculaire de son action.
\item Justifie, à partir des données structurales de la synapse chimique, le sens unidirectionnel de la transmission.
\item Explique pourquoi une intoxication grave au curare peut entraîner la mort par asphyxie.
\end{enumerate}
\end{exercice}

\begin{exercice}[Type Bac : cocaïne, dopamine et sensibilisation au tramadol]
La dopamine est un neuromédiateur libéré dans les synapses du « circuit de la récompense ». Normalement, après sa libération, elle est recapturée par le neurone présynaptique grâce à un transporteur membranaire (DAT). On mesure la concentration de dopamine dans la fente synaptique chez un animal de laboratoire :

\begin{center}
\begin{tabularx}{\linewidth}{|l|X|X|X|X|}
\hline
\rowcolor{popblueL} Temps (min) & 0 & 10 & 30 & 60 \\
\hline
Témoin (sans drogue) & 5 & 5 & 5 & 5 \\
\hline
Après injection de cocaïne & 5 & 40 & 35 & 12 \\
\hline
\end{tabularx}
\end{center}
(Concentrations en unités arbitraires ; la cocaïne est injectée à $t = 0$.)

\begin{enumerate}
\item Trace sur un même graphique les deux courbes d'évolution de la concentration de dopamine synaptique en fonction du temps.
\item Explique, à partir de tes connaissances sur le fonctionnement de la synapse, le mécanisme d'action de la cocaïne responsable des variations observées.
\item Explique l'origine de la sensation de plaisir intense puis la chute de l'effet, et montre comment ces phénomènes conduisent à la dépendance.
\item Le tramadol, médicament antalgique détourné de son usage et très consommé par certains jeunes au Cameroun, agit comme un agoniste des récepteurs opioïdes $\mu$ et inhibe la recapture de la sérotonine et de la noradrénaline. Lors d'un surdosage, il déprime les centres respiratoires du tronc cérébral et abaisse le seuil épileptogène. Rédige un texte argumenté d'environ quinze lignes, destiné à un camarade, expliquant scientifiquement pourquoi la consommation abusive de tramadol peut entraîner des convulsions, une dépression respiratoire puis le coma. Tu t'appuieras sur tes connaissances de la transmission synaptique (excitation/inhibition, sommation, neuromédiateurs) et tu concluras par un conseil de prévention.
\end{enumerate}
\end{exercice}

\newpage

\coursec{Corrigés des exercices}

\courssub{Je vérifie mes connaissances}

\begin{corrige}[Exercice 1 — QCM : structure et types de synapses]
\begin{enumerate}
\item \textbf{Réponse a)} : la fente synaptique. C'est l'espace extracellulaire d'environ \SI{20}{nm} qui sépare la membrane présynaptique de la membrane postsynaptique dans une synapse chimique. La gaine de myéline entoure les axones, le bouton synaptique est l'extrémité renflée de l'axone présynaptique, et la plaque motrice est la synapse neuro-musculaire.
\item \textbf{Réponse b)} : des neuromédiateurs. Les vésicules synaptiques stockent les molécules de neuromédiateur (acétylcholine, dopamine, GABA...) qui seront libérées par exocytose. Les ions \ce{Ca^{2+}} arrivent du milieu extracellulaire, les récepteurs sont postsynaptiques, et l'ATP est produit par les mitochondries.
\item \textbf{Réponse c)} : plaque motrice. C'est la synapse entre un motoneurone et une fibre musculaire squelettique ; le neuromédiateur y est l'acétylcholine.
\item \textbf{Réponse b)} : des canaux jonctionnels (connexons). Dans les jonctions gap, des canaux protéiques (connexons, formés de six sous-unités de connexine) relient directement les deux cytoplasmes, permettant le passage direct du courant ionique, sans neuromédiateur.
\end{enumerate}
\end{corrige}

\begin{corrige}[Exercice 2 — Vrai ou faux justifié]
\begin{enumerate}
\item \textbf{Faux.} La transmission chimique est \emph{unidirectionnelle} : seul l'élément présynaptique possède les vésicules de neuromédiateur, et seul l'élément postsynaptique porte les récepteurs spécifiques. Le message ne peut donc passer que du pré- vers le postsynaptique.
\item \textbf{Vrai.} En stimulant le nerf vague d'un cœur A puis en transférant son liquide de perfusion sur un cœur B dénervé, Loewi observe le ralentissement de B : une substance chimique (l'acétylcholine, « Vagusstoff ») libérée par le nerf est responsable de l'effet, ce qui prouve la nature chimique de la transmission.
\item \textbf{Vrai.} L'arrivée du potentiel d'action ouvre les canaux calciques voltage-dépendants de la membrane présynaptique ; l'entrée de \ce{Ca^{2+}} déclenche la fusion des vésicules avec la membrane et l'exocytose du neuromédiateur. Le calcium est le signal couplant l'influx à la libération.
\item \textbf{Faux.} La stimulation simultanée de deux afférences différentes définit la sommation \emph{spatiale}. La sommation temporelle correspond à l'addition de PPSE provoqués par des stimulations successives rapprochées de la \emph{même} afférence.
\item \textbf{Vrai.} Le curare est un antagoniste compétitif des récepteurs nicotiniques de l'acétylcholine au niveau de la plaque motrice : il occupe les récepteurs sans les activer, bloquant ainsi la transmission neuro-musculaire.
\end{enumerate}
\end{corrige}

\begin{corrige}[Exercice 3 — Définitions]
\begin{itemize}
\item \cle{Synapse} : zone de contact fonctionnel entre un neurone et une autre cellule (neurone, cellule musculaire ou glandulaire), où l'influx nerveux est transmis d'un élément présynaptique à un élément postsynaptique, par voie chimique ou électrique.
\item \cle{Neuromédiateur} : substance chimique (acétylcholine, dopamine, GABA, sérotonine...) synthétisée et stockée dans les vésicules du bouton présynaptique, libérée par exocytose dans la fente synaptique et capable de se fixer sur des récepteurs postsynaptiques spécifiques pour y modifier la perméabilité membranaire.
\item \cle{Plaque motrice} : synapse neuro-musculaire entre la terminaison d'un motoneurone et une fibre musculaire squelettique ; l'acétylcholine libérée provoque une dépolarisation de la membrane musculaire (plaque terminale) déclenchant la contraction.
\item \cle{Sommation spatiale} : addition algébrique, sur le neurone postsynaptique, des potentiels postsynaptiques générés simultanément par plusieurs afférences distinctes ; si la somme atteint le seuil au niveau du cône d'émergence, un potentiel d'action naît.
\item \cle{Potentiel postsynaptique excitateur} (PPSE) : dépolarisation locale, graduée et non propagée, de la membrane postsynaptique, résultant de l'ouverture de canaux cationiques (entrée de \ce{Na+}) après fixation d'un neuromédiateur excitateur ; il rapproche le potentiel de membrane du seuil.
\end{itemize}
\end{corrige}

\begin{corrige}[Exercice 4 — Calcul direct : vitesse de transmission]
\begin{enumerate}
\item Durée du trajet axonal :
\[ t_{\text{axone}} = \frac{L}{v} = \frac{\SI{1,2}{m}}{\SI{80}{m.s^{-1}}} = \SI{0,015}{s} = \SI{15}{ms}. \]
\item Durée totale de la transmission :
\[ t_{\text{total}} = t_{\text{axone}} + \Delta t = 15 + 0,5 = \SI{15,5}{ms}. \]
\item Part due au franchissement synaptique :
\[ p = \frac{\Delta t}{t_{\text{total}}} \times 100 = \frac{0,5}{15,5} \times 100 \approx 3,2\,\%. \]
\textbf{Commentaire :} le retard synaptique ne représente ici qu'environ 3\,\% du temps total car l'axone est long. Mais à l'échelle d'un circuit comportant de nombreuses synapses en série (réflexes complexes), chaque synapse ajoutant environ \SI{0,5}{ms}, le cumul des retards synaptiques devient prépondérant : c'est ce qui explique la durée mesurable des réflexes polysynaptiques.
\end{enumerate}
\end{corrige}

\courssub{Je m'entraîne}

\begin{corrige}[Exercice 5 — Schématisation d'une synapse chimique]
Le schéma attendu doit respecter l'organisation observée en microscopie électronique : bouton présynaptique renflé contenant vésicules et mitochondries, fente synaptique étroite, membrane postsynaptique porteuse de récepteurs.

\begin{popfigure}
\begin{center}
\begin{tikzpicture}[scale=1.0, font=\small]
% Bouton présynaptique
\fill[poporangeL] (0,0) ellipse (2.6 and 1.6);
\draw[popdark, thick] (0,0) ellipse (2.6 and 1.6);
% Vésicules
\foreach \x/\y in {-1.2/0.5,-0.4/0.8,0.5/0.6,1.2/0.3,-0.8/-0.2,0.2/0.1,1.0/-0.5,-1.5/-0.5,0.6/-0.9,-0.3/-0.7}
  {\draw[poporange, fill=white] (\x,\y) circle (0.16);}
% Mitochondries
\draw[popgreen, fill=popgreenL, thick] (-1.7,0.9) ellipse (0.55 and 0.28);
\draw[popgreen] (-1.7,0.9) ellipse (0.3 and 0.12);
\draw[popgreen, fill=popgreenL, thick] (1.6,0.9) ellipse (0.55 and 0.28);
\draw[popgreen] (1.6,0.9) ellipse (0.3 and 0.12);
% Membrane présynaptique (face active)
\draw[popdark, very thick] (-2.2,-1.55) -- (2.2,-1.55);
% Fente synaptique
\fill[popblueL] (-2.6,-2.1) rectangle (2.6,-1.55);
% Points de neuromédiateur dans la fente
\foreach \x in {-1.5,-0.9,-0.3,0.3,0.9,1.5}
  {\fill[poporange] (\x,-1.85) circle (0.07);}
% Membrane postsynaptique
\draw[popdark, very thick] (-2.6,-2.1) -- (2.6,-2.1);
% Récepteurs
\foreach \x in {-1.8,-0.9,0,0.9,1.8}
  {\draw[popblue, very thick] (\x,-2.1) -- (\x,-2.35) (\x-0.12,-2.35) -- (\x+0.12,-2.35);}
% Élément postsynaptique
\fill[popblue!10] (-2.6,-3.2) rectangle (2.6,-2.35);
\draw[popblue, thick] (-2.6,-3.2) rectangle (2.6,-2.35);
% Flèche sens de transmission
\draw[-{Stealth[length=4mm]}, poppink, very thick] (3.4,-0.3) -- (3.4,-2.8);
\node[poppink, rotate=90] at (3.8,-1.55) {\textbf{sens de transmission}};
% Légendes
\node[anchor=west] at (-6.4,1.6) {\textbf{1}};
\draw[-, popmuted] (-6.2,1.6) -- (-1.9,1.0);
\node[anchor=west] at (-6.4,0.9) {\textbf{2}};
\draw[-, popmuted] (-6.2,0.9) -- (-0.5,0.75);
\node[anchor=west] at (-6.4,0.2) {\textbf{3}};
\draw[-, popmuted] (-6.2,0.2) -- (-1.75,0.85);
\node[anchor=west] at (-6.4,-1.55) {\textbf{4}};
\draw[-, popmuted] (-6.2,-1.55) -- (-2.3,-1.55);
\node[anchor=west] at (-6.4,-2.35) {\textbf{5}};
\draw[-, popmuted] (-6.2,-2.35) -- (-2.3,-1.85);
\node[anchor=west] at (-6.4,-3.05) {\textbf{6}};
\draw[-, popmuted] (-6.2,-3.05) -- (-1.0,-2.25);
\end{tikzpicture}
\end{center}
\legende{Ultrastructure d'une synapse chimique neuro-neuronique. \textbf{1} : bouton présynaptique ; \textbf{2} : vésicules synaptiques (contenant le neuromédiateur) ; \textbf{3} : mitochondrie (fournit l'ATP) ; \textbf{4} : membrane présynaptique ; \textbf{5} : fente synaptique ; \textbf{6} : membrane postsynaptique porteuse de récepteurs. La flèche rose indique le sens unidirectionnel de transmission (pré $\rightarrow$ post).}
\end{popfigure}

\textbf{Points de vigilance :} la fente synaptique doit être visible et étroite ; les vésicules ne figurent que du côté présynaptique ; les récepteurs uniquement sur la membrane postsynaptique ; le sens de transmission doit être fléché du pré- vers le postsynaptique. Évite de représenter des vésicules des deux côtés, ce qui contredirait l'unidirectionnalité.
\end{corrige}

\begin{corrige}[Exercice 6 — Identification de synapses sur électronographies]
\begin{enumerate}
\item \textbf{Document A : synapse à transmission chimique.} Justification : fente synaptique large (\SI{20}{nm}) et présence de vésicules d'un seul côté (présynaptique), signe d'une libération de neuromédiateur par exocytose.
\textbf{Document B : synapse à transmission électrique (jonction gap).} Justification : espace intercellulaire très étroit (\SI{3}{nm}), absence de vésicules, et présence de canaux jonctionnels (connexons) reliant les deux cytoplasmes, permettant le passage direct du courant ionique.
\item \textbf{Avantages :} la synapse chimique permet la \emph{modulation} du signal (excitation ou inhibition, amplification, intégration par sommation) ; la synapse électrique assure une transmission \emph{quasi instantanée} (pas de retard synaptique) et souvent bidirectionnelle, permettant la synchronisation de populations cellulaires.
\item \textbf{Localisations :} synapse chimique : jonction nerf-muscle squelettique (plaque motrice), synapses du système nerveux central et du système nerveux autonome (ex. nerf vague-cœur). Synapse électrique : certaines jonctions entre neurones du système nerveux central, entre cellules du muscle cardiaque et du muscle lisse (assurant la contraction synchronisée).
\end{enumerate}
\end{corrige}

\begin{corrige}[Exercice 7 — Sommations synaptiques]
\begin{enumerate}
\item \textbf{Situation 2 : sommation temporelle.} Une seule et même afférence A est stimulée plusieurs fois à intervalles rapprochés : les PPSE successifs de \SI{2}{mV} s'additionnent ($5 \times 2 = \SI{10}{mV}$, dépassé avec la décroissance partielle, atteignant \SI{11}{mV}) et franchissent le seuil de \SI{10}{mV}.
\textbf{Situation 3 : sommation spatiale puis spatio-temporelle.} La stimulation simultanée de deux afférences distinctes A et B illustre la sommation spatiale ($2 + 3 = \SI{5}{mV}$, sous le seuil) ; la répétition de cette stimulation conjointe ajoute une composante temporelle.
\item C'est la \textbf{situation 3} (lors des trois stimulations répétées de A et B ensemble) qui illustre la sommation \emph{spatio-temporelle} : elle combine l'addition de signaux issus de plusieurs afférences (espace) et de stimulations successives rapprochées (temps), atteignant \SI{15}{mV} $>$ \SI{10}{mV}.
\item \textbf{Rôle intégrateur :} le neurone postsynaptique additionne algébriquement, à chaque instant, l'ensemble des signaux excitateurs (PPSE) et inhibiteurs (PPSI) qu'il reçoit de ses multiples afférences. Ce n'est que si la somme des dépolarisations atteint le seuil au niveau du cône d'émergence qu'un potentiel d'action est généré (loi du tout ou rien). Le neurone fonctionne donc comme un \emph{centre d'intégration} qui « décide » de transmettre ou non l'information.
\end{enumerate}
\end{corrige}

\begin{corrige}[Exercice 8 — Effets de substances sur la synapse]
\begin{center}
\begin{tabularx}{\linewidth}{|l|X|X|}
\hline
\rowcolor{popblueL} \textbf{Substance} & \textbf{Site d'action} & \textbf{Conséquence sur la transmission} \\
\hline
Curare & Récepteurs nicotiniques de l'acétylcholine sur la membrane postsynaptique (plaque motrice) & Antagoniste compétitif : l'acétylcholine ne peut plus se fixer ; blocage de la transmission neuro-musculaire, paralysie flasque. \\
\hline
Tétrodotoxine (TTX) & Canaux \ce{Na+} voltage-dépendants de la membrane axonale et présynaptique & Blocage de la naissance et de la conduction du potentiel d'action : l'influx n'atteint pas le bouton synaptique, pas de libération de neuromédiateur ; transmission supprimée. \\
\hline
Cocaïne & Transporteur de recapture de la dopamine (DAT) sur la membrane présynaptique & Accumulation de dopamine dans la fente synaptique : stimulation excessive et prolongée des récepteurs postsynaptiques ; transmission dopaminergique exagérée. \\
\hline
Benzodiazépine & Récepteurs du GABA (site modulateur) sur la membrane postsynaptique & Potentialisation de l'effet inhibiteur du GABA (entrée accrue de \ce{Cl-}, hyperpolarisation) : diminution de l'excitabilité neuronale, effet anxiolytique et sédatif. \\
\hline
\end{tabularx}
\end{center}
\textbf{Point de vigilance :} toujours distinguer le \emph{site} (présynaptique, fente, postsynaptique) et le \emph{sens} de l'effet (facilitation ou blocage). La TTX agit en amont de la synapse (conduction), le curare au niveau postsynaptique, la cocaïne sur la recapture : trois cibles différentes, trois mécanismes distincts.
\end{corrige}

\courssub{Je me prépare au Bac}

\begin{corrige}[Exercice 9 — Type Bac : l'expérience de Loewi]
\begin{enumerate}
\item \textbf{Problème :} comment l'influx nerveux du nerf vague est-il transmis au muscle cardiaque ? La transmission est-elle électrique (passage direct de courant) ou chimique (libération d'une substance) ? \textbf{Hypothèse de Loewi :} la stimulation du nerf vague provoque la libération d'une substance chimique diffusible, responsable du ralentissement cardiaque.
\item \textbf{Analyse :} à l'étape 2, la stimulation du nerf vague fait chuter la fréquence du cœur A de 60 à 40 battements/min (diminution de 20 battements/min, soit $-33\,\%$) : le nerf vague exerce un effet cardio-modérateur. À l'étape 3, le cœur B, \emph{dénervé} et non stimulé, subit la même chute (60 $\rightarrow$ 40 battements/min) après réception du liquide de perfusion de A. À l'étape 4, sans transfert de liquide, B reste à 60 battements/min malgré la stimulation de A.
\item \textbf{Interprétation :} l'étape 4 est l'expérience témoin : elle montre que le ralentissement de B n'est possible que si le liquide de perfusion de A est transféré. Ce liquide contient donc une \emph{substance chimique} libérée par le nerf vague stimulé, capable à elle seule de reproduire l'effet de la stimulation nerveuse. Le message transmis est de nature chimique.
\item \textbf{Conclusion :} la transmission entre le nerf vague et le cœur est de nature \emph{chimique}. Le neuromédiateur mis en jeu est l'\emph{acétylcholine} (le « Vagusstoff » de Loewi), libérée par les terminaisons parasympathiques.
\item \textbf{Étapes de la libération (ultrastructure) :} (1) l'arrivée du potentiel d'action au bouton présynaptique dépolarise la membrane ; (2) les canaux calciques voltage-dépendants s'ouvrent et les ions \ce{Ca^{2+}} entrent dans le bouton ; (3) l'élévation du calcium intracellulaire provoque la migration et la fusion des vésicules synaptiques avec la membrane présynaptique ; (4) l'acétylcholine est libérée par exocytose dans la fente synaptique ; (5) elle diffuse et se fixe sur les récepteurs postsynaptiques des cellules cardiaques, provoquant le ralentissement ; (6) elle est ensuite dégradée par l'acétylcholinestérase, ce qui limite la durée du signal.
\end{enumerate}
\textbf{Barème indicatif (sur 10) :} problème et hypothèse : 2 pts ; analyse chiffrée des étapes 2-3-4 : 2 pts ; rôle témoin de l'étape 4 : 1,5 pt ; conclusion (nature chimique + acétylcholine) : 1,5 pt ; étapes de la libération : 3 pts.
\textbf{Points de vigilance :} comparer explicitement les valeurs chiffrées ; identifier l'étape 4 comme \emph{témoin} ; ne pas écrire que le liquide « transporte l'influx » : il transporte une substance chimique ; citer le rôle de \ce{Ca^{2+}} dans l'exocytose.
\end{corrige}

\begin{corrige}[Exercice 10 — Type Bac : le curare et la plaque motrice]
\begin{enumerate}
\item \textbf{Définition :} la plaque motrice est la synapse neuro-musculaire unissant la terminaison d'un motoneurone à une fibre musculaire squelettique ; le neuromédiateur y est l'acétylcholine. \textbf{Schéma attendu :} ramification terminale de l'axone moteur avec vésicules, fente synaptique, membrane musculaire (sarcolemme) repliée en gouttières porteuses de récepteurs cholinergiques, noyaux et mitochondries. Légendes minimales : bouton présynaptique avec vésicules d'acétylcholine ; fente synaptique ; membrane postsynaptique (sarcolemme) avec récepteurs ; fibre musculaire.
\item \textbf{Le curare n'agit ni sur le muscle ni sur la conduction nerveuse :} l'expérience 3 montre que la stimulation \emph{directe} du muscle curarisé provoque une contraction normale : la fibre musculaire est donc fonctionnelle, le curare ne la paralyse pas. L'expérience 2 montre que la stimulation du nerf devient inefficace alors que le muscle est sain : le blocage se situe donc entre le nerf et le muscle, c'est-à-dire au niveau de la \emph{jonction} (plaque motrice). L'expérience 5 confirmera que la conduction axonale est intacte.
\item \textbf{Localisation et mécanisme :} l'expérience 5 montre que l'influx parcourt normalement l'axone et que l'acétylcholine est normalement libérée : l'élément présynaptique est intact. L'expérience 4 montre que même de l'acétylcholine injectée en excès ne provoque aucune contraction : le blocage se situe donc au niveau des \emph{récepteurs postsynaptiques}. Mécanisme moléculaire : le curare est un \emph{antagoniste compétitif} des récepteurs nicotiniques de l'acétylcholine : il se fixe sur les récepteurs sans les activer et empêche l'acétylcholine de s'y lier ; les canaux cationiques ne s'ouvrent pas, aucune dépolarisation de la plaque terminale ne se produit, et la contraction est impossible.
\item \textbf{Sens unidirectionnel :} l'asymétrie structurale de la synapse chimique l'impose : les vésicules de neuromédiateur et les sites d'exocytose n'existent que dans l'élément présynaptique, tandis que les récepteurs spécifiques ne sont présents que sur la membrane postsynaptique. Le message ne peut donc passer que du neurone vers le muscle, jamais en sens inverse.
\item \textbf{Mort par asphyxie :} les muscles respiratoires (diaphragme et muscles intercostaux) sont des muscles squelettiques commandés par des plaques motrices cholinergiques. En cas d'intoxication grave, le curare bloque ces jonctions : les muscles respiratoires sont paralysés, la ventilation s'arrête, et l'organisme meurt par asphyxie alors que la victime reste consciente (le curare ne traverse pas la barrière hémato-encéphalique).
\end{enumerate}
\textbf{Barème indicatif (sur 10) :} définition + schéma légendé : 2,5 pts ; analyse des expériences 2 et 3 : 2 pts ; localisation et mécanisme moléculaire : 2,5 pts ; unidirectionnalité justifiée par la structure : 1,5 pt ; explication de l'asphyxie : 1,5 pt.
\textbf{Points de vigilance :} raisonner par élimination en citant chaque expérience ; employer le terme « antagoniste compétitif » ; ne pas écrire que le curare « bloque l'influx » : il bloque la \emph{transmission}, pas la conduction.
\end{corrige}

\begin{corrige}[Exercice 11 — Type Bac : cocaïne, dopamine et sensibilisation au tramadol]
\begin{enumerate}
\item \textbf{Graphique :}

\begin{popfigure}
\begin{center}
\begin{tikzpicture}
\begin{axis}[
  width=0.85\linewidth, height=6.2cm,
  xlabel={Temps (min)}, ylabel={Concentration de dopamine (u. a.)},
  xmin=0, xmax=62, ymin=0, ymax=45,
  xtick={0,10,30,60}, ytick={0,5,12,35,40},
  grid=major, grid style={popline!50},
  legend style={at={(0.02,0.97)}, anchor=north west, font=\small}
]
\addplot[popcurve, color=popblue, mark=*, mark size=1.5pt] coordinates {(0,5) (10,5) (30,5) (60,5)};
\addlegendentry{Témoin (sans drogue)}
\addplot[popcurve, color=poppink, mark=square*, mark size=1.5pt] coordinates {(0,5) (10,40) (30,35) (60,12)};
\addlegendentry{Après injection de cocaïne}
\end{axis}
\end{tikzpicture}
\end{center}
\legende{Évolution de la concentration de dopamine dans la fente synaptique en fonction du temps, chez le témoin (bleu) et après injection de cocaïne à $t = 0$ (rose).}
\end{popfigure}

\item \textbf{Mécanisme d'action de la cocaïne :} normalement, la dopamine libérée dans la fente synaptique est rapidement recapturée par le neurone présynaptique grâce au transporteur DAT, ce qui limite la durée de son action sur les récepteurs postsynaptiques. La cocaïne \emph{bloque le transporteur DAT} : la recapture est inhibée, la dopamine s'accumule dans la fente synaptique (pic à 40 u. a. à $t = 10$ min, soit huit fois la valeur témoin) et stimule de façon excessive et prolongée les récepteurs postsynaptiques. La transmission dopaminergique est donc fortement amplifiée.
\item \textbf{Plaisir, chute de l'effet et dépendance :} la stimulation excessive et prolongée des synapses dopaminergiques du circuit de la récompense provoque une sensation de plaisir intense (euphorie). Puis la dopamine est progressivement éliminée (dégradation enzymatique, diffusion) : la concentration retombe (12 u. a. à $t = 60$ min) et l'effet disparaît, laissant un état de manque. De plus, la stimulation répétée entraîne une désensibilisation (diminution du nombre ou de la sensibilité des récepteurs postsynaptiques) : le sujet doit consommer à nouveau, et à doses croissantes, pour retrouver l'effet initial. C'est le cercle vicieux de la \emph{dépendance} (tolérance puis addiction).
\item \textbf{Texte de sensibilisation sur le tramadol (exemple de rédaction) :}

Le tramadol est un médicament antalgique qui ne doit être pris que sur ordonnance. Détourné de son usage, il agit sur les synapses du système nerveux central de plusieurs façons. Comme agoniste des récepteurs opioïdes $\mu$, il mime l'action des endorphines et déprime l'activité de nombreux neurones, notamment ceux des centres respiratoires du tronc cérébral : en surdosage, ces centres n'intègrent plus correctement leurs afférences, la commande des muscles respiratoires devient insuffisante, et la ventilation ralentit jusqu'à la dépression respiratoire, privant le cerveau d'oxygène. Par ailleurs, en inhibant la recapture de la sérotonine et de la noradrénaline, le tramadol provoque l'accumulation de ces neuromédiateurs excitateurs dans les fentes synaptiques, comme le fait la cocaïne pour la dopamine. L'excès d'excitation, par sommation spatiale et temporelle de PPSE anormalement nombreux, fait franchir le seuil de décharge à de vastes populations de neurones : c'est la crise convulsive généralisée, d'autant plus facile que le tramadol abaisse le seuil épileptogène. L'association d'une hyperexcitation cérébrale et d'une dépression des centres vitaux conduit au coma, voire au décès. La « force » recherchée n'est qu'une perturbation dangereuse de tes synapses. Refuse le tramadol sans prescription, parle-en à un adulte ou à un personnel de santé : ton cerveau vaut plus qu'un faux moment de courage.
\end{enumerate}
\textbf{Barème indicatif (sur 12) :} graphique correct (axes légendés, deux courbes) : 2,5 pts ; mécanisme de la cocaïne (blocage de la recapture) : 2,5 pts ; plaisir, tolérance et dépendance : 2 pts ; texte argumenté (mécanismes synaptiques corrects, cohérence, conseil de prévention) : 5 pts.
\textbf{Points de vigilance :} légender les axes avec les unités ; citer explicitement la recapture et le transporteur DAT ; dans le texte, mobiliser les notions d'excitation/inhibition, de sommation et de seuil ; rester scientifique tout en adoptant un ton de sensibilisation adapté à un camarade.
\end{corrige}

\newpage

\coursec{Fiche bilan}

\begin{popfigure}
\begin{center}\tcbox[colback=popblue,colframe=popblue,coltext=white,arc=9pt,boxsep=4pt,left=6pt,right=6pt]{\bfseries\small La synapse et la transmission synaptique}\end{center}
\vspace{-2pt}
\begin{tcbraster}[raster columns=3,raster equal height=rows,raster column skip=6pt,raster row skip=6pt,enhanced,blankest]
\begin{tcolorbox}[enhanced,colback=white,colframe=popgreen,boxrule=0.9pt,arc=3pt,title={Structure et types anatomiques},fonttitle=\bfseries\scriptsize,coltitle=white,colbacktitle=popgreen,left=4pt,right=4pt,top=2pt,bottom=2pt,boxsep=1pt,toptitle=1.5pt,bottomtitle=1.5pt,halign=left]
\begin{itemize}[nosep,leftmargin=*,itemsep=1pt,font=\scriptsize]
\item Neuro-neuronique
\item Neuro-musculaire (plaque motrice)
\item Neuro-glandulaire
\end{itemize}
\end{tcolorbox}
\begin{tcolorbox}[enhanced,colback=white,colframe=poporange,boxrule=0.9pt,arc=3pt,title={Nature de la transmission},fonttitle=\bfseries\scriptsize,coltitle=white,colbacktitle=poporange,left=4pt,right=4pt,top=2pt,bottom=2pt,boxsep=1pt,toptitle=1.5pt,bottomtitle=1.5pt,halign=left]
\begin{itemize}[nosep,leftmargin=*,itemsep=1pt,font=\scriptsize]
\item Expérience de Loewi (1921)
\item Chimique : médiateur
\item Électrique : jonctions gap
\end{itemize}
\end{tcolorbox}
\begin{tcolorbox}[enhanced,colback=white,colframe=poppurple,boxrule=0.9pt,arc=3pt,title={Synapse chimique : étapes},fonttitle=\bfseries\scriptsize,coltitle=white,colbacktitle=poppurple,left=4pt,right=4pt,top=2pt,bottom=2pt,boxsep=1pt,toptitle=1.5pt,bottomtitle=1.5pt,halign=left]
\begin{itemize}[nosep,leftmargin=*,itemsep=1pt,font=\scriptsize]
\item PA $\rightarrow$ exocytose
\item Liaison médiateur--récepteur
\item PPS / PPI puis recyclage
\end{itemize}
\end{tcolorbox}
\begin{tcolorbox}[enhanced,colback=white,colframe=popgold,boxrule=0.9pt,arc=3pt,title={Sommations},fonttitle=\bfseries\scriptsize,coltitle=white,colbacktitle=popgold,left=4pt,right=4pt,top=2pt,bottom=2pt,boxsep=1pt,toptitle=1.5pt,bottomtitle=1.5pt,halign=left]
\begin{itemize}[nosep,leftmargin=*,itemsep=1pt,font=\scriptsize]
\item Spatiale
\item Temporelle
\item Rôle intégrateur du neurone
\end{itemize}
\end{tcolorbox}
\begin{tcolorbox}[enhanced,colback=white,colframe=poppink,boxrule=0.9pt,arc=3pt,title={Substances et synapse},fonttitle=\bfseries\scriptsize,coltitle=white,colbacktitle=poppink,left=4pt,right=4pt,top=2pt,bottom=2pt,boxsep=1pt,toptitle=1.5pt,bottomtitle=1.5pt,halign=left]
\begin{itemize}[nosep,leftmargin=*,itemsep=1pt,font=\scriptsize]
\item Agonistes / antagonistes
\item Toxines : tétanos, botulisme
\item Drogues : dépendance
\end{itemize}
\end{tcolorbox}
\begin{tcolorbox}[enhanced,colback=white,colframe=popblue!70,boxrule=0.9pt,arc=3pt,title={Sensibilisation},fonttitle=\bfseries\scriptsize,coltitle=white,colbacktitle=popblue!70,left=4pt,right=4pt,top=2pt,bottom=2pt,boxsep=1pt,toptitle=1.5pt,bottomtitle=1.5pt,halign=left]
\begin{itemize}[nosep,leftmargin=*,itemsep=1pt,font=\scriptsize]
\item Médicaments : ordonnance
\item Refus des drogues
\end{itemize}
\end{tcolorbox}
\end{tcbraster}

\legende{Carte mentale de la leçon : la synapse, sa structure, ses mécanismes de transmission, l'intégration des messages et l'action des substances (drogues, médicaments, toxines).}
\end{popfigure}

\begin{bilan}
\textbf{Définitions clés}
\begin{itemize}
  \item \cle{Synapse} : zone de contact fonctionnel entre un neurone et une autre cellule (neurone, fibre musculaire, cellule glandulaire), où le message nerveux est transmis sans continuité cytoplasmique.
  \item \cle{Médiateur chimique} (neurotransmetteur) : substance libérée par l'élément présynaptique (ex. acétylcholine, noradrénaline) qui se fixe sur des récepteurs postsynaptiques spécifiques.
  \item \cle{Plaque motrice} : synapse neuro-musculaire entre un motoneurone et une fibre musculaire squelettique ; le médiateur y est l'acétylcholine.
  \item \cle{PPS} (potentiel postsynaptique) : variation du potentiel de membrane postsynaptique ; \cle{PPSE} excitateur (dépolarisation), \cle{PPSI} inhibiteur (hyperpolarisation).
  \item \cle{Sommation} : addition des PPS au niveau du cône d'émergence ; spatiale (plusieurs synapses simultanées), temporelle (influx successifs d'une même synapse), spatio-temporelle (les deux).
\end{itemize}

\textbf{Connaissances à maîtriser par cœur}
\begin{center}
\begin{tabularx}{\linewidth}{|l|X|X|}
\hline
\rowcolor{popblueL} \textbf{Critère} & \textbf{Synapse chimique} & \textbf{Synapse électrique} \\
\hline
Contact & Discontinu (fente synaptique de 20 à \SI{40}{nm}) & Continu (jonctions gap, $\approx$ \SI{2}{nm}) \\
\hline
Support du message & Médiateur chimique & Courant ionique direct \\
\hline
Vitesse & Retard synaptique ($\approx$ \SI{0,5}{ms}) & Quasi instantanée \\
\hline
Sens de transmission & Unidirectionnel & Souvent bidirectionnel \\
\hline
\end{tabularx}
\end{center}

\textbf{Étapes de la transmission chimique} : arrivée du PA $\rightarrow$ ouverture des canaux Ca$^{2+}$ $\rightarrow$ exocytose des vésicules $\rightarrow$ diffusion du médiateur $\rightarrow$ fixation sur récepteurs $\rightarrow$ PPS $\rightarrow$ destruction ou recapture du médiateur.

\textbf{Méthodes express}
\begin{itemize}
  \item Sur une électronographie : vésicules + fente large = synapse \cle{chimique} ; membranes accolées sans vésicules = synapse \cle{électrique}.
  \item Expérience de Loewi : le liquide du cœur stimulé reproduit l'effet sur un cœur dénervé $\Rightarrow$ transmission \cle{chimique}.
  \item Substances : \cle{agoniste} (mimique le médiateur), \cle{antagoniste} (bloque le récepteur), inhibiteur d'enzyme (prolonge l'action, ex. anticholinestérasiques), toxines (tétanospasmine : spasmes ; toxine botulique : paralysie flasque).
\end{itemize}
\end{bilan}

\begin{aretenir}[Checklist avant le Bac]
\begin{itemize}
  \item $\square$ Je sais schématiser et légender l'ultrastructure d'une synapse chimique (vésicules, fente, récepteurs, mitochondries).
  \item $\square$ Je distingue les trois types anatomiques de synapses : neuro-neuronique, neuro-musculaire, neuro-glandulaire.
  \item $\square$ J'identifie une synapse chimique et une synapse électrique sur une électronographie.
  \item $\square$ Je sais décrire le protocole de Loewi et conclure à la nature chimique de la transmission.
  \item $\square$ Je décris dans l'ordre les étapes de la libération et de l'action d'un médiateur.
  \item $\square$ Je distingue PPSE et PPSI et j'explique leur rôle.
  \item $\square$ Je différencie sommation spatiale, temporelle et spatio-temporelle sur un document.
  \item $\square$ J'explique le rôle intégrateur du neurone au niveau du cône d'émergence.
  \item $\square$ Je cite au moins trois substances (drogue, médicament, toxine) et leur mode d'action synaptique précis.
  \item $\square$ Je sais expliquer le mécanisme de la dépendance et argumenter un message de sensibilisation.
  \item $\square$ Je rédige mes réponses avec le vocabulaire exact : présynaptique, postsynaptique, exocytose, recapture.
\end{itemize}
\end{aretenir}

\findecours

\end{document}
